% ============================================================
% RMC_Paper.tex
% Reasoning-Margin Certification: Theory, Framework, and
% Empirical Evaluation on 3-SAT and Vulnerability Detection
% ============================================================
\documentclass[conference]{IEEEtran}
\IEEEoverridecommandlockouts
% The preceding line is only needed to identify funding in the first footnote. If that is unneeded, please comment it out.
%Template version as of 6/27/2024

\usepackage{cite}
\usepackage{amsmath,amssymb,amsfonts}
\usepackage{graphicx}
\usepackage{textcomp}
\usepackage{xcolor}
\def\BibTeX{{\rm B\kern-.05em{\sc i\kern-.025em b}\kern-.08em
    T\kern-.1667em\lower.7ex\hbox{E}\kern-.125emX}}


\usepackage{amsmath,amssymb,amsthm}
\usepackage[T1]{fontenc}
\usepackage[utf8]{inputenc}
\usepackage[expansion=false]{microtype}
\usepackage[margin=1in]{geometry}
\usepackage{booktabs}
\usepackage{tabularx}
\usepackage{longtable}
\usepackage{array}
\usepackage{multirow}
\usepackage{makecell}
\usepackage{xcolor}
\usepackage{colortbl}
\usepackage{graphicx}
\usepackage{hyperref}
\usepackage{algorithm}
\usepackage{algpseudocode}
\usepackage{cleveref}
\usepackage{enumitem}
\usepackage{caption}
\usepackage{subcaption}
\usepackage{mathtools}
\usepackage[colorinlistoftodos,textsize=small]{todonotes}
\hypersetup{
  colorlinks=true, linkcolor=blue, citecolor=blue, urlcolor=blue
}

% Theorem environments
\newtheorem{proposition}{Proposition}
\newtheorem{lemma}{Lemma}
\newtheorem{definition}{Definition}
\newtheorem{corollary}{Corollary}
\newtheorem{observation}{Observation}
\theoremstyle{remark}
\newtheorem{remark}{Remark}

% Shorthand
\newcommand{\LCB}{\mathrm{LCB}}
\newcommand{\UCB}{\mathrm{UCB}}
\newcommand{\Sep}{\mathrm{Sep}}
\newcommand{\Inv}{\mathrm{Inv}}
\newcommand{\Margin}{\mathrm{Margin}}
\newcommand{\RMC}{\mathrm{RMC}}
\newcommand{\accBase}{b_{\mathrm{acc}}}
\newcommand{\AccDelta}{\mathrm{Acc}_{\Delta}}
\newcommand{\nAcc}{\mathrm{nAcc}}
\newcommand{\RMCSAT}{\mathrm{RMC}_{\mathrm{SAT}}}
\newcommand{\RMCVD}{\mathrm{RMC}_{\mathrm{VD}}}
\newcommand{\thetaP}{\hat{\theta}_P}
\newcommand{\thetaN}{\hat{\theta}_N}
\newcommand{\shat}{\hat{s}_G}
% Semantic names for the shortcut/side-channel adversary notation.
% Change only these display definitions if a different notation is preferred.
\newcommand{\shortcutDisc}{d_{\mathrm{shortcut}}^{G}}
\newcommand{\hatshortcutDisc}{\hat d_{\mathrm{shortcut}}^{G}}
\newcommand{\shortcutHypCls}{\mathcal{H}_{\mathrm{shortcut}}}
\newcommand{\shortcutPosRate}[1]{q_{+}(#1)}
\newcommand{\shortcutNegRate}[1]{q_{-}(#1)}
\newcommand{\shortcutCeiling}{C_{\mathrm{sc}}^{G}}
\newcommand{\hatshortcutCeiling}{\widehat{C}_{\mathrm{sc}}^{G}}
% Backward-compatible aliases used by earlier drafts.
\newcommand{\deltaG}{\shortcutDisc}
\newcommand{\hatdeltaG}{\hatshortcutDisc}
\newcommand{\CG}{\shortcutCeiling}
\newcommand{\hatCG}{\hatshortcutCeiling}
\newcommand{\qpos}[1]{\shortcutPosRate{#1}}
\newcommand{\qneg}[1]{\shortcutNegRate{#1}}
\newcommand{\shortcutPosRateVar}{q_{+}}
\newcommand{\shortcutNegRateVar}{q_{-}}
\newcommand{\Hc}{\shortcutHypCls}
\newcommand{\Ic}{\mathcal{I}}
\newcommand{\pp}{p_{+}}
\newcommand{\nm}{n_{-}}
\newcommand{\reasonFac}{r_{\mathrm{sem}}}
\newcommand{\shortcutFac}{q_{\mathrm{shortcut}}}
\newcommand{\presentationFac}{\pi_{\mathrm{pres}}}
\newcommand{\EvalSet}{\mathcal{D}_{\mathrm{eval}}}
\newcommand{\CatMetric}{M_{\mathrm{cat}}}
\newcommand{\CatScore}{\operatorname{Score}_{\mathrm{cat}}}

\title{%
  \textbf{Reasoning-Margin Certification (RMC):}\\
  A Statistical Framework for Shortcut-Adjusted Reasoning Evaluation\\
  in LLMs via Minimal-Edit Groups and Confidence Bounds
}

\author{
  [Author Names Omitted for Blind Review]
}

\date{}


\begin{document}

\maketitle

\section{Supplementary Roadmap}
This supplementary material contains the material moved out of the ICSE-length main paper: the full problem setup, detailed RMC theory, derivations for Case A and Case B, the table-column mapping, the shortcut-ceiling spectrum, full Case A per-$N$ SAT tables, additional SAT diagnostics, the full vulnerability-detection table, and implementation-level notes.  The main paper contains the problem claim, the RMC core definition, the experimental design summary, the central Case B by-$N$ trend figures, and the central VD results.  The complete Case B, Case C, and Case D per-$N$ SAT tables are provided here for numerical auditability.  Throughout both documents, the roles of the treatments are fixed: Case B (permutation-calibrated estimated ceiling) is the primary shortcut-adjusted certificate for the SAT domains; Case C (complexity-debiased estimated ceiling) is a supplementary sensitivity analysis that first residualizes the final Case-B feature battery against the raw linear complexity covariates $(\text{num\_vars}, \text{num\_clauses}, \text{num\_literals}, \text{file\_size\_bytes})$, then permutation-calibrates the residual max-KS estimate, and finally blends it 50/50 with the Case A reference $\delta=0$; Case D is a separability-only diagnostic that removes both shortcut subtraction and witness gating and reports the factorized separability lower bound directly; Case A (fixed $0.25$ ceiling) is the idealized construction-level reference that defines the boundary $N^*$; and the vulnerability-detection certificate is construction-level Case A because no code-level residual battery is available yet (\Cref{sec:vd_rmc}).

\section{Problem Revealed by Core Claim I and the RMC Response}
\label{sec:setup}

This section states the measurement problem revealed by the paper's first core claim and then derives the design role of RMC.  The claim is a contribution of the paper rather than generic background: ordinary category metrics overestimate reasoning evidence because they count every correct final label in the same way, regardless of whether the label comes from semantic reasoning, shortcut cues, random guessing, class priors, or one-sided sensitivity.  RMC is introduced as the response to this measurement problem.  The rest of the section turns the claim into a formal evaluation problem and explains why RMC uses category outcomes inside matched groups rather than as isolated instance-level scores.
% ============================================================

\subsection{Evaluation Instances, Semantic Factors, and Non-Semantic Factors}
\label{sec:factor}

Following the causal and shortcut-learning view that observed instances can contain both task-relevant semantic factors and spurious factors~\cite{pearl2009causality,arjovsky2019invariant,geirhos2020shortcut,damour2022underspecification}, we model each benchmark instance as the result of a generation process
\begin{equation}\label{eq:formula-01}
x = R(\reasonFac, \shortcutFac, \presentationFac), \qquad y = g(\reasonFac),
\end{equation}
where $\reasonFac$ denotes the semantic factor that determines the intended label, $\shortcutFac$ denotes label-correlated non-semantic information that can support shortcut prediction, and $\presentationFac$ denotes label-preserving presentation variation.  For SAT, $\reasonFac$ is the satisfiability structure of the CNF, including the constraints that make an instance satisfiable or unsatisfiable.  The shortcut factor $\shortcutFac$ may include formula length, clause statistics, literal balance, variable naming, generator template, or prompt artifacts when those cues correlate with the label.  The presentation factor $\presentationFac$ covers clause order, literal order, whitespace, and other formatting choices that change the surface form while preserving the semantic problem.

RMC assigns different roles to these factors.  The reasoning factor $\reasonFac$ defines the semantic contrast that a model should track.  The shortcut factor $\shortcutFac$ defines the shortcut hypothesis class $\shortcutHypCls$ and the shortcut ceiling subtracted by RMC.  The presentation factor $\presentationFac$ defines the controlled presentation protocol under which the RMC margin is interpreted.  Thus, the term \emph{shortcut ceiling} refers to the best group-level separability achievable by a shortcut-only category predictor under the specified shortcut class.  It is defined over $\shortcutFac$.  The presentation factor $\presentationFac$ belongs to the construction and stability protocol: presentation variation is controlled by construction, canonicalization, randomization, and stability checks, while shortcut effects enter the ceiling term.

This separation keeps the statistical target precise.  Merging $\shortcutFac$ and $\presentationFac$ into a single residual factor would place harmless formatting changes inside the shortcut adversary even when the changes have no systematic label relation.  Separating the two factors gives RMC a cleaner interpretation: the inferential correction targets label-correlated non-semantic evidence, while presentation variation is controlled so that category outcomes inside a group remain comparable.  When presentation robustness is a research question, metamorphic perturbations such as clause shuffling, literal reordering, and whitespace changes can be evaluated as an auxiliary stability audit.  The RMC score reported in this paper is a shortcut-adjusted group statistic based on final category outcomes.

The intended evaluation target, which follows from the first core claim, is attribution of the model's success to the semantic factor $\reasonFac$ rather than to $\shortcutFac$, random guessing, class priors, or one-sided sensitivity.  A model can output labels that are often correct under the benchmark distribution while relying on non-semantic evidence.  This distinction is central to the paper.  Accuracy and F1 are computed from final category predictions, so they count every correct category as evidence of ability.  RMC solves the resulting measurement problem by giving category outcomes a stricter evidential role: they are first organized inside matched groups, then converted into a positive and negative arm separability statistic, then compared against a shortcut ceiling under a confidence bound.  Thus, the object of inference is shortcut-adjusted category discrimination under a specified group construction, presentation-control protocol, shortcut class, and confidence level.

\subsection{Core Claim I: Ordinary Category Metrics Overestimate Reasoning Evidence}
\label{sec:ordinary-metrics}

Let $\EvalSet=\{(x_i,y_i)\}_{i=1}^n$ denote the evaluation set, and let
$\hat y_i=f(x_i)$ be the final category predicted by a model.  A conventional
category metric can be written as
\begin{equation}\label{eq:formula-02}
\CatMetric(f,\EvalSet)
=
\CatScore\bigl((y_1,\ldots,y_n),(\hat y_1,\ldots,\hat y_n)\bigr),
\end{equation}
where $\CatMetric$ denotes a category-outcome metric and $\CatScore$ denotes its label-vector scoring rule.  Accuracy, precision, recall, and F1 all have this form.  They observe the final decision vector while leaving the source of the decision unobserved.  This is the same failure mode that motivates behavioral and contrastive evaluations beyond aggregate accuracy~\cite{ribeiro2020checklist,gardner2020contrast}.

\begin{observation}[Category agreement leaves reasoning unidentified]
\label{obs:nonid}
If two models produce the same final category prediction on every instance in
$\EvalSet$, that is,
\begin{equation}\label{eq:formula-03}
f_1(x_i)=f_2(x_i)\qquad\text{for all }i,
\end{equation}
then every metric of the form in \cref{eq:formula-02} assigns them the same
score:
\begin{equation}\label{eq:formula-04}
\CatMetric(f_1,\EvalSet)=\CatMetric(f_2,\EvalSet).
\end{equation}
Consequently, ordinary category metrics assign the same evidential status to a model whose labels are produced by the semantic factor $\reasonFac$ and to a model whose labels are produced by shortcut cues $\shortcutFac$, as long as the two models agree on the observed category vector.
\end{observation}

\begin{proof}
The metric receives only the ground truth category vector and the predicted
category vector.  If two models induce the same predicted vector on
$\EvalSet$, the inputs to $\CatScore$ are identical.  Therefore the returned score
is identical.
\end{proof}

Observation~\ref{obs:nonid} states the paper's first core finding in its simplest form.  Acc and F1 can be high even when the evidence for reasoning is weak.  A correct category may arise because the model solved the semantic constraints, but it may also arise because the model learned a class prior, guessed the majority class, recognized a construction artifact, or became sensitive to only one side of the task.  Ordinary category metrics merge these cases into the same correctness count.  They therefore tend to overestimate reasoning ability when shortcut success and semantic success are both counted as the same kind of correctness.

ADR$^{+w}$ provides a useful reference point for SAT because it requires more than a correct SAT label.  On satisfiable instances, the model must also provide a valid assignment, and the assignment can be checked directly against the CNF, reflecting the classical NP verification view of SAT~\cite{cook2023complexity}.  This is meaningful evidence: a valid assignment shows that the model has found values for the variables that jointly satisfy the clauses, which is stronger than predicting the word ``SAT.''  ADR$^{+w}$ should be treated as a SAT-side witness-aware proxy rather than an absolute definition of reasoning ability.  It checks constructive evidence on the SAT side and leaves a checkable refutation or proof object on the UNSAT side outside the score.  Since verifying a model's understanding of why a CNF is UNSAT is substantially harder than checking a satisfying assignment and typically requires a normalized proof object or refutation checker~\cite{heule2016drat,cruz2017lrat}, ADR$^{+w}$ is best understood as a SAT-side witness-aware proxy, and in this paper it is used as an external sanity check rather than as part of the RMC computation.

This interpretation makes the overestimation claim conservative.  If Acc and F1 are substantially higher than ADR$^{+w}$, then they are higher than a reference that already checks constructive SAT-side evidence while omitting the stricter UNSAT-side proof requirement.  A fully proof-aware SAT metric that required both a satisfying assignment for SAT instances and a valid refutation for UNSAT instances would be no larger than ADR$^{+w}$ under the same predictions.  Thus, large gaps between Acc or F1 and ADR$^{+w}$ support the claim that ordinary category metrics overstate the model's reasoning ability.

\subsection{Core Claim II: RMC Addresses the Overestimation Problem}
\label{sec:rmc-target-setup}

RMC addresses the problem revealed by Core Claim I by changing the evaluation unit and the shortcut-only reference model while leaving assignments outside the RMC formula.  Its group construction is related in spirit to ADR-style paired evaluation, metamorphic testing, behavioral testing, and contrast-set evaluation, but the target is a confidence-bounded shortcut-adjusted margin rather than only pairwise agreement or a pass/fail relation~\cite{chen2020metamorphic,barr2015oracle,ribeiro2020checklist,gardner2020contrast}.  This scope is deliberate.  If a task provides a complete and machine-checkable certificate, such as a satisfying assignment for a SAT instance or a proof object for a theorem, then verifying the certificate is usually the strongest available evidence.  Many evaluations expose only final labels.  A vulnerability detector may output only ``vulnerable'' or ``non-vulnerable'' without a validated exploit, patch, or causal trace.  A bug triage system may output only whether a report is valid.  A flaky test predictor may output only whether a test is flaky.  Even for CNF SAT, assignment checking is naturally available only for satisfiable formulas unless one also elicits and verifies UNSAT refutations.  RMC is designed for this common black-box situation, where the evaluator observes final categories and can build matched groups while complete witnesses remain unavailable or unreliable.

The RMC calculation is therefore based on final category predictions.  Its purpose is to estimate, with a specified confidence level, how much of the observed category discrimination remains after accounting for shortcut-only explanations, random guessing, class priors, and one-sided category behavior.  The presentation factor $\presentationFac$ has a different role: it defines what the group construction must hold stable or normalize so that residual category discrimination can be attributed to $\reasonFac$ after discounting $\shortcutFac$.  This makes RMC a method for shortcut-adjusted reasoning evaluation under a controlled presentation protocol, rather than a single statistic that simultaneously estimates every form of invariance.

The basic intervention is a minimal edit that changes the semantic factor $\reasonFac$ while keeping $\shortcutFac$ matched and $\presentationFac$ controlled as much as the construction permits.  For SAT, a positive-arm instance $x_i^+$ is satisfiable and a matched negative-arm instance $x_i^-$ is unsatisfiable.  These instances are collected into a minimal-edit group $G$ with positive arm $P_G$ and negative arm $N_G$.  ADR uses a pair of highly similar inputs with different labels to test whether a model can distinguish a semantic difference that survives superficial similarity.  RMC extends this idea from a single pair to a group of many such matched pairs.  The group provides repeated observations of the same evaluation question, which supports factorized separability and finite-sample confidence certification.

RMC avoids rewarding a model merely for predicting many positives correctly or many negatives correctly.  It estimates the two arm sensitivities separately and combines them through a factorized separability statistic.  This product form is important because it penalizes degenerate behavior.  An all-SAT predictor may perform well on $P_G$ but fails on $N_G$; an all-UNSAT predictor has the opposite failure; a random or prior-driven predictor has little chance of clearing the shortcut ceiling reliably.

\paragraph{Class imbalance and prior baselines.}
This point is especially important for imbalanced evaluations.  Raw accuracy must be interpreted against the majority-class no-information baseline, but RMC itself does not need to be rescaled by the empirical positive/negative ratio.  The arm sensitivities are conditional quantities: $\hat\theta_P$ is computed within the positive arm and $\hat\theta_N$ within the negative arm.  Therefore, if the model always predicts the majority label, one factor may be large but the other collapses to zero.  Unequal arm sizes affect the width of the CP and Wilson lower bounds through the actual counts $|P_{\mathrm{eval}}|$ and $|N_{\mathrm{eval}}|$; they do not change the target separability $\hat\theta_P\hat\theta_N$.  The fixed Case A ceiling $0.25$ is also conservative under imbalance: any prior-only randomized predictor that outputs the positive label with probability $r$ has expected product $r(1-r)\leq0.25$.  Using the empirical prior product instead would lower the ceiling and make certification easier.

The RMC target is therefore a relative certificate:
\begin{quote}
under the specified minimal-edit construction, presentation-control protocol, shortcut hypothesis class $\shortcutHypCls$, and confidence level, the model's within-group category discrimination should exceed the strongest evaluated shortcut account that lacks access to the semantic factor $\reasonFac$.
\end{quote}
This target is weaker than proving that the model reasons in an absolute sense, but it is stronger than reporting raw category agreement.  It is the kind of claim that can be supported by finite black-box evaluations.

ADR$^{+w}$ and RMC play different roles.  ADR$^{+w}$ audits whether SAT-side correct labels have constructive support from valid assignments.  RMC estimates shortcut-adjusted reasoning evidence from category outcomes alone.  Thus, ADR$^{+w}$ serves as an external diagnostic rather than an input or gate in the RMC formula.  In our SAT results, RMC is typically slightly lower than ADR$^{+w}$, which is the expected pattern if RMC is discounting category success that may be due to shortcut cues, random guessing, class priors, one-sided behavior, or finite-sample uncertainty.  As $N$ increases, the gap between RMC and ADR$^{+w}$ can grow because the model's understanding of the CNF structure becomes weaker.  In particular, a correct UNSAT label may be less reliable evidence of true refutation-like understanding at larger $N$, since the model may predict UNSAT without identifying the contradiction that makes the formula unsatisfiable.  RMC captures this degradation through the negative-arm sensitivity and the confidence-corrected separability statistic, while ADR$^{+w}$ mainly audits constructive evidence on the SAT side.

\subsection{Interventional View and Scope of RMC}
\label{sec:counterfactual-reasoning}

The desired behavior can be stated as an interventional property of the composed system $f\circ R$:
\begin{equation}\label{eq:formula-05}
\mathsf{RE}(f;R)
\approx
\mathsf{Sens}_{\mathrm{sem}}(f)
\wedge
\mathsf{Adj}_{\mathrm{shortcut}}(f \mid \mathcal{C}_{\presentationFac}),
\end{equation}
where $\mathcal{C}_{\presentationFac}$ denotes the benchmark protocol that controls presentation-level variation through matching, canonicalization, randomization, or auxiliary stability checks.  The term $\mathsf{Sens}_{\mathrm{sem}}(f)$ means that when $\reasonFac$ changes under a controlled construction, the prediction $\hat y$ tracks the true label $y$.  The term $\mathsf{Adj}_{\mathrm{shortcut}}(f \mid \mathcal{C}_{\presentationFac})$ means that the observed sensitivity remains above the shortcut ceiling after accounting for shortcut strategies that use $\shortcutFac$ but lack access to $\reasonFac$, with presentation variation controlled by $\mathcal{C}_{\presentationFac}$.

This formulation clarifies the role of the three factors.  The semantic factor $\reasonFac$ defines the target ability.  The shortcut factor $\shortcutFac$ defines the adversarial shortcut-only reference model that RMC subtracts through the shortcut ceiling.  The presentation factor $\presentationFac$ defines the presentation-control condition under which the RMC margin is interpreted.  Consequently, the RMC margin has a precise scope: it estimates reasoning evidence beyond shortcut, guessing, class-prior, and one-sided explanations for category success, under a benchmark construction that keeps presentation variation stable or separately audited.

This section gives the rest of the paper its formal target and connects the first core claim to the second core claim.  Category outcomes alone provide limited evidence about the internal reasoning process of an LLM.  When category outcomes are placed inside minimal-edit groups and compared with a shortcut ceiling using a 95\% lower confidence bound, they provide a conservative estimate of reasoning evidence that avoids the most common overestimation failure of Acc and F1.










% ============================================================



\section{RMC Theory Framework}
\label{sec:theory}
% ============================================================

\subsection{Minimal-Edit Interventions and the Certification Target}
\label{sec:groups}

\paragraph{From ADR pairs to RMC groups.}
The basic unit that motivates our construction is an ADR-style \emph{minimal-edit pair}: a satisfiable instance and an unsatisfiable instance generated from the same seed by a small semantic intervention.  The pair-level view is intuitive: a model succeeds on a pair when it predicts the satisfiable instance as SAT and the unsatisfiable instance as UNSAT.  Such a pair asks whether the model can distinguish two highly similar inputs whose labels differ for a semantic reason.

RMC turns this pair idea into a group evaluation design.  Many minimal-edit pairs generated from the same seed family are collected into one \emph{minimal-edit group}.  In a group, the positive arm collects the satisfiable members from those pairs, and the negative arm collects the corresponding unsatisfiable members.  This group-level formulation lets us estimate the model's positive-arm and negative-arm sensitivities separately and combine them through a factorized separability statistic.  The group view preserves the ADR intuition of paired semantic contrast while adding enough repeated observations to support confidence-bounded certification.

\begin{definition}[Minimal-edit pair and group]
\label{def:group}
A \textbf{minimal-edit pair} is a paired intervention
$(x_i^+,x_i^-)$ generated from a common seed or seed variant, where $x_i^+$ is
satisfiable and $x_i^-$ is unsatisfiable.
The edit is designed to flip the semantic label while holding non-semantic side
channels $\shortcutFac$, such as length, variable count, clause template,
literal balance, and prompt format, as invariant as the construction permits.

A \textbf{minimal-edit group} is a finite collection of such pairs generated
within the same seed family:
\begin{equation}\label{eq:minimal-edit-group}
G=\{(x_i^+,x_i^-)\}_{i=1}^{m_G}.
\end{equation}
Equivalently, the group induces a positive arm
\begin{equation}\label{eq:formula-06}
P_G=\{x_i^+:1\le i\le m_G\}
\end{equation}
and a negative arm
\begin{equation}\label{eq:formula-07}
N_G=\{x_i^-:1\le i\le m_G\}.
\end{equation}
The positive arm $P_G$ contains instances with ground-truth label positive, whereas
the negative arm $N_G$ contains instances with ground-truth label negative.
For evaluation, $P_{\mathrm{eval}}\subseteq P_G$ and
$N_{\mathrm{eval}}\subseteq N_G$ denote instances with valid model predictions.
\end{definition}

This pair-to-group transition is necessary because a single matched pair can illustrate semantic contrast, while a group provides repeated positive-arm and negative-arm observations under the same construction regime.  The repeated observations allow the two semantic sensitivities to be estimated separately and support a stable confidence certificate.

The estimand is a \emph{relative certificate} of reasoning evidence rather than an absolute philosophical quantity.  It asks whether the observed within-group discrimination exceeds the behavior achievable by an explicitly specified class of side-channel-only shortcut predictors that lack access to the semantic factor $\reasonFac$.

\paragraph{Why this is consistent with the critique of class-label metrics.}
RMC uses correctness events derived from model outputs in a stricter evidential role than accuracy or F1.  Accuracy and F1 aggregate isolated
instance-level labels and therefore treat any correct label as equally
informative, regardless of whether it was produced by semantic solving, class
imbalance, random guessing, or a side channel.  RMC treats label
correctness as a raw observation that must pass three additional evidential
filters: matched minimal-edit grouping, conjunctive positive/negative
separability, and a shortcut-adjusted lower-confidence margin.  The RMC formula uses no satisfying assignments.  Consequently, the claim has a precise scope: ordinary metrics computed from category predictions are non-identifying and often overestimate reasoning, whereas RMC uses the same kind of observable category outcomes inside a controlled statistical test designed to avoid overclaiming.

\paragraph{Example.}
Consider the following satisfiable 2-SAT formula:
\begin{equation}\label{eq:formula-08}
F_{\mathrm{SAT}}
=
(x_1)
\wedge
(\neg x_1\vee x_2)
\wedge
(\neg x_2\vee x_3).
\end{equation}
It is satisfiable, for example by the assignment
$x_1=x_2=x_3=\texttt{true}$.
A negative-arm formula can be obtained by a small semantic edit that adds the
clause $\neg x_3$:
\begin{equation}\label{eq:formula-09}
F_{\mathrm{UNSAT}}
=
(x_1)
\wedge
(\neg x_1\vee x_2)
\wedge
(\neg x_2\vee x_3)
\wedge
(\neg x_3).
\end{equation}
This formula is unsatisfiable because $x_1$ forces $x_2$, $x_2$ forces $x_3$,
but the last clause requires $\neg x_3$.

The two formulas differ in the semantic factor $\reasonFac$: one has a
satisfying assignment and the other has an implication-chain contradiction.
At the same time, a minimal-edit construction attempts to keep side-channel
features $\shortcutFac$ similar, such as variable names, clause width, literal
balance, and prompt format.

If a model predicts $F_{\mathrm{SAT}}$ as SAT and $F_{\mathrm{UNSAT}}$ as
UNSAT, the result shows correct discrimination on this minimal-edit pair.
A full group repeats this construction over multiple paired variants from the
same seed family, producing several SAT-side and UNSAT-side instances whose
surface features are intended to remain matched within the group.

However, correctness on one pair, or even high agreement over many pairs, does
not by itself establish that the model used the semantic implication chain
rather than a shortcut feature.
For instance, if negative-arm formulas systematically contain one extra clause
or a distinctive final unit clause, a shortcut-only predictor in
$\shortcutHypCls$ might exploit that pattern without performing SAT reasoning.
The RMC certificate asks whether the model's observed within-group
discrimination exceeds what can be explained by the best side-channel-only
predictor in the specified class $\shortcutHypCls$.




\begin{definition}[Within-group side-channel distinguishability]
\label{def:deltaG}
Consider a minimal-edit group
\begin{equation}\label{eq:formula-10}
G=\{(x_i^+,x_i^-)\}_{i=1}^{m_G}
\end{equation}
with induced arms $(P_G,N_G)$, where $P_G$ contains instances with ground-truth
label SAT and $N_G$ contains instances with ground-truth label UNSAT.
The positive arm and the negative arm differ in the semantic factor that
determines the label, while non-semantic side-channel features are intended to
remain as similar as possible.

Let $\shortcutHypCls$ be a class of shortcut predictors that can observe only the
side-channel information $\shortcutFac$, but not the semantic factor $\reasonFac$.
For a shortcut predictor $h\in\shortcutHypCls$, define its positive-output rates
on the two arms within pair as
\begin{equation}\label{eq:formula-11}
\shortcutPosRate{h}
=
\Pr[h(\shortcutFac)=1\mid Y=1,G],
\end{equation}
and
\begin{equation}\label{eq:formula-12}
\shortcutNegRate{h}
=
\Pr[h(\shortcutFac)=1\mid Y=0,G].
\end{equation}
Equivalently, in a finite group these arm-level rates are estimated by averaging
over instances within each arm:
\begin{equation}\label{eq:formula-13}
\widehat{\shortcutPosRate{h}}
=
\frac{1}{|P_G|}
\sum_{x\in P_G}
\mathbf{1}\{h(\shortcutFac(x))=1\},
\end{equation}
and
\begin{equation}\label{eq:formula-14}
\widehat{\shortcutNegRate{h}}
=
\frac{1}{|N_G|}
\sum_{x\in N_G}
\mathbf{1}\{h(\shortcutFac(x))=1\}.
\end{equation}
Thus, the group first aggregates predictions within the positive and negative
arms separately; it avoids multiplying or crossing all pairs at this stage.

The corresponding negative-output rates are
\begin{equation}\label{eq:formula-15}
1-\shortcutPosRate{h}
=
\Pr[h(\shortcutFac)=0\mid Y=1,G],
\end{equation}
and
\begin{equation}\label{eq:formula-16}
1-\shortcutNegRate{h}
=
\Pr[h(\shortcutFac)=0\mid Y=0,G].
\end{equation}
Here, $1-\shortcutPosRate{h}$ is the probability that the shortcut predictor
outputs the negative label on positive-arm instances, whereas
$1-\shortcutNegRate{h}$ is the probability that it outputs the negative label on
negative-arm instances.

The within-group side-channel distinguishability is
\begin{equation}\label{eq:formula-17}
\shortcutDisc
=
\sup_{h\in\shortcutHypCls}
\left|
\shortcutPosRate{h}-\shortcutNegRate{h}
\right|.
\end{equation}
Thus, $\shortcutDisc$ measures how well the positive and negative arms can be
distinguished using side-channel information alone.
When $\shortcutDisc=0$, the side-channel distribution is indistinguishable across
the two arms: no predictor in $\shortcutHypCls$ can tell positive and negative
instances apart by looking only at $\shortcutFac$.
Larger values of $\shortcutDisc$ indicate stronger residual shortcut signal
within the group.
\end{definition}





\begin{definition}[Shortcut ceiling]
\label{def:shortcutceiling}
For a side-channel-only predictor $h\in\shortcutHypCls$, define its within-group
separability as
\begin{equation}\label{eq:formula-18}
\Sep_G(h)=\shortcutPosRate{h}(1-\shortcutNegRate{h}).
\end{equation}
This is the same product structure used for model separability: the predictor
must output the positive label on positive-arm instances and avoid outputting
the positive label on negative-arm instances.

The shortcut ceiling is the largest separability that can be achieved
without using the semantic factor:
\begin{equation}\label{eq:formula-19}
\shortcutCeiling
=
\sup_{h\in\shortcutHypCls}\Sep_G(h).
\end{equation}
If the two arms differ by at most $\shortcutDisc$ in side-channel space, then
\begin{equation}\label{eq:formula-20}
\shortcutCeiling
\le
\left(\frac{1+\shortcutDisc}{2}\right)^2.
\end{equation}
\end{definition}

\begin{proof}
Fix any side-channel-only predictor $h\in\shortcutHypCls$ and write
$\shortcutPosRateVar=\shortcutPosRate{h}$ and $\shortcutNegRateVar=\shortcutNegRate{h}$.
By the definition of $\shortcutDisc$,
\begin{equation}\label{eq:formula-21}
|\shortcutPosRateVar-\shortcutNegRateVar|\le \shortcutDisc .
\end{equation}
Therefore,
\begin{equation}\label{eq:formula-22}
\shortcutNegRateVar\ge \shortcutPosRateVar-\shortcutDisc .
\end{equation}
For certification, the most favorable case for the shortcut predictor is the
one that makes $\Sep_G(h)=\shortcutPosRateVar(1-\shortcutNegRateVar)$ as large as possible.
Using the inequality above,
\begin{equation}\label{eq:formula-23}
\Sep_G(h)
=
\shortcutPosRateVar(1-\shortcutNegRateVar)
\le
\shortcutPosRateVar(1-\shortcutPosRateVar+\shortcutDisc).
\end{equation}
The right-hand side is a concave quadratic function of $\shortcutPosRateVar$:
\begin{equation}\label{eq:formula-24}
\shortcutPosRateVar(1-\shortcutPosRateVar+\shortcutDisc).
\end{equation}
Its maximum over $\shortcutPosRateVar\in[0,1]$ is attained at
\begin{equation}\label{eq:formula-25}
\shortcutPosRateVar=\frac{1+\shortcutDisc}{2},
\end{equation}
which gives
\begin{equation}\label{eq:formula-26}
\Sep_G(h)
\le
\left(\frac{1+\shortcutDisc}{2}\right)^2.
\end{equation}
Since this bound holds for every $h\in\shortcutHypCls$, it also holds for the supremum over
$\shortcutHypCls$:
\begin{equation}\label{eq:formula-27}
\shortcutCeiling
=
\sup_{h\in\shortcutHypCls}\Sep_G(h)
\le
\left(\frac{1+\shortcutDisc}{2}\right)^2.
\end{equation}
This proves the claim.
\end{proof}

\paragraph{Operational Case-B estimator for $\shortcutCeiling$.}
The empirical Case-B analysis instantiates the shortcut class
$\shortcutHypCls$ as a battery of one-dimensional threshold adversaries over
label-free CNF surface features. For each instance $x$, the implementation
extracts a feature vector $\phi(x)$ from the CNF content only, excluding filename,
directory, group id, pair id, model output, and the ground-truth label. The
pre-registered final feature set used for the reported Case-B ceiling is
\begingroup\small
\begin{equation}
\label{eq:caseb-feature-set}
\begin{aligned}
\mathcal{F}_{\mathrm{CB}}
=\{&
\texttt{mean\_clause\_len},
\texttt{min\_clause\_len},\\
&
\texttt{max\_clause\_len},
\texttt{pos\_ratio},
\texttt{var\_coverage}
\}.
\end{aligned}
\end{equation}
\endgroup
A diagnostic run may additionally record features such as variable count,
clause count, literal counts, distinct-variable count, and file size, but the
reported final Case-B tables use the residual feature set in
\Cref{eq:caseb-feature-set}.

For a group $G$ and feature $f\in\mathcal{F}_{\mathrm{CB}}$, let
\begin{equation}\label{eq:caseb-ks-feature}
D_{G,f}
=
\sup_t
\left|
\widehat{\Pr}[f(x)\le t\mid x\in P_G]
-
\widehat{\Pr}[f(x)\le t\mid x\in N_G]
\right|
\end{equation}
be the empirical two-sample Kolmogorov--Smirnov distance between the SAT and
UNSAT arms. The observed one-feature shortcut distinguishability is
\begin{equation}\label{eq:caseb-delta-observed}
\hat\delta^{\mathrm{obs}}_G
=
\max_{f\in\mathcal{F}_{\mathrm{CB}}} D_{G,f}.
\end{equation}
Because selecting the maximum over features introduces upward bias in small
groups, the implementation uses permutation calibration. It repeatedly permutes
the SAT/UNSAT labels within the same group, recomputes the same max-KS statistic,
and subtracts the $q$-quantile of this permutation reference distribution:
\begin{equation}\label{eq:caseb-delta-calibrated}
\hat\delta_G
=
\max\left\{
0,\,
\hat\delta^{\mathrm{obs}}_G
-
Q_q\!\left(\hat\delta^{\mathrm{perm}}_G\right)
\right\}.
\end{equation}
The reported experiments use $q=0.95$ and 200 permutations by default. If
permutation calibration is disabled, the same estimator reduces to the raw
max-KS value $\hat\delta^{\mathrm{obs}}_G$.

The group-level Case-B shortcut ceiling is then
\begin{equation}\label{eq:caseb-group-shortcut-ceiling}
\hatshortcutCeiling
=
\left(\frac{1+\hat\delta_G}{2}\right)^2 .
\end{equation}
For a model--$N$ slice containing evaluable groups $\mathcal{G}_{\mathrm{eval}}$,
the table entry labeled $\shortcutCeiling$ is the macro-average of these
group-level ceilings:
\begin{equation}\label{eq:caseb-reported-shortcut-ceiling}
\hatshortcutCeiling
=
\frac{1}{|\mathcal{G}_{\mathrm{eval}}|}
\sum_{G\in\mathcal{G}_{\mathrm{eval}}}
\hatshortcutCeiling .
\end{equation}
The reported diagnostics also record the macro-mean $\hat\delta_G$, the
uncalibrated observed max-KS value, the subtracted permutation bias, and the most
frequent top leakage feature. This is why the Case-B ceiling can vary by $N$ and
by model slice: it is computed over the valid groups actually entering that
slice, rather than imposed as a fixed theoretical constant.

The shortcut ceiling $\shortcutCeiling$ is the benchmark against which RMC compares the model.
Standard metrics such as accuracy and F1 compare a model against dataset-level baselines while leaving the source of correctness unresolved.  RMC compares the model's within-group separability against the best separability that remains achievable by side-channel-only predictors.
A positive RMC margin therefore means that the model exceeds the strongest
shortcut-only explanation allowed by the specified adversary class $\shortcutHypCls$.

\subsection{Separability, Pseudo-Replication, and Factorized Inference}
\label{sec:lcb}

For each group, define the two one-sided semantic sensitivities
\begin{equation}\label{eq:formula-28}
\thetaP=\frac{\pp}{|P_{\mathrm{eval}}|},
\qquad
\thetaN=\frac{\nm}{|N_{\mathrm{eval}}|},
\qquad
\shat=\thetaP\thetaN,
\end{equation}
where $\pp$ is the number of SAT instances predicted SAT and $\nm$ is the
number of UNSAT instances predicted UNSAT.
The product is deliberately conjunctive: a model receives credit only when it is
sensitive to satisfiable instances and to unsatisfiable instances in the same
evaluation regime.  All-SAT behavior yields $\thetaP\approx1$ but
$\thetaN\approx0$, while all-UNSAT behavior yields the reverse; either collapses
$\shat$ toward zero even when accuracy or F1 is non-trivial.



\paragraph{How ADR$^{+w}$ relates to RMC.}
The tables report ADR and ADR$^{+w}$ because they provide intuitive pair-level diagnostics.  ADR asks whether the model makes the correct paired category decision on a SAT/UNSAT contrast.  ADR$^{+w}$ is stricter on the SAT side: the SAT-side decision is credited only when the model also supplies a machine-checkable satisfying assignment.  This requirement is meaningful because a valid assignment is a constructive certificate for a satisfiable CNF formula.  The proposed variable values must jointly satisfy all clauses, so the model has done more than merely emit the word ``SAT.''  In this limited but important sense, ADR$^{+w}$ can be read as an approximation to SAT-side CNF understanding.

This approximation is deliberately one-sided.  ADR$^{+w}$ verifies constructive understanding on satisfiable formulas and leaves proof of unsatisfiability outside the witness score.  Although ADR$^{+w}$ still requires the model to label the UNSAT member correctly, the score stops short of requiring a checkable refutation certificate.  This is a pragmatic design choice: satisfying assignments are easy to verify, while UNSAT proofs from LLM outputs are substantially harder to elicit, normalize, and validate.  Therefore ADR$^{+w}$ is best interpreted as a generous SAT-side witness-aware proxy for reasoning evidence.  A fully proof-aware evaluator that required both SAT assignments and UNSAT refutations would be at least as strict and would usually produce a lower rate.

RMC has a different computational role.  It uses category outcomes in matched groups and then applies factorized lower confidence bounds and a shortcut ceiling.  Its purpose is to estimate, with an explicit confidence budget, how much of the observed category success remains after discounting shortcut features, random guessing, class priors, and one-sided sensitivity.  This design is motivated by the black-box setting in which complete certificates are unavailable or unreliable.  When a valid assignment, proof trace, refutation, exploit, or patch can be checked directly, that evidence should be used.  RMC addresses the harder and more common evaluation condition in which the reliable observable is the final category.  ADR$^{+w}$ therefore serves as an external sanity check rather than an input to RMC.  When RMC lies slightly below ADR$^{+w}$, the relationship is conceptually coherent: ADR$^{+w}$ shows that many SAT-side successes have constructive support, while RMC further asks whether the paired category behavior exceeds shortcut and uncertainty baselines.  The fact that Acc and F1 are substantially higher than both quantities indicates that ordinary label metrics overestimate reasoning ability.

The distinction also avoids pseudo-replication.  Expanding a group into all cross-arm SAT--UNSAT comparisons creates the tempting but invalid binomial count $k_G=\pp\nm$ out of $|P_{\mathrm{eval}}||N_{\mathrm{eval}}|$ pairwise comparisons.  Each SAT prediction is reused against every UNSAT prediction, so these cross-arm comparisons share observations.  Treating them as independent shrinks confidence intervals at the wrong rate.  RMC therefore estimates $\theta_P$ and $\theta_N$ as two binomial proportions, multiplies their conservative lower bounds, and then compares the result with the shortcut ceiling.

\begin{definition}[Component lower confidence bounds(LCB)]
\label{def:component_lcb}
For $k$ successes in $n$ independent Bernoulli trials, let
$\hat{\theta}=k/n$ denote the empirical success rate.
We use two alternative one-sided lower confidence bounds for the unknown
success probability $\theta$.

\paragraph{Clopper--Pearson(CP) lower bound\cite{clopper1934use}.}
The exact Clopper--Pearson lower confidence bound at level $\alpha'$ is
\begin{equation}\label{eq:formula-29}
\LCB^{\mathrm{CP}}_{\alpha'}(k,n)
=
\begin{cases}
0, & k=0,\\[2pt]
B^{-1}(\alpha';\,k,\,n-k+1), & 0<k<n,\\[2pt]
(\alpha')^{1/n}, & k=n,
\end{cases}
\end{equation}
where $B^{-1}(\cdot;\,a,b)$ denotes the inverse CDF of a Beta distribution with
shape parameters $a$ and $b$.

\paragraph{Wilson-score lower bound\cite{wilson1927probable,brown2001interval}.}
As a robustness check, we also compute the Wilson-score lower confidence bound:
\begin{equation}\label{eq:formula-30}
\LCB^{\mathrm{W}}_{z}(k,n)
=
\frac{
\hat{\theta}+z^2/(2n)
-
z\sqrt{\hat{\theta}(1-\hat{\theta})/n+z^2/(4n^2)}
}
{1+z^2/n}.
\end{equation}
In the implementation, we set $z=1.96$, corresponding to the usual normal
quantile for a 95\% two-sided interval.
\end{definition}





\begin{definition}[Factorized separability lower confidence bounds]
\label{def:factlcb}
The within-group separability statistic is
\begin{equation}\label{eq:formula-31}
\shat
=
\thetaP\thetaN
=
\frac{\pp}{|P_{\mathrm{eval}}|}
\cdot
\frac{\nm}{|N_{\mathrm{eval}}|},
\end{equation}
where $\pp$ is the number of positive-arm instances predicted positive, and
$\nm$ is the number of negative-arm instances predicted negative.

Rather than treating the $|P_{\mathrm{eval}}|\times |N_{\mathrm{eval}}|$
cross-label comparisons as independent Bernoulli trials, we estimate the two
component success probabilities separately and then multiply their lower
confidence bounds.

\paragraph{CP factorized LCB.}
With $\alpha'=\alpha/2$, the Clopper--Pearson factorized lower bound is
\begin{equation}\label{eq:formula-32}
\LCB^{\mathrm{CP}}_{\alpha}(\shat)
=
\LCB^{\mathrm{CP}}_{\alpha'}(\pp,|P_{\mathrm{eval}}|)
\cdot
\LCB^{\mathrm{CP}}_{\alpha'}(\nm,|N_{\mathrm{eval}}|).
\end{equation}

\paragraph{Wilson factorized LCB.}
The Wilson-score analog is
\begin{equation}\label{eq:formula-33}
\LCB^{\mathrm{W}}_{\alpha}(\shat)
=
\LCB^{\mathrm{W}}_{1.96}(\pp,|P_{\mathrm{eval}}|)
\cdot
\LCB^{\mathrm{W}}_{1.96}(\nm,|N_{\mathrm{eval}}|).
\end{equation}

The CP version is used as the primary conservative estimator, while the Wilson
version is reported as an asymptotic robustness check.
\end{definition}

\begin{proposition}[Coverage of the factorized CP bound]
\label{prop:coverage}
Under independent model evaluations across instances,
$\LCB^{\mathrm{CP}}_{\alpha}(\shat)$ is a conservative finite-sample
$(1-\alpha)$ lower confidence bound on the true separability
\begin{equation}\label{eq:formula-34}
s_G^*=\theta_P^*\theta_N^* .
\end{equation}
\end{proposition}

\begin{proof}
The CP component bounds are exact one-sided binomial lower confidence bounds.
With $\alpha'=\alpha/2$, the positive-arm bound
\begin{equation}\label{eq:formula-35}
\LCB^{\mathrm{CP}}_{\alpha'}(\pp,|P_{\mathrm{eval}}|)
\end{equation}
fails to lower-bound $\theta_P^*$ with probability at most $\alpha/2$, and the
negative-arm bound
\begin{equation}\label{eq:formula-36}
\LCB^{\mathrm{CP}}_{\alpha'}(\nm,|N_{\mathrm{eval}}|)
\end{equation}
fails to lower-bound $\theta_N^*$ with probability at most $\alpha/2$.
By Bonferroni, both component lower bounds hold simultaneously with probability
at least $1-\alpha$.
On that event,
\begin{equation}\label{eq:formula-37}
\LCB^{\mathrm{CP}}_{\alpha'}(\pp,|P_{\mathrm{eval}}|)
\le \theta_P^*
\quad\text{and}\quad
\LCB^{\mathrm{CP}}_{\alpha'}(\nm,|N_{\mathrm{eval}}|)
\le \theta_N^* .
\end{equation}
Since all quantities are nonnegative, multiplying the two component inequalities
yields
\begin{equation}\label{eq:formula-38}
\LCB^{\mathrm{CP}}_{\alpha}(\shat)
\le
\theta_P^*\theta_N^*
=
s_G^* .
\end{equation}
Therefore $\LCB^{\mathrm{CP}}_{\alpha}(\shat)$ is a conservative
$(1-\alpha)$ lower confidence bound on the true separability.~$\blacksquare$
\end{proof}

\begin{remark}[Role of the Wilson factorized bound]
\label{rem:wilson_role}
The Wilson-score factorized bound uses the same product structure:
\begin{equation}\label{eq:formula-39}
\LCB^{\mathrm{W}}_{\alpha}(\shat)
=
\LCB^{\mathrm{W}}_{1.96}(\pp,|P_{\mathrm{eval}}|)
\cdot
\LCB^{\mathrm{W}}_{1.96}(\nm,|N_{\mathrm{eval}}|).
\end{equation}
Unlike the CP bound, however, the Wilson bound is an approximate score-based
confidence bound rather than an exact finite-sample binomial bound.
We therefore use $\LCB^{\mathrm{CP}}_{\alpha}(\shat)$ as the primary
conservative estimator for certification, and report
$\LCB^{\mathrm{W}}_{\alpha}(\shat)$ as an asymptotic robustness check.
Agreement between the CP and Wilson margins indicates that the certification
conclusion remains stable across lower-bound implementations.
\end{remark}

\subsection{Reasoning-Margin Certification}
\label{sec:rmc}

Here, $\alpha$ controls the uncertainty in estimating the model's separability:
smaller $\alpha$ gives a more conservative lower bound.
Similarly, $\beta$ controls the uncertainty in estimating the side-channel
distinguishability: smaller $\beta$ gives a more conservative upper bound on
the shortcut ceiling.

In the Case A reference setting, the shortcut ceiling is fixed by design at
$\shortcutCeiling=0.25$, so the certification uncertainty comes from the
model-separability lower bound.  Setting $\alpha=0.05$ yields a conservative
95\% lower-confidence certificate for $s_G^*$ above the fixed ceiling.  In
Case B, where $\shortcutDisc$ is estimated from an adversary battery, the strict
form of the certificate is relative to both the model-evaluation confidence
budget and the shortcut-estimation confidence budget, as stated in
Proposition~\ref{prop:relid}; the operational form reported in this paper uses
the permutation-calibrated point estimate of $\hatshortcutDisc$ in place of its
upper confidence bound (see Remark~\ref{rem:caseb-variants} after
Definition~\ref{def:rmc}).

\begin{proposition}[Relative identifiability]
\label{prop:relid}
Suppose that two conservative estimates are available for group $G$:
\begin{equation}\label{eq:formula-40}
\LCB_{\alpha}(\shat),
\end{equation}
a lower confidence bound on the model's true within-group separability, and
\begin{equation}\label{eq:formula-41}
\UCB_{\beta}(\hatshortcutDisc),
\end{equation}
an upper confidence bound on the within-group side-channel distinguishability.
If
\begin{equation}\label{eq:formula-42}
\LCB_{\alpha}(\shat)
>
\left(\frac{1+\UCB_{\beta}(\hatshortcutDisc)}{2}\right)^2,
\end{equation}
then, with confidence at least $1-\alpha-\beta$, the model's within-group
discrimination exceeds every side-channel-only predictor in
$\shortcutHypCls$.
Equivalently, the model separates the positive and negative arms better than
the strongest shortcut-only account allowed by the specified adversary class.
\end{proposition}

\begin{proof}
The quantity $\LCB_{\alpha}(\shat)$ lower-bounds the model's true separability
with failure probability at most $\alpha$.
The quantity $\UCB_{\beta}(\hatshortcutDisc)$ upper-bounds the true side-channel
distinguishability with failure probability at most $\beta$.
Therefore, with probability at least $1-\alpha-\beta$, both bounds hold
simultaneously.

On this joint event, the model's true separability is at least
$\LCB_{\alpha}(\shat)$, while the maximum separability achievable by a
side-channel-only predictor is at most
\begin{equation}\label{eq:formula-43}
\left(\frac{1+\UCB_{\beta}(\hatshortcutDisc)}{2}\right)^2 .
\end{equation}
If the former is strictly larger than the latter, then no predictor in $\shortcutHypCls$
that observes only side-channel information can reproduce the model's
within-group discrimination.
This proves the claim.~$\blacksquare$
\end{proof}

\begin{definition}[RMC discrimination margin]
\label{def:rmc}
Let $E(\shat)$ denote the separability estimate used for certification or
reporting. In this paper, $E(\shat)$ can be instantiated in three ways:
\begin{equation}\label{eq:formula-44}
E(\shat)\in
\left\{
\shat,\;
\LCB^{\mathrm{CP}}_{\alpha}(\shat),\;
\LCB^{\mathrm{W}}_{\alpha}(\shat)
\right\}.
\end{equation}
Here, $\shat$ is the raw empirical separability, 
$\LCB^{\mathrm{CP}}_{\alpha}(\shat)$ is the primary conservative CP lower
confidence bound, and $\LCB^{\mathrm{W}}_{\alpha}(\shat)$ is the Wilson-score
robustness check.

Given a shortcut ceiling
\begin{equation}\label{eq:formula-45}
\shortcutCeiling=
\left(\frac{1+\shortcutDisc}{2}\right)^2,
\end{equation}
or its conservative estimated version
\begin{equation}\label{eq:formula-46}
\hatshortcutCeiling=
\left(\frac{1+\UCB_{\beta}(\hatshortcutDisc)}{2}\right)^2,
\end{equation}
the signed RMC margin is
\begin{equation}\label{eq:formula-47}
\RMC_E^{\pm}(G)
=
E(\shat)-\shortcutCeiling .
\end{equation}
When $\shortcutDisc$ is estimated rather than fixed by construction, we use the
conservative version
\begin{equation}\label{eq:formula-48}
\RMC_E^{\pm}(G)
=
E(\shat)-\hatshortcutCeiling .
\end{equation}

The clipped and normalized versions are defined separately as
\begin{equation}\label{eq:formula-49}
\RMC_E^{+}(G)
=
\max(0,\RMC_E^{\pm}(G)),
\end{equation}
and
\begin{equation}\label{eq:formula-50}
\mathrm{nRMC}_E(G)
=
\frac{\RMC_E^{+}(G)}{1-\shortcutCeiling}.
\end{equation}
When the conservative estimated ceiling $\hatshortcutCeiling$ is used, the normalized
score is computed analogously as
\begin{equation}\label{eq:formula-51}
\mathrm{nRMC}_E(G)
=
\frac{\RMC_E^{+}(G)}{1-\hatshortcutCeiling}.
\end{equation}
\end{definition}

The signed margin $\RMC_E^{\pm}$ is the primary certification quantity.
Negative values mean that the model's separability estimate is below the
shortcut ceiling; zero means that no excess over the shortcut ceiling is
certified; positive values mean that the model exceeds the specified
side-channel-only account.

\begin{remark}[Strict versus reported Case B ceiling]\label{rem:caseb-variants}
Two instantiations of the estimated ceiling are used.  The \emph{strict}
instantiation is Eq.~\eqref{eq:formula-46}, which inflates the adversary side
with the one-sided upper confidence bound $\UCB_{\beta}(\hatshortcutDisc)$ and
yields the $1-\alpha-\beta$ joint guarantee of Proposition~\ref{prop:relid}.
The \emph{reported} instantiation, used in all Case B tables and figures of
this paper, replaces $\UCB_{\beta}(\hatshortcutDisc)$ with the
permutation-calibrated point estimate $\hatshortcutDisc$ itself
(\Cref{sec:ceiling_estimation}); the permutation calibration removes the upward
selection bias of the maximum over the feature battery, but the reported
ceiling does not additionally absorb the sampling uncertainty of
$\hatshortcutDisc$.  The strict instantiation is uniformly harder to pass, so
every negative Case B margin in this paper remains negative under the strict
form, while positive Case B margins could shrink.  We state this distinction
explicitly so that the certificate remains auditable, and we recommend the
strict instantiation whenever a positive Case B certificate supports
high-stakes claims.
\end{remark}

% In the main certification setting, we use
% \[
% E(\shat)=\LCB^{\mathrm{CP}}_{\alpha}(\shat),
% \]
% so the margin is a conservative lower-bound comparison:
% \[
% \RMC_{\mathrm{CP}}^{\pm}(G)
% =
% \LCB^{\mathrm{CP}}_{\alpha}(\shat)-\shortcutCeiling .
% \]
% The raw estimate $\shat-\shortcutCeiling$ is reported only as an optimistic diagnostic, and
% the Wilson version is reported as a robustness check.

When $\shortcutDisc=0$, the shortcut ceiling is
\begin{equation}\label{eq:formula-52}
\shortcutCeiling=\left(\frac{1+0}{2}\right)^2=0.25,
\end{equation}
so the signed RMC range is $[-0.25,0.75]$ and the normalized score satisfies
$\mathrm{nRMC}\in[0,1]$.
The normalized score is secondary; it only rescales the positive headroom above
the shortcut ceiling while preserving the signed margin.


\subsection{RMC in SAT Evaluation}
\label{sec:witness}

In the SAT experiments, RMC is used to certify within-group SAT/UNSAT
discrimination.
The certification target is the same as in the general framework: a model is
certified when its conservative separability lower bound exceeds the
shortcut ceiling.
Thus, RMC applies beyond SAT; SAT is one application domain in which
the positive and negative arms correspond to satisfiable and unsatisfiable
instances.

For each group $G$, the signed RMC margin is
\begin{equation}\label{eq:formula-53}
\RMC^{\pm}(G)
=
\underbrace{\LCB^{\mathrm{CP}}_{\alpha}(\shat)-\shortcutCeiling}_{\text{discrimination margin}}.
\end{equation}
The model certifies group $G$ iff $\RMC^{\pm}(G)>0$.

This definition certifies whether the model's within-group discrimination
exceeds the maximum separability attainable by the specified side-channel-only
adversary class.
Because the score factorizes positive-arm and negative-arm correctness,
one-sided label bias fails certification: in the SAT setting, all-SAT
behavior gives $\theta_P\approx 1$ but $\theta_N\approx 0$, while all-UNSAT
behavior gives the reverse.

The SAT witness is reported separately through ADR$^{+w}$ and related audit columns.  This separation is deliberate: RMC evaluates whether category-level SAT/UNSAT behavior exceeds shortcut and guessing baselines at 95\% confidence, while ADR$^{+w}$ asks whether SAT-side successes have constructive assignment support.  Keeping the two quantities separate makes the interpretation clean.  RMC is the shortcut-adjusted category-outcome certificate; ADR$^{+w}$ is a witness-aware reference point for the SAT side.

The implementation exports macro-averaged per-group quantities as
\texttt{RMC\_signed}, \texttt{RMC\_clipped}, and
\texttt{RMC\_certified\_fraction}; the certified capability boundary $N^*$ in
\Cref{tab:com_pern,tab:os_pern} is the largest $N$ whose per-group CP RMC
macro-average is positive.







\subsection{Global-Pooled Estimators Used in the 3-SAT and 2-SAT Tables}
\label{sec:table_estimators}

The main per-$N$ tables report global-pooled discrimination margins; the
boundary $N^*$ is defined from the same global-pooled RMC(CP) column under the
Case A ceiling, while per-group macro-averages are retained as diagnostics.
This choice is deliberate: it gives a stable, comparable view of the
CP-vs.-Wilson-vs.-no-LCB effect at each $N$, while ADR$^{+w}$ separately exposes
the witness-sensitive decision rate.

For each model and $N$, the code pools all valid SAT-side and UNSAT-side
predictions:
\begin{equation}\label{eq:formula-54}
\theta_P^{\mathrm{gp}}
=
\frac{\sum_G p_{+,G}}{\sum_G |P_{\mathrm{eval},G}|},
\qquad
\theta_N^{\mathrm{gp}}
=
\frac{\sum_G n_{-,G}}{\sum_G |N_{\mathrm{eval},G}|}.
\end{equation}
The columns in \Cref{tab:com_pern,tab:os_pern} are then computed as follows.
First define the pooled component lower bounds:
\begin{equation}\label{eq:formula-55}
\ell_P^{\mathrm{CP}}
=
\mathrm{CP}_{0.025}
\!\left(
\sum_G p_{+,G},
\sum_G |P_{\mathrm{eval},G}|
\right),
\end{equation}
\begin{equation}\label{eq:formula-56}
\ell_N^{\mathrm{CP}}
=
\mathrm{CP}_{0.025}
\!\left(
\sum_G n_{-,G},
\sum_G |N_{\mathrm{eval},G}|
\right),
\end{equation}
and the corresponding Wilson-score lower bounds:
\begin{equation}\label{eq:formula-57}
\ell_P^{\mathrm{W}}
=
\mathrm{W}_{1.96}
\!\left(
\sum_G p_{+,G},
\sum_G |P_{\mathrm{eval},G}|
\right),
\end{equation}
\begin{equation}\label{eq:formula-58}
\ell_N^{\mathrm{W}}
=
\mathrm{W}_{1.96}
\!\left(
\sum_G n_{-,G},
\sum_G |N_{\mathrm{eval},G}|
\right).
\end{equation}
The reported RMC columns are:
\begin{equation}\label{eq:formula-59}
\mathrm{RMC}(\mathrm{no\ LCB})
=
\theta_P^{\mathrm{gp}}\theta_N^{\mathrm{gp}}
-
0.25,
\end{equation}
\begin{equation}\label{eq:formula-60}
\mathrm{RMC}(\mathrm{CP})
=
\ell_P^{\mathrm{CP}}\ell_N^{\mathrm{CP}}
-
0.25,
\end{equation}
\begin{equation}\label{eq:formula-61}
\mathrm{RMC}(\mathrm{Wilson})
=
\ell_P^{\mathrm{W}}\ell_N^{\mathrm{W}}
-
0.25.
\end{equation}
Their adjacent $\mathrm{nRMC}$ columns are
\begin{equation}\label{eq:formula-62}
\mathrm{nRMC}
=
\frac{\max(0,\RMC^{\pm})}{0.75},
\end{equation}
under the Case-A ceiling.
This notation matches the analysis scripts:
\texttt{global\_pooled\_RMC\_without\_LCB},
\texttt{global\_pooled\_RMC\_signed}, and
\texttt{wilson\_pooled\_RMC\_signed}.
These pooled bounds treat instance-level outcomes as independent binomial
draws; the clustering induced by seed families and the fixed prompt template is
discussed under threats to validity (\Cref{sec:threats}).

\begin{table*}[t]
\centering
\caption{How the theory maps to the reported 3-SAT and 2-SAT table columns.
The displayed RMC columns are global-pooled discrimination margins.  $N^*$ is
the largest $N$ with positive global-pooled RMC(CP) under the Case A ceiling.  The RMC columns are based on category outcomes; ADR$^{+w}$ is a separate SAT-side witness-aware diagnostic.}
\label{tab:column_map}
\renewcommand{\arraystretch}{1.18}
\begin{tabular}{p{3.0cm}p{5.0cm}p{8.0cm}}
\toprule
Table column & Analysis field & Meaning \\
\midrule
F1$_{\mathrm{SAT}}$, F1$_{\mathrm{UNSAT}}$
& \texttt{sat\_positive\_f1}, \texttt{unsat\_positive\_f1}
& Standard per-class binary-prediction metrics; useful diagnostics but not certification quantities. \\
ADR
& \texttt{ADR}
& Pair-level decision agreement over SAT/UNSAT pairs. \\
ADR$^{+w}$
& \texttt{ADR\_with\_assignment}
& Pair success requiring correct labels and a verified satisfying assignment on the SAT side; this is a SAT-side witness-aware proxy and an external reference for RMC; the RMC formula uses category outcomes. \\
RMC(no LCB)
& \texttt{global\_pooled\_RMC\_without\_LCB}
& Raw pooled separability margin, $\hat{s}^{\mathrm{gp}}-0.25$; optimistic diagnostic. \\
RMC(CP)
& \texttt{global\_pooled\_RMC\_signed}
& Exact CP factorized pooled lower-bound margin; primary reported conservative estimator. \\
RMC(Wilson)
& \texttt{wilson\_pooled\_RMC\_signed}
& Wilson factorized pooled lower-bound margin; asymptotic robustness check. \\
$N^*$
& \texttt{global\_pooled\_RMC\_signed} by $N$
& Largest $N$ with positive global-pooled RMC(CP) under the Case A ceiling; this boundary uses category outcomes and excludes assignment validity.  The per-group macro CP RMC (\texttt{RMC\_signed}) is retained as a per-group diagnostic. \\
\bottomrule
\end{tabular}
\end{table*}

\subsection{Shortcut-Ceiling Spectrum}
\label{sec:ceiling_spectrum}

RMC makes explicit that every evaluation
protocol has a shortcut ceiling.
The reason standard metrics overestimate reasoning is that their ceilings are
high and often uncontrolled, whereas the minimal-edit separability ceiling is
construction-controlled.

\begin{table*}[t]
\centering
\caption{Shortcut-ceiling spectrum.  $\pi$ is the positive-class rate and $\rho$
denotes dataset-level side-channel/label agreement.  The RMC ceiling is the
controlled shortcut baseline induced by the minimal-edit group protocol.}
\label{tab:ceilings}
\renewcommand{\arraystretch}{1.25}
\begin{tabular}{p{3.0cm}p{5.0cm}p{4.0cm}p{4.5cm}}
\toprule
Protocol & Shortcut ceiling & In balanced SAT tables & Consequence \\
\midrule
F1 
& $\ge 2\pi/(1+\pi)$, higher with leakage 
& $\ge 0.667$ 
& Can reward one-sided label bias. \\
Accuracy 
& $\rho$ 
& $\approx0.5$ if balanced 
& Blind to mechanism and witness validity. \\
ADR 
& Pair decision chance/shortcut rate 
& varies by pairing 
& Useful diagnostic, not a certificate. \\
\textbf{RMC separability} 
& $\bigl((1+\shortcutDisc)/2\bigr)^2$ 
& $\mathbf{0.25}$ under Case A 
& Lowest controlled shortcut ceiling. \\
\bottomrule
\end{tabular}
\end{table*}










% ============================================================
\section{Experimental Setup}
\label{sec:setup_exp}
% ============================================================

\subsection{Dataset: Minimal-Edit 3-SAT and 2-SAT Groups}

We use the \textbf{minimal-edit SAT benchmarks} from the SAT Group Evaluation project: 3-SAT groups for the commercial API models and 2-SAT groups for the open-weight models.
Each group $G$ is generated from a seed random CNF formula (3-CNF in the 3-SAT benchmark, 2-CNF in the 2-SAT benchmark) with $N$ variables.
The SAT variant is the seed itself; the UNSAT variant is obtained by a minimal clause modification that flips satisfiability.
Surface features (number of clauses, variable count, literal balance, template style) are held as constant as possible between the SAT and UNSAT instances within each group.

\paragraph{Task assignment and cross-family comparability.}
Commercial models are evaluated on 3-SAT; open-weight models are evaluated on 2-SAT.
Ten of the thirteen open-weight configurations already exhibit degenerate single-label behavior on 2-SAT at every tested size (\Cref{sec:open}), so 2-SAT is the simplest regime in which thinking-mode open-weight models produce a non-degenerate certification range, and 3-SAT would be strictly harder for this family.
Because 2-SAT is polynomial-time decidable while 3-SAT is NP-complete, certified boundaries are comparable only within the same task family, and no cross-family model comparison is made in this paper.

\paragraph{Group sizes evaluated.}
Commercial API models (3-SAT): $N \in \{5, 10, 25, 30, 40, 50, 60\}$ variables.
Open-weight models (2-SAT): $N \in \{3, 5, 8, 10, 15, 20, 25, 30, 40, 50\}$ variables (extended range to probe lower-$N$ behaviour and allow finer-grained certification boundaries).
Each model is evaluated on 5--110 groups per $N$ value, depending on data availability.

\subsection{Models}

\paragraph{Commercial API models (11 configurations).}
\begin{itemize}[leftmargin=*, itemsep=2pt]
\item \textbf{OpenAI}: GPT-4o-mini, GPT-5.4-mini, GPT-5
\item \textbf{Anthropic}: Claude Haiku 4.5, Claude Sonnet 4, Claude Opus 4.7 (standard batch), Claude Opus 4.7 (direct), Claude Opus 4.7 (medium-thinking), Claude Opus 4.7 (full-thinking, limited data)
\item \textbf{DeepSeek}: DeepSeek-Chat, DeepSeek-Reasoner
\end{itemize}

\paragraph{Open-source/open-weight models (13 configurations).}
\begin{itemize}[leftmargin=*, itemsep=2pt]
\item \textbf{Meta LLaMA}: LLaMA-3.1-8B-Instruct, LLaMA-3.1-70B-Instruct
\item \textbf{Google}: Gemma-4-31B-it
\item \textbf{DeepSeek (distilled)}: DeepSeek-R1-Distill-Qwen-32B (two prompt configurations)
\item \textbf{Alibaba Qwen}: Qwen2.5-32B-Instruct; Qwen3-30B-A3B-Instruct (four variants: standard, thinking, extended-$N$, extended-$N$+thinking); Qwen3-14B (thinking); Qwen3-32B (thinking); QwQ-32B (thinking); Qwen3-30B (tool-assisted SCC+CDCL, limited data)
\end{itemize}

All models are evaluated in zero-shot mode with a descriptive prompt.
Thinking-mode variants use extended chain-of-thought reasoning budgets.
For each instance, we record the SAT/UNSAT prediction.  When the model predicts SAT, we also record the proposed variable assignment for the separate ADR$^{+w}$ witness audit.

\paragraph{Successful open-weight 2-SAT coverage actually used by the metrics.}
Because the open-weight runs are not perfectly rectangular across all providers and all $N$ values, the effective sample sizes used by the reported 2-SAT metrics are the \emph{successful} complete-pair counts rather than the nominal requested workload.  \Cref{tab:open_2sat_success_counts} reports, for each supplementary-paper 2-SAT model and each tested $N$, the number of groups with at least one complete pair, the total number of complete pairs, the corresponding number of evaluated samples (two per complete pair), and the mean complete pairs/samples per successful group.  Failed requests and records without usable final results are excluded from these counts.

\setlength{\LTleft}{0pt}
\setlength{\LTright}{0pt}
\begin{longtable}{llrrrrr}
\caption{Successful open-weight 2-SAT evaluation coverage used by the reported metrics.
``Successful groups'' counts groups with at least one complete SAT/UNSAT pair having usable final predictions.
``Successful pairs'' and ``Successful samples'' count only complete returned results, excluding failed requests and missing-result records.}
\label{tab:open_2sat_success_counts}\\
\toprule
Model & $N$ & Successful groups & Successful pairs & Successful samples & Mean pairs/group & Mean samples/group \\
\midrule
\endfirsthead
\toprule
Model & $N$ & Successful groups & Successful pairs & Successful samples & Mean pairs/group & Mean samples/group \\
\midrule
\endhead
\bottomrule
\endfoot
Llama-Nemotron-Super-49B (r-on) & 3 & 2 & 22 & 44 & 11.0 & 22.0 \\
Llama-Nemotron-Super-49B (r-on) & 5 & 2 & 22 & 44 & 11.0 & 22.0 \\
Llama-Nemotron-Super-49B (r-on) & 8 & 2 & 22 & 44 & 11.0 & 22.0 \\
Llama-Nemotron-Super-49B (r-on) & 10 & 2 & 22 & 44 & 11.0 & 22.0 \\
Llama-Nemotron-Super-49B (r-on) & 15 & 2 & 22 & 44 & 11.0 & 22.0 \\
Llama-Nemotron-Super-49B (r-on) & 20 & 2 & 22 & 44 & 11.0 & 22.0 \\
Llama-Nemotron-Super-49B (r-on) & 25 & 2 & 22 & 44 & 11.0 & 22.0 \\
Llama-Nemotron-Super-49B (r-on) & 30 & 2 & 22 & 44 & 11.0 & 22.0 \\
Llama-Nemotron-Super-49B (r-on) & 40 & 2 & 22 & 44 & 11.0 & 22.0 \\
Llama-Nemotron-Super-49B (r-on) & 50 & 2 & 22 & 44 & 11.0 & 22.0 \\
Llama-Nemotron-Super-49B (r-on) & 60 & 2 & 13 & 26 & 6.5 & 13.0 \\
Nemotron-H-47B (r-on) & 3 & 86 & 946 & 1892 & 11.0 & 22.0 \\
Nemotron-H-47B (r-on) & 5 & 86 & 946 & 1892 & 11.0 & 22.0 \\
Nemotron-H-47B (r-on) & 8 & 86 & 946 & 1892 & 11.0 & 22.0 \\
Nemotron-H-47B (r-on) & 10 & 86 & 946 & 1892 & 11.0 & 22.0 \\
Nemotron-H-47B (r-on) & 15 & 86 & 946 & 1892 & 11.0 & 22.0 \\
Nemotron-H-47B (r-on) & 20 & 85 & 935 & 1870 & 11.0 & 22.0 \\
Nemotron-H-47B (r-on) & 25 & 85 & 935 & 1870 & 11.0 & 22.0 \\
Nemotron-H-47B (r-on) & 30 & 85 & 935 & 1870 & 11.0 & 22.0 \\
Nemotron-H-47B (r-on) & 40 & 85 & 935 & 1870 & 11.0 & 22.0 \\
Nemotron-H-47B (r-on) & 50 & 85 & 935 & 1870 & 11.0 & 22.0 \\
Nemotron-H-47B (r-on) & 60 & 85 & 935 & 1870 & 11.0 & 22.0 \\
QwQ-32B (think) & 3 & 9 & 99 & 198 & 11.0 & 22.0 \\
QwQ-32B (think) & 5 & 9 & 99 & 198 & 11.0 & 22.0 \\
QwQ-32B (think) & 8 & 9 & 99 & 198 & 11.0 & 22.0 \\
QwQ-32B (think) & 10 & 9 & 99 & 198 & 11.0 & 22.0 \\
QwQ-32B (think) & 15 & 9 & 99 & 198 & 11.0 & 22.0 \\
QwQ-32B (think) & 20 & 9 & 99 & 198 & 11.0 & 22.0 \\
QwQ-32B (think) & 25 & 9 & 99 & 198 & 11.0 & 22.0 \\
QwQ-32B (think) & 30 & 9 & 99 & 198 & 11.0 & 22.0 \\
QwQ-32B (think) & 40 & 9 & 90 & 180 & 10.0 & 20.0 \\
QwQ-32B (think) & 50 & 8 & 88 & 176 & 11.0 & 22.0 \\
QwQ-32B (think) & 60 & 8 & 88 & 176 & 11.0 & 22.0 \\
Qwen3-14B (think) & 3 & 14 & 154 & 308 & 11.0 & 22.0 \\
Qwen3-14B (think) & 5 & 14 & 154 & 308 & 11.0 & 22.0 \\
Qwen3-14B (think) & 8 & 14 & 154 & 308 & 11.0 & 22.0 \\
Qwen3-14B (think) & 10 & 14 & 154 & 308 & 11.0 & 22.0 \\
Qwen3-14B (think) & 15 & 14 & 154 & 308 & 11.0 & 22.0 \\
Qwen3-14B (think) & 20 & 14 & 154 & 308 & 11.0 & 22.0 \\
Qwen3-14B (think) & 25 & 14 & 154 & 308 & 11.0 & 22.0 \\
Qwen3-14B (think) & 30 & 14 & 154 & 308 & 11.0 & 22.0 \\
Qwen3-14B (think) & 40 & 13 & 143 & 286 & 11.0 & 22.0 \\
Qwen3-14B (think) & 50 & 13 & 143 & 286 & 11.0 & 22.0 \\
Qwen3-14B (think) & 60 & 13 & 143 & 286 & 11.0 & 22.0 \\
Qwen3-32B (think) & 3 & 11 & 121 & 242 & 11.0 & 22.0 \\
Qwen3-32B (think) & 5 & 11 & 121 & 242 & 11.0 & 22.0 \\
Qwen3-32B (think) & 8 & 11 & 121 & 242 & 11.0 & 22.0 \\
Qwen3-32B (think) & 10 & 11 & 121 & 242 & 11.0 & 22.0 \\
Qwen3-32B (think) & 15 & 11 & 121 & 242 & 11.0 & 22.0 \\
Qwen3-32B (think) & 20 & 11 & 121 & 242 & 11.0 & 22.0 \\
Qwen3-32B (think) & 25 & 11 & 120 & 240 & 10.9 & 21.8 \\
Qwen3-32B (think) & 30 & 11 & 121 & 242 & 11.0 & 22.0 \\
Qwen3-32B (think) & 40 & 11 & 121 & 242 & 11.0 & 22.0 \\
Qwen3-32B (think) & 50 & 11 & 119 & 238 & 10.8 & 21.6 \\
Qwen3-32B (think) & 60 & 11 & 114 & 228 & 10.4 & 20.7 \\
gpt-oss-20b (reason) & 3 & 3 & 33 & 66 & 11.0 & 22.0 \\
gpt-oss-20b (reason) & 5 & 3 & 33 & 66 & 11.0 & 22.0 \\
gpt-oss-20b (reason) & 8 & 3 & 33 & 66 & 11.0 & 22.0 \\
gpt-oss-20b (reason) & 10 & 3 & 33 & 66 & 11.0 & 22.0 \\
gpt-oss-20b (reason) & 15 & 3 & 33 & 66 & 11.0 & 22.0 \\
gpt-oss-20b (reason) & 20 & 3 & 33 & 66 & 11.0 & 22.0 \\
gpt-oss-20b (reason) & 25 & 3 & 33 & 66 & 11.0 & 22.0 \\
gpt-oss-20b (reason) & 30 & 3 & 33 & 66 & 11.0 & 22.0 \\
gpt-oss-20b (reason) & 40 & 3 & 26 & 52 & 8.7 & 17.3 \\
gpt-oss-20b (reason) & 50 & 2 & 22 & 44 & 11.0 & 22.0 \\
gpt-oss-20b (reason) & 60 & 2 & 22 & 44 & 11.0 & 22.0 \\
\end{longtable}

\subsection{RMC Computation Protocol}

\begin{algorithm}[t]
\small
\caption{$\RMCSAT$ Computation for Group $G$}
\label{alg:rmc}
\begin{algorithmic}[1]
\Require $G=(P_{\mathrm{eval}},N_{\mathrm{eval}},\text{pairs})$; 
         confidence level $\alpha=0.05$
\Require Side-channel distinguishability $\shortcutDisc$
         \Statex \hspace{\algorithmicindent} Case A uses $\shortcutDisc=0$

\State $\pp \gets |\{x\in P_{\mathrm{eval}}:\hat{y}(x)=\mathrm{SAT}\}|$
\State $\nm \gets |\{x\in N_{\mathrm{eval}}:\hat{y}(x)=\mathrm{UNSAT}\}|$

\State $\thetaP \gets \pp/|P_{\mathrm{eval}}|$
\State $\thetaN \gets \nm/|N_{\mathrm{eval}}|$
\State $\shat \gets \thetaP\thetaN$

\State $\ell_P \gets \mathrm{CP}_{\alpha/2}(\pp,|P_{\mathrm{eval}}|)$
\State $\ell_N \gets \mathrm{CP}_{\alpha/2}(\nm,|N_{\mathrm{eval}}|)$
\Statex \hspace{\algorithmicindent} where CP is the Clopper--Pearson lower bound

\State $\ell_{\mathrm{fact}} \gets \ell_P\ell_N$
\State $\shortcutCeiling \gets \bigl((1+\shortcutDisc)/2\bigr)^2$
\Statex \hspace{\algorithmicindent} In Case A, $\shortcutCeiling=0.25$

\State $d_G \gets \ell_{\mathrm{fact}}-\shortcutCeiling$

\State $\RMCSAT(G) \gets d_G$
\State \textbf{certify} $G$ iff $\RMCSAT(G)>0$
\Statex \hspace{\algorithmicindent} Assignment validity is evaluated separately by ADR$^{+w}$, not by RMC.
\end{algorithmic}
\end{algorithm}

\paragraph{Confidence parameters.}
We use $\alpha = 0.05$ throughout (95\% joint coverage with Bonferroni split $\alpha'=0.025$ per component).
We compute three SAT ceiling treatments plus one separability-only diagnostic with fixed roles.  Case A fixes $\shortcutDisc=0$ under the idealized assumption that the minimal-edit construction fully decouples $\shortcutFac$ from $\reasonFac$; it yields the fixed $0.25$ ceiling and defines the construction-level boundary $N^*$.  Case B drops this assumption and estimates $\hatshortcutDisc$ from the reported \emph{final residual} side-channel adversary battery with permutation calibration (\Cref{sec:ceiling_estimation}); Case B is the primary shortcut-adjusted certificate reported in the main paper for the SAT domains.  Case C keeps the same ceiling transform as Case B but first residualizes that same final battery against generic structural-complexity covariates, then permutation-calibrates the residual max-KS estimate, and finally blends it 50/50 with the Case A reference $\delta=0$; Case C is therefore a supplementary complexity-debiased sensitivity analysis rather than a replacement for Case B.  Case D sets the ceiling to zero and removes witness gating, so the reported quantity is the factorized separability lower bound itself; it is intentionally easier than Case A--C and is used only as a separability-only diagnostic.  The VD experiments use Case A because a code-level residual battery is not yet available (\Cref{sec:vd_rmc}).

\paragraph{Aggregate metrics.}
For per-group diagnostics, we report the macro-mean of $\RMCSAT^{\pm}$ over groups with sufficient data ($|P_{\mathrm{eval}}|>0$ and $|N_{\mathrm{eval}}|>0$) and the fraction of groups certified ($\RMCSAT^{\pm}>0$).  The \textbf{certified capability boundary} $N^*$ is defined as the largest $N$ with positive global-pooled RMC(CP) under the Case A ceiling; the corresponding Case B boundary applies the same criterion with the permutation-calibrated estimated ceiling (\Cref{sec:ceiling_estimation}).
For the compact per-$N$ tables, we additionally report the global-pooled discrimination margins defined in \Cref{sec:table_estimators}: RMC(no LCB), RMC(CP), and RMC(Wilson), each with its normalized $\mathrm{nRMC}$ companion.
For Case D specifically, $\shortcutCeiling=0$, so the normalized quantities are numerically identical to their unclipped counterparts; the dedicated Case D tables therefore report the separability estimate and the two lower bounds directly rather than repeating redundant normalized columns.

\paragraph{Majority-adjusted accuracy.}
For ordinary accuracy on an imbalanced evaluation set, let $n_{+}$ and $n_{-}$ be the numbers of positive and negative instances and $N=n_{+}+n_{-}$.  The no-information accuracy baseline is the majority-class rate
\begin{equation}\label{eq:acc-baseline-supp}
\accBase=\frac{\max(n_{+},n_{-})}{N}.
\end{equation}
We therefore interpret raw accuracy through
\begin{equation}\label{eq:acc-delta-supp}
\AccDelta=\mathrm{Acc}-\accBase,
\end{equation}
or, when a normalized scale is useful,
\begin{equation}\label{eq:nacc-supp}
\nAcc=\frac{\mathrm{Acc}-\accBase}{1-\accBase}.
\end{equation}
For balanced binary datasets this reduces to subtracting $0.5$.  For imbalanced datasets it prevents an always-majority predictor from being mistaken for a capable classifier.  This adjustment is only for interpreting conventional accuracy; the RMC statistics remain the arm-normalized quantities defined above.

\subsection{Shortcut-Ceiling Estimation Protocol (Case B)}
\label{sec:ceiling_estimation}
% NOTE TO AUTHORS: please verify that the description below matches the
% implementation in the ceiling-estimation scripts (feature set, KS statistic,
% number of permutations, and per-$N$ aggregation).

The reported Case B ceiling is estimated from the \emph{final residual} CNF side-channel battery used by the current code path: mean clause length, minimum clause length, maximum clause length, literal-polarity ratio, and variable coverage.
For each feature, the battery instantiates the family of one-dimensional threshold adversaries; the strongest threshold adversary's arm distinguishability equals the Kolmogorov--Smirnov (KS) statistic between the positive-arm and negative-arm feature distributions within the group construction at a given $N$.
The raw estimate $\hat{d}_{\mathrm{raw}}$ takes the maximum of these KS statistics over that final battery.
Because a maximum over several in-sample statistics is upward biased, $\hat{d}_{\mathrm{raw}}$ is calibrated against a permutation null: arm labels are permuted within matched pairs, the max-over-features statistic is recomputed for each permutation, and the calibrated estimate $\hatshortcutDisc$ corrects $\hat{d}_{\mathrm{raw}}$ by the permutation-null level of the same max statistic.
The reported per-$N$ residual ceiling is therefore
\[
\shortcutCeiling
=
\left(\frac{1+\hatshortcutDisc}{2}\right)^2-0.25,
\]
macro-averaged over groups at that $N$.
Certification in the reported Case B tables compares the global-pooled factorized LCB against this calibrated residual ceiling; the strict variant of Remark~\ref{rem:caseb-variants} would additionally replace $\hatshortcutDisc$ with $\UCB_{\beta}(\hatshortcutDisc)$.  The older all-feature max-KS estimator, which still exists in the code as a \texttt{diagnostic} leakage audit, is not the reported certificate.

Because $\shortcutCeiling=\tfrac{1}{4}+\shortcutDisc/2+{\shortcutDisc}^{2}/4$ always contains the $0.25$ chance floor of the product statistic, calibrated ceilings of $0.30$--$0.41$ on 3-SAT correspond to residual distinguishability $\hatshortcutDisc\approx0.10$--$0.28$, and 2-SAT ceilings up to $0.554$ correspond to $\hatshortcutDisc\approx0.49$.
These nonzero values quantify how far the minimal-edit construction falls short of the Case A idealization, and they are the empirical reason that Case B, rather than Case A, is the primary reported certificate.

\subsection{Complexity-Debiased Ceiling Protocol (Case C)}
\label{sec:ceiling_estimation_casec}

Case C keeps the same residual ceiling transform as Case B but changes the estimator for $\hatshortcutDisc$.
The input feature battery is \emph{not} expanded: Case C starts from the same final Case-B residual battery, so the only change is how that battery is debiased before the shortcut adversary is evaluated.
For each feature $f_k$ in that battery, the code fits a pooled per-feature ridge regression
\[
f_k \sim 1 + (\texttt{num\_vars},\texttt{num\_clauses},\texttt{num\_literals},\texttt{file\_size\_bytes})
\]
across all instances with the required covariates.
The fitted value is subtracted instance-wise, and the resulting residualized feature replaces the raw feature value before the same max-KS adversary is run.
The implementation deliberately stops at this raw linear design: higher-order, interaction, and log-expanded controls are not part of the reported Case C estimator.

The residualized max-KS score is then permutation-calibrated exactly as in Case B.
If $\hatshortcutDisc_{\mathrm{obs,resid}}$ is the observed residualized max-KS statistic and $\hat b_{\mathrm{perm,resid}}$ is the permutation-null bias estimate, the calibrated residualized distinguishability is
\[
\hatshortcutDisc_{\mathrm{FWL,cal}}
=
\max\!\bigl(0,\hatshortcutDisc_{\mathrm{obs,resid}}-\hat b_{\mathrm{perm,resid}}\bigr).
\]
The reported Case C estimate is not this calibrated residualized value itself, but its 50/50 blend with the Case A construction-level reference $\delta=0$:
\[
\hatshortcutDisc^{(\mathrm{C})}
=
\tfrac{1}{2}\hatshortcutDisc_{\mathrm{FWL,cal}}.
\]
The conceptual difference from Case B is therefore precise: Case B charges the full calibrated residual leakage captured by the final shortcut battery to the shortcut account, whereas Case C charges only the component of that same leakage that remains after partialing out the part linearly predictable from generic structural growth, and then shrinks that residualized estimate halfway toward the Case A idealization.  This is why Case C is typically less punitive than Case B, yet still stricter than pure Case A.

\subsection{VD Domain: Dataset}
\label{sec:vd_dataset}

We apply RMC to a second domain---\textbf{vulnerability detection} (VD)---to demonstrate that the framework generalizes beyond SAT to settings without efficiently verifiable witnesses.

\paragraph{BigVul benchmark.}
We use the BigVul dataset, a real-world software vulnerability benchmark drawn from open-source C/C++ projects.
Each evaluation instance is a \emph{vulnerability pair}: a function version \emph{before} the security patch (\texttt{func\_before}, label $y=1$, vulnerable) and the corresponding patched version \emph{after} the fix (\texttt{func\_after}, label $y=0$, non-vulnerable).
The test split contains \textbf{1,055 vulnerability pairs}, each sharing a CVE identifier.
This structure directly satisfies the RMC group construction: the before/after variants differ only in semantically meaningful patch content, holding almost all surface features constant ($\shortcutDisc\approx 0$).

\paragraph{Group formats.}
We evaluate two complementary scenarios for each model:
\begin{itemize}[leftmargin=*, itemsep=2pt]
\item \textbf{1 pair/group} (\texttt{bigvul\_test\_csv}): The standard BigVul test split, with each CVE pair forming its own group ($|P_{\mathrm{eval}}|=1$, $|N_{\mathrm{eval}}|=1$), for 1,055 total pairs.
\item \textbf{10 pairs/group} (\texttt{group\_test}): An augmented format where each base vulnerability is paired with 10 redundant-statement-augmented variants, yielding $|P_{\mathrm{eval}}|=10$, $|N_{\mathrm{eval}}|=10$ per group, for 10,550 total pairs.
\end{itemize}

\subsection{VD Domain: Models}
\label{sec:vd_models}

We evaluate \textbf{27 fine-tuned LLaMA-3.1-8B variants} across the two scenarios (49 model/scenario pairs total).
All models are fine-tuned from \texttt{LLaMA-3.1-8B-Instruct} on subsets of the BigVul training set, differing in training data composition:

\begin{itemize}[leftmargin=*, itemsep=2pt]
\item \textbf{\texttt{add\_after\_ratio\_*}} (20 scenario pairs): Balanced before/after training with varied vulnerable-to-non-vulnerable ratios; multiple seeds (v01--v10) and two base ratio variants.
\item \textbf{\texttt{add\_after\_percent\_70}} (2 pairs): 70\% non-vulnerable training instances.
\item \textbf{\texttt{add\_after\_10epoch}} (2 pairs): Base balanced model trained for 10 epochs.
\item \textbf{\texttt{add\_after\_percent\_50}} (2 pairs): Equal proportions of before/after (borderline case).
\item \textbf{\texttt{func\_before\_*}} (14 pairs): Trained exclusively on pre-patch (vulnerable) functions, inducing a non-vulnerable prediction bias.
\end{itemize}

\subsection{VD-Specific RMC: Global Aggregate Factorized LCB}
\label{sec:vd_rmc}

The VD domain differs from SAT in two key respects.

\paragraph{No witness gate.}
Vulnerability lacks the efficient certificate structure available for satisfying assignments: practical evaluation typically lacks a polynomial-time certificate for the presence of a vulnerability.
We therefore omit the witness gate, and $\RMCVD^{\pm}$ reduces to the discrimination margin alone:
\begin{equation}\label{eq:formula-63}
\RMCVD^{\pm} = \LCB_{\mathrm{fact}}^{\mathrm{global}} - \shortcutCeiling.
\end{equation}

\paragraph{Why per-group LCB fails for singleton groups.}
With $|P_{\mathrm{eval}}|=1$ per group (1 pair/group scenario), the Clopper--Pearson LCB at $\alpha'=0.025$ is at most $\approx0.025$ regardless of the observed outcome.
Macro-averaging across 1,055 such groups yields $\LCB_{\mathrm{fact}}\approx 0.025^2\approx 0.0006$---far below the $0.25$ shortcut ceiling for any model.
Per-group RMC is uncertifiable by construction for single-pair groups, irrespective of model quality.

\paragraph{Global aggregate approach.}
We instead pool all pairs across groups into a single aggregate count and apply the factorized LCB once:
\begin{equation}\label{eq:formula-64}
\LCB_{\mathrm{fact}}^{\mathrm{global}}
= \mathrm{CP}_{1-\alpha/2}\!\Bigl(\textstyle\sum_G p_{+,G},\;\sum_G |P_G|\Bigr)
\;\times\;
\mathrm{CP}_{1-\alpha/2}\!\Bigl(\textstyle\sum_G n_{-,G},\;\sum_G |N_G|\Bigr).
\end{equation}
This is appropriate when groups are too small to carry individual LCBs: the correct estimation unit is the individual prediction, and pooling pools evidence before applying a single bound.
With $\alpha=0.05$, $\shortcutDisc=0$ (Case A), the VD margin is:
\begin{equation}\label{eq:formula-65}
\RMCVD^{\pm} = \LCB_{\mathrm{fact}}^{\mathrm{global}} - 0.25.
\end{equation}

\paragraph{Case A scope of the VD certificate.}
The VD certificate is construction-level (Case A).
The before/after construction holds surface features closely matched, but we have not yet estimated a code-level residual side-channel battery for this domain; candidate features include function-length and token-count deltas, diff-size cues, comment and identifier statistics, and patch-pattern fingerprints.
Until such a battery is calibrated, VD certificates should be read with the same caveat that applies to the idealized Case A SAT reference, and constructing a VD Case B ceiling is immediate future work.

% ============================================================

\section{Results by Research Question}
\label{sec:results}
% ============================================================

\subsection{RQ1: Conventional class-label metrics overestimate reasoning ability}
\label{sec:rq1}

RQ1 asks whether conventional metrics computed from isolated final labels provide reliable evidence of reasoning.  Across SAT and vulnerability detection, accuracy and F1 often remain moderate or high precisely when SAT-side witness-aware ADR$^{+w}$ and RMC indicate little certified reasoning evidence.  This is the central negative claim of the paper: ordinary label metrics systematically overestimate reasoning when they reward label agreement produced by one-sided behavior, shortcuts, class priors, or random guessing.

The theoretical reason follows from Observation~\ref{obs:nonid}.  Accuracy and F1 are functions of the confusion matrix.  The confusion matrix records how many labels were predicted correctly, while leaving unrecorded whether the prediction was produced by semantic reasoning, label bias, or a shortcut correlated with the label.  A model that always predicts SAT on a balanced SAT benchmark obtains around $0.50$ accuracy and a nonzero SAT-class F1.  A model that recognizes only UNSAT-like surface patterns can obtain high F1$_{\mathrm{UNSAT}}$ while failing to construct satisfying assignments and failing to discriminate matched positive instances.  Thus, conventional class-label metrics are non-identifying proxies for reasoning.

The SAT results make this gap visible.  In \Cref{tab:com_pern}, GPT-5 at $N=40$ has Acc $=0.645$ and F1$_{\mathrm{UNSAT}}=0.735$, yet ADR$^{+w}=0.164$ and RMC(CP) $=-0.040$.  DeepSeek-Reasoner at $N=25$ has Acc $=0.626$ and F1$_{\mathrm{UNSAT}}=0.716$, yet ADR$^{+w}=0.144$ and RMC(CP) $=-0.053$.  ADR$^{+w}$ is used here as a stricter SAT-side witness-aware proxy and a generous reference point rather than an absolute definition of reasoning ability.  It checks constructive understanding for satisfiable formulas while omitting a comparable UNSAT refutation.  The point is therefore conservative: if ordinary label metrics greatly exceed ADR$^{+w}$, they are already higher than a permissive witness-aware proxy for SAT reasoning evidence; if they also exceed RMC, they additionally fail the confidence-bounded shortcut test.  At larger $N$, several models retain nonzero accuracy because one label remains easy to predict or because the benchmark is balanced, but RMC collapses once one of the two semantic sensitivities collapses.

The same phenomenon appears in vulnerability detection.  In \Cref{tab:vd}, vulnerable-only training yields $\theta_P$ above $0.92$ but $\theta_N$ below $0.09$.  Such models appear competent under one-sided recall but almost always predict vulnerable.  RMC treats this as a failure of semantic discrimination because the factorized statistic requires both vulnerable and non-vulnerable recognition.

\subsection{RQ2: RMC gives a 95\% confidence estimate without treating category accuracy as reasoning}
\label{sec:rq2}

RQ2 asks whether RMC can estimate reasoning-like discrimination beyond shortcut behavior at approximately 95\% confidence without repeating the overestimation problem of Acc and F1.  Under Case A, the shortcut ceiling is fixed at $0.25$ and the CP factorized lower bound uses $\alpha=0.05$.  A positive RMC(CP) margin therefore means that, with conservative 95\% lower-confidence coverage for the separability term, the model's matched within-group category discrimination exceeds the fixed shortcut ceiling.  RMC remains output-based and category-outcome-based while adding the design elements that ordinary category metrics lack: matched groups, a shortcut ceiling, and a confidence-bounded test against shortcut, prior, and guessing explanations.  ADR$^{+w}$ is used in this section only as an independent witness-aware reference point for the SAT side.

We evaluate 11 commercial API configurations on 3-SAT instances across $N\in\{5,10,25,30,40,50,60\}$ variables.  Most commercial configurations receive no certificate at any tested size.  The certified configurations show a phase boundary: performance is strong at small $N$, then drops abruptly as the formulas become harder.  This boundary is the core empirical artifact produced by RMC.

\begin{table*}[t]
\centering
\caption{Per-$N$ RMC metrics for the three certified commercial reasoning models on 3-SAT.
ADR$^{+w}$: pair-level ADR with verified satisfying assignment on the SAT side; this is reported separately from RMC.
F1$_{\mathrm{SAT}}$/F1$_{\mathrm{UNSAT}}$: per-class F1 scores.
RMC(no LCB): global pooled point estimate without LCB.
RMC(CP): global pooled Clopper--Pearson LCB ($\alpha=0.05$).
RMC(Wilson): global pooled Wilson-score LCB.
$\mathrm{nRMC}=\max(0,\RMC^{\pm})/0.75\in[0,1]$.
\colorbox{green!15}{Green cells}: that estimator's $\RMC^{\pm}>0$.
$N^*=$ largest $N$ with positive global-pooled RMC(CP) under the Case A ceiling; this construction-level boundary is the headline boundary in the main paper, and the calibrated Case B certificate can tighten it (\Cref{tab:com_pern_caseB}).}
\label{tab:com_pern}
\renewcommand{\arraystretch}{1.12}
\resizebox{\textwidth}{!}{%
\begin{tabular}{p{2.2cm}r rrrr rr rr rr rr rr}
\toprule
\multirow{2}{*}{Model} & \multirow{2}{*}{$N$}
  & \multirow{2}{*}{Acc}
  & \multicolumn{2}{c}{F1 (per class)}
  & \multirow{2}{*}{MCC}
  & \multirow{2}{*}{G-mean}
  & \multirow{2}{*}{BalAcc}
  & \multirow{2}{*}{ADR} & \multirow{2}{*}{ADR$^{+w}$}
  & \multicolumn{2}{c}{RMC(no LCB)}
  & \multicolumn{2}{c}{RMC(CP)}
  & \multicolumn{2}{c}{RMC(Wilson)} \\
\cmidrule(lr){4-5}\cmidrule(lr){11-12}\cmidrule(lr){13-14}\cmidrule(lr){15-16}
 & & & F1$_{\mathrm{SAT}}$ & F1$_{\mathrm{UNSAT}}$ & & & & &
  & $\RMC$ & $\mathrm{nRMC}$
  & $\RMC^{\pm}$ & $\mathrm{nRMC}$
  & $\RMC^{\pm}$ & $\mathrm{nRMC}$ \\
\midrule
\multirow{7}{*}{\makecell[l]{\textbf{GPT-5}\\($N^*=30$)}}
  & 5  & 1.000 & 1.000 & 1.000 & 1.000 & 1.000 & 1.000 & 1.000 & 1.000 & $+$0.750 & 1.000 & \cellcolor{green!15}$+$0.679 & \cellcolor{green!15}0.905 & \cellcolor{green!15}$+$0.677 & \cellcolor{green!15}0.903 \\
  & 10 & 0.990 & 0.991 & 0.989 & 0.980 & 0.991 & 0.991 & 0.989 & 0.989 & $+$0.732 & 0.976 & \cellcolor{green!15}$+$0.647 & \cellcolor{green!15}0.863 & \cellcolor{green!15}$+$0.647 & \cellcolor{green!15}0.862 \\
  & 25 & 0.864 & 0.847 & 0.877 & 0.745 & 0.857 & 0.864 & 0.736 & 0.673 & $+$0.484 & 0.645 & \cellcolor{green!15}$+$0.362 & \cellcolor{green!15}0.482 & \cellcolor{green!15}$+$0.365 & \cellcolor{green!15}0.487 \\
  & 30 & 0.735 & 0.651 & 0.787 & 0.533 & 0.694 & 0.734 & 0.486 & 0.440 & $+$0.232 & 0.309 & \cellcolor{green!15}$+$0.117 & \cellcolor{green!15}0.156 & \cellcolor{green!15}$+$0.122 & \cellcolor{green!15}0.163 \\
  & 40 & 0.645 & 0.466 & 0.735 & 0.393 & 0.551 & 0.645 & 0.309 & 0.164 & $+$0.053 & 0.071 & $-$0.040 & 0.000 & $-$0.034 & 0.000 \\
  & 50 & 0.573 & 0.288 & 0.695 & 0.242 & 0.410 & 0.573 & 0.173 & 0.064 & $-$0.082 & 0.000 & $-$0.151 & 0.000 & $-$0.145 & 0.000 \\
  & 60 & 0.555 & 0.222 & 0.688 & 0.210 & 0.353 & 0.555 & 0.127 & 0.000 & $-$0.125 & 0.000 & $-$0.183 & 0.000 & $-$0.178 & 0.000 \\
\midrule
\multirow{5}{*}{\makecell[l]{\textbf{Opus 4.7}\\\textbf{(med.\ think)}\\($N^*=25$)}}
  & 5  & 0.945 & 0.947 & 0.942 & 0.892 & 0.944 & 0.945 & 0.897 & 0.813 & $+$0.641 & 0.855 & \cellcolor{green!15}$+$0.532 & \cellcolor{green!15}0.710 & \cellcolor{green!15}$+$0.534 & \cellcolor{green!15}0.712 \\
  & 10 & 0.959 & 0.960 & 0.958 & 0.918 & 0.958 & 0.959 & 0.917 & 0.861 & $+$0.669 & 0.892 & \cellcolor{green!15}$+$0.566 & \cellcolor{green!15}0.754 & \cellcolor{green!15}$+$0.567 & \cellcolor{green!15}0.756 \\
  & 25 & 0.935 & 0.956 & 0.875 & 0.838 & 0.889 & 0.894 & 0.759 & 0.655 & $+$0.541 & 0.721 & \cellcolor{green!15}$+$0.342 & \cellcolor{green!15}0.456 & \cellcolor{green!15}$+$0.352 & \cellcolor{green!15}0.469 \\
  & 50 & 0.527 & 0.620 & 0.372 & 0.134 & 0.467 & 0.555 & 0.200 & 0.060 & $-$0.032 & 0.000 & $-$0.126 & 0.000 & $-$0.119 & 0.000 \\
  & 60 & 0.530 & 0.656 & 0.256 & 0.157 & 0.382 & 0.549 & 0.156 & 0.000 & $-$0.104 & 0.000 & $-$0.172 & 0.000 & $-$0.166 & 0.000 \\
\midrule
\multirow{6}{*}{\makecell[l]{\textbf{DeepSeek-}\\\textbf{Reasoner}\\($N^*=10$)}}
  & 5  & 0.986 & 0.986 & 0.987 & 0.973 & 0.986 & 0.986 & 0.973 & 0.973 & $+$0.723 & 0.964 & \cellcolor{green!15}$+$0.642 & \cellcolor{green!15}0.856 & \cellcolor{green!15}$+$0.642 & \cellcolor{green!15}0.856 \\
  & 10 & 0.876 & 0.865 & 0.886 & 0.774 & 0.872 & 0.879 & 0.780 & 0.780 & $+$0.511 & 0.681 & \cellcolor{green!15}$+$0.391 & \cellcolor{green!15}0.522 & \cellcolor{green!15}$+$0.394 & \cellcolor{green!15}0.526 \\
  & 25 & 0.626 & 0.454 & 0.716 & 0.349 & 0.541 & 0.633 & 0.289 & 0.144 & $+$0.043 & 0.057 & $-$0.053 & 0.000 & $-$0.047 & 0.000 \\
  & 30 & 0.515 & 0.154 & 0.660 & 0.075 & 0.288 & 0.519 & 0.073 & 0.010 & $-$0.167 & 0.000 & $-$0.214 & 0.000 & $-$0.209 & 0.000 \\
  & 50 & 0.486 & 0.194 & 0.622 & $-$0.011 & 0.324 & 0.496 & 0.090 & 0.000 & $-$0.145 & 0.000 & $-$0.198 & 0.000 & $-$0.193 & 0.000 \\
  & 60 & 0.509 & 0.160 & 0.653 & 0.049 & 0.294 & 0.513 & 0.086 & 0.010 & $-$0.164 & 0.000 & $-$0.211 & 0.000 & $-$0.206 & 0.000 \\
\bottomrule
\multicolumn{16}{l}{\colorbox{green!15}{\phantom{X}} = that estimator's $\RMC^{\pm}>0$.}
\end{tabular}%
}
\end{table*}



\paragraph{RQ2 result on commercial models.}
GPT-5 obtains the strongest commercial certificate.  It is near perfect at $N=5$ and $N=10$, remains certified at $N=25$, and then loses certification as the positive-arm sensitivity weakens.  The important point is the form of the evidence: RMC(CP) discounts raw within-group separability by a finite-sample lower confidence bound and then compares the result with a shortcut ceiling.  RMC itself uses category outcomes; satisfying assignments are reported separately through ADR$^{+w}$.  Consequently, GPT-5's positive certificate at smaller $N$ is stronger evidence of reasoning than its accuracy alone because it survives group-level factorization, shortcut discounting, and the 95\% confidence discipline.

Claude Opus 4.7 medium-thinking and DeepSeek-Reasoner also certify at small sizes but show steeper capability cliffs.  DeepSeek-Reasoner illustrates why RMC should be preferred over raw separability: at $N=25$, the raw RMC estimate is still positive ($+0.043$), while RMC(CP) falls below zero ($-0.053$).  The point estimate suggests residual capability; the confidence-bounded certificate rejects it.  This difference is exactly the role of the 95\% confidence discipline.

\paragraph{Case B sensitivity: CNF side-channel adversary battery.}
The Case A results assume $\shortcutDisc=0$, meaning that the minimal-edit construction has removed the observable side channels included in the design.  Case B relaxes this assumption by estimating $\hatshortcutDisc$ from the calibrated \emph{final residual} CNF side-channel battery (mean/min/max clause length, literal-polarity ratio, variable coverage) and replacing the fixed ceiling with the residual shortcut ceiling $\hatshortcutCeiling=\bigl((1+\hatshortcutDisc)/2\bigr)^2-0.25$.  Thus Case B measures the extra shortcut-only separability above the Case-A chance floor, while remaining explicitly tied to the stated adversary class.

\Cref{tab:com_pern_caseB} shows that under this updated residual Case B ceiling, GPT-5 and Claude Opus 4.7 medium-thinking remain globally certified across all reported $N$ values, while DeepSeek-Reasoner remains certified through $N=25$ and fails from $N=30$ onward.  The global Case B margins are therefore positive but substantially smaller than the corresponding raw separability gaps, which is exactly the intended role of RMC: preserve the semantic discrimination target while discounting what can be attributed to the measured residual shortcut class.  We therefore report Case B as the primary shortcut-adjusted certificate in the main paper, with Case A retained as the idealized construction-level reference that defines $N^*$; the ceiling-estimation protocol is given in \Cref{sec:ceiling_estimation}.

\begin{table*}[ht]
\centering
\caption{Per-$N$ RMC metrics under \textbf{Case B} (estimated $\shortcutCeiling$) for the three certified commercial reasoning models on 3-SAT.
All column definitions match \Cref{tab:com_pern}.
$\shortcutCeiling$: per-$N$ mean of the permutation-calibrated estimated shortcut ceiling from the final residual CNF side-channel adversary battery (mean/min/max clause length, literal-polarity ratio, variable coverage; calibrated point estimate; the strict UCB-inflated variant is discussed in Remark~\ref{rem:caseb-variants}).
Here $\shortcutCeiling=\bigl((1+\hatshortcutDisc)/2\bigr)^2-0.25=\hatshortcutDisc/2+{\hatshortcutDisc}^{2}/4$, so it reports only the residual shortcut-only contribution above the Case-A chance floor; e.g.\ $\shortcutCeiling=0.159$ corresponds to $\hatshortcutDisc\approx0.28$.
ADR$^{+w}$ values copied from Case A (CDCL solver unavailable in Case B runtime).
$\hat{s}$ is the point-estimate separability used by RMC(no LCB), so RMC(no LCB)$=\hat{s}-\shortcutCeiling$.
$\mathrm{nRMC}=\max(0,\RMC^{\pm})/(1-\shortcutCeiling)$ (normalized relative to available headroom; $0.000$ when $\RMC^{\pm}\le0$).
Thus nRMC measures the fraction of remaining headroom above the estimated shortcut ceiling rather than the raw margin itself: for GPT-5, $N=5$ has $\hat{s}=1.000$, so $0.841/(1-0.159)=1.000$, whereas $N=10$ has a slightly larger raw RMC but $\hat{s}=0.982<1$, so $0.862/(1-0.119)=0.979$.
Positive RMC(CP) or RMC(Wilson) entries indicate that the model clears the estimated Case-B side-channel ceiling under that estimator.}
\label{tab:com_pern_caseB}
\renewcommand{\arraystretch}{1.12}
\resizebox{\textwidth}{!}{%
\begin{tabular}{p{2.2cm}r rrrr rr rr rr rr rr rr}
\toprule
\multirow{2}{*}{Model} & \multirow{2}{*}{$N$}
  & \multirow{2}{*}{Acc}
  & \multicolumn{2}{c}{F1 (per class)}
  & \multirow{2}{*}{MCC}
  & \multirow{2}{*}{G-mean}
  & \multirow{2}{*}{BalAcc}
  & \multirow{2}{*}{ADR} & \multirow{2}{*}{ADR$^{+w}$}
  & \multirow{2}{*}{$\shortcutCeiling$}
  & \multirow{2}{*}{$\hat{s}$}
  & \multicolumn{2}{c}{RMC(no LCB)}
  & \multicolumn{2}{c}{RMC(CP)}
  & \multicolumn{2}{c}{RMC(Wilson)} \\
\cmidrule(lr){4-5}\cmidrule(lr){13-14}\cmidrule(lr){15-16}\cmidrule(lr){17-18}
 & & & F1$_{\mathrm{SAT}}$ & F1$_{\mathrm{UNSAT}}$ & & & & & & &
  & $\RMC$ & $\mathrm{nRMC}$
  & $\RMC^{\pm}$ & $\mathrm{nRMC}$
  & $\RMC^{\pm}$ & $\mathrm{nRMC}$ \\
\midrule
\multirow{7}{*}{\makecell[l]{\textbf{GPT-5}\\(Case B)}}
  & 5  & 1.000 & 1.000 & 1.000 & 1.000 & 1.000 & 1.000 & 1.000 & 1.000 & 0.159 & 1.000 & $+$0.841 & 1.000 & $+$0.770 & 0.915 & $+$0.768 & 0.914 \\
  & 10 & 0.990 & 0.991 & 0.989 & 0.980 & 0.991 & 0.991 & 0.989 & 0.989 & 0.119 & 0.982 & $+$0.862 & 0.979 & $+$0.778 & 0.883 & $+$0.777 & 0.883 \\
  & 25 & 0.864 & 0.847 & 0.877 & 0.745 & 0.857 & 0.864 & 0.736 & 0.673 & 0.145 & 0.734 & $+$0.589 & 0.689 & $+$0.467 & 0.546 & $+$0.470 & 0.550 \\
  & 30 & 0.735 & 0.651 & 0.787 & 0.533 & 0.694 & 0.734 & 0.486 & 0.440 & 0.140 & 0.482 & $+$0.342 & 0.398 & $+$0.227 & 0.264 & $+$0.232 & 0.270 \\
  & 40 & 0.645 & 0.466 & 0.735 & 0.393 & 0.551 & 0.645 & 0.309 & 0.164 & 0.156 & 0.303 & $+$0.148 & 0.175 & $+$0.054 & 0.064 & $+$0.060 & 0.071 \\
  & 50 & 0.573 & 0.288 & 0.695 & 0.242 & 0.410 & 0.573 & 0.173 & 0.064 & 0.089 & 0.168 & $+$0.079 & 0.087 & $+$0.010 & 0.011 & $+$0.016 & 0.018 \\
  & 60 & 0.555 & 0.222 & 0.688 & 0.210 & 0.353 & 0.555 & 0.127 & 0.000 & 0.046 & 0.125 & $+$0.079 & 0.083 & $+$0.021 & 0.022 & $+$0.026 & 0.028 \\
\midrule
\multirow{5}{*}{\makecell[l]{\textbf{Opus 4.7}\\\textbf{(med.\ think)}\\(Case B)}}
  & 5  & 0.945 & 0.947 & 0.942 & 0.892 & 0.944 & 0.945 & 0.897 & 0.813 & 0.139 & 0.891 & $+$0.751 & 0.873 & $+$0.643 & 0.747 & $+$0.644 & 0.749 \\
  & 10 & 0.959 & 0.960 & 0.958 & 0.918 & 0.958 & 0.959 & 0.917 & 0.861 & 0.117 & 0.919 & $+$0.802 & 0.908 & $+$0.699 & 0.791 & $+$0.700 & 0.793 \\
  & 25 & 0.935 & 0.956 & 0.875 & 0.838 & 0.889 & 0.894 & 0.759 & 0.655 & 0.086 & 0.791 & $+$0.705 & 0.771 & $+$0.506 & 0.553 & $+$0.516 & 0.564 \\
  & 50 & 0.527 & 0.620 & 0.372 & 0.134 & 0.467 & 0.555 & 0.200 & 0.060 & 0.065 & 0.218 & $+$0.154 & 0.164 & $+$0.059 & 0.063 & $+$0.066 & 0.070 \\
  & 60 & 0.530 & 0.656 & 0.256 & 0.157 & 0.382 & 0.549 & 0.156 & 0.000 & 0.031 & 0.146 & $+$0.115 & 0.119 & $+$0.047 & 0.049 & $+$0.053 & 0.055 \\
\midrule
\multirow{7}{*}{\makecell[l]{\textbf{DeepSeek-}\\\textbf{Reasoner}\\(Case B)}}
  & 5  & 0.986 & 0.986 & 0.987 & 0.973 & 0.986 & 0.986 & 0.973 & 0.973 & 0.134 & 0.973 & $+$0.839 & 0.969 & $+$0.758 & 0.875 & $+$0.758 & 0.875 \\
  & 10 & 0.876 & 0.865 & 0.886 & 0.774 & 0.872 & 0.879 & 0.780 & 0.780 & 0.133 & 0.761 & $+$0.628 & 0.724 & $+$0.509 & 0.587 & $+$0.512 & 0.590 \\
  & 25 & 0.626 & 0.454 & 0.716 & 0.349 & 0.541 & 0.633 & 0.289 & 0.144 & 0.137 & 0.293 & $+$0.156 & 0.181 & $+$0.061 & 0.070 & $+$0.066 & 0.077 \\
  & 30 & 0.515 & 0.154 & 0.660 & 0.075 & 0.288 & 0.519 & 0.073 & 0.010 & 0.125 & 0.083 & $-$0.042 & 0.000 & $-$0.089 & 0.000 & $-$0.084 & 0.000 \\
  & 40 & 0.524 & 0.183 & 0.664 & 0.072 & 0.317 & 0.520 & 0.093 & 0.031 & 0.143 & 0.101 & $-$0.043 & 0.000 & $-$0.096 & 0.000 & $-$0.090 & 0.000 \\
  & 50 & 0.486 & 0.194 & 0.622 & $-$0.011 & 0.324 & 0.496 & 0.090 & 0.000 & 0.073 & 0.105 & $+$0.032 & 0.034 & $-$0.021 & 0.000 & $-$0.016 & 0.000 \\
  & 60 & 0.509 & 0.160 & 0.653 & 0.049 & 0.294 & 0.513 & 0.086 & 0.010 & 0.047 & 0.086 & $+$0.039 & 0.041 & $-$0.008 & 0.000 & $-$0.003 & 0.000 \\
\bottomrule
\multicolumn{18}{l}{Positive $\RMC^{\pm}$ entries indicate model discrimination above the estimated Case-B side-channel ceiling.}
\end{tabular}%
}
\end{table*}

\paragraph{Case C sensitivity: complexity-debiased side-channel ceiling.}
Case C keeps the same residual shortcut-ceiling formula as Case B,
$\hatshortcutCeiling=\bigl((1+\hatshortcutDisc)/2\bigr)^2-0.25$,
but changes how $\hatshortcutDisc$ is estimated.
Specifically, each feature in the final Case-B residual battery is first residualized by a pooled per-feature ridge regression on the raw linear complexity covariates $(\text{num\_vars}, \text{num\_clauses}, \text{num\_literals}, \text{file\_size\_bytes})$, and the resulting per-instance residuals replace the raw feature values before the max-KS adversary is run.
The residualized max-KS score is then permutation-calibrated within group exactly as in Case B.  This is a feature-space Frisch--Waugh--Lovell partial-out construction: it asks how much shortcut leakage remains after removing the component linearly predictable from generic structural scaling, rather than scaling the final scalar leakage by a post-hoc attenuation factor.  The implemented design deliberately stops at these raw linear controls, without higher-order, interaction, or log-expanded covariates.

The reported Case C $\hat\shortcutDisc$ is then the 50/50 construction-level blend of the permutation-calibrated residual estimate and the Case A reference $\delta=0$:
\[
\hat\shortcutDisc^{(\text{C})}
=
\tfrac{1}{2}\max\!\bigl(0,\hat\shortcutDisc_{\text{obs,resid}}-\hat b_{\text{perm,resid}}\bigr).
\]
Equivalently, if $\hat\shortcutDisc_{\text{FWL,cal}}:=\max(0,\hat\shortcutDisc_{\text{obs,resid}}-\hat b_{\text{perm,resid}})$ denotes the calibrated residualized estimate, then $\hat\shortcutDisc^{(\text{C})}=\tfrac{1}{2}\hat\shortcutDisc_{\text{FWL,cal}}$.  This positions Case C halfway between the calibrated residualized estimate and the idealized construction-level $\delta=0$, and avoids letting a single residualized leakage estimate dominate the certificate.

\Cref{tab:com_pern_caseC} shows that with this blend, the Case C ceiling is materially smaller than the raw Case B ceiling and sits a notch above the Case A construction-level floor: across the three certified commercial 3-SAT models, the per-$N$ mean residual $\shortcutCeiling$ lies in the $0.035$--$0.091$ range (compared with $0.10$--$0.30$ for raw Case B).
Under the global-pooled CP estimator with this ceiling, GPT-5 remains Case C certified at every tested $N\in\{5,10,25,30,40,50,60\}$ (CP margin decays from $+0.838$ at $N=5$ to $+0.006$ at $N=60$); Claude Opus 4.7 (medium thinking) is also certified at every tested $N\in\{5,10,25,50,60\}$ (from $+0.705$ at $N=5$ to $+0.028$ at $N=60$); DeepSeek-Reasoner stays certified through $N=25$ ($+0.110$), fails at $N=30$ ($-0.041$), recovers a borderline $+0.001$ at $N=50$, and finishes negative at $N=60$.
The Wilson factorized variant gives essentially the same boundaries.
The effect of Case C is therefore not to change the underlying discriminative trend, but to isolate the portion of the measured side-channel leakage that is not already explained by generic structural growth.

\begin{table*}[ht]
\centering
\caption{Per-$N$ RMC metrics under \textbf{Case C} (complexity-debiased estimated $\shortcutCeiling$, feature-level residualization) for the three certified commercial reasoning models on 3-SAT.
All column definitions match \Cref{tab:com_pern}.
$\shortcutCeiling$: per-$N$ mean of the complexity-debiased residual shortcut ceiling.
Case C keeps the same residual ceiling transform as Case B, $\hatshortcutCeiling=\bigl((1+\hatshortcutDisc)/2\bigr)^2-0.25$, but estimates $\hatshortcutDisc$ by residualizing each feature in the final Case-B residual battery against the raw linear covariates (num\_vars, num\_clauses, num\_literals, file\_size\_bytes), applying the max-KS adversary and permutation calibration to those residuals, and then taking the default 50/50 blend with the Case A reference $\delta=0$.
$\hat{s}$ is the point-estimate separability used by RMC(no LCB), so RMC(no LCB)$=\hat{s}-\shortcutCeiling$.
$\mathrm{nRMC}=\max(0,\RMC^{\pm})/(1-\shortcutCeiling)$.
Positive RMC(CP) or RMC(Wilson) entries indicate that the model clears the estimated Case-C side-channel ceiling under that estimator.}
\label{tab:com_pern_caseC}
\renewcommand{\arraystretch}{1.12}
\resizebox{\textwidth}{!}{%
\begin{tabular}{p{2.2cm}r rrrr rr rr rr rr rr rr}
\toprule
\multirow{2}{*}{Model} & \multirow{2}{*}{$N$}
  & \multirow{2}{*}{Acc}
  & \multicolumn{2}{c}{F1 (per class)}
  & \multirow{2}{*}{MCC}
  & \multirow{2}{*}{G-mean}
  & \multirow{2}{*}{BalAcc}
  & \multirow{2}{*}{ADR} & \multirow{2}{*}{ADR$^{+w}$}
  & \multirow{2}{*}{$\shortcutCeiling$}
  & \multirow{2}{*}{$\hat{s}$}
  & \multicolumn{2}{c}{RMC(no LCB)}
  & \multicolumn{2}{c}{RMC(CP)}
  & \multicolumn{2}{c}{RMC(Wilson)} \\
\cmidrule(lr){4-5}\cmidrule(lr){13-14}\cmidrule(lr){15-16}\cmidrule(lr){17-18}
 & & & F1$_{\mathrm{SAT}}$ & F1$_{\mathrm{UNSAT}}$ & & & & & & &
  & $\RMC$ & $\mathrm{nRMC}$
  & $\RMC^{\pm}$ & $\mathrm{nRMC}$
  & $\RMC^{\pm}$ & $\mathrm{nRMC}$ \\
\midrule
\multirow{7}{*}{\makecell[l]{\textbf{GPT-5}\\(Case C)}}
  & 5  & 1.000 & 1.000 & 1.000 & 1.000 & 1.000 & 1.000 & 1.000 & 1.000 & 0.091 & 1.000 & \cellcolor{green!15}$+$0.909 & \cellcolor{green!15}1.000 & \cellcolor{green!15}$+$0.838 & \cellcolor{green!15}0.922 & \cellcolor{green!15}$+$0.836 & \cellcolor{green!15}0.920 \\
  & 10 & 0.990 & 0.991 & 0.989 & 0.980 & 0.991 & 0.991 & 0.989 & 0.989 & 0.065 & 0.982 & \cellcolor{green!15}$+$0.916 & \cellcolor{green!15}0.980 & \cellcolor{green!15}$+$0.832 & \cellcolor{green!15}0.890 & \cellcolor{green!15}$+$0.831 & \cellcolor{green!15}0.890 \\
  & 25 & 0.864 & 0.847 & 0.877 & 0.745 & 0.857 & 0.864 & 0.736 & 0.673 & 0.085 & 0.734 & \cellcolor{green!15}$+$0.649 & \cellcolor{green!15}0.709 & \cellcolor{green!15}$+$0.527 & \cellcolor{green!15}0.576 & \cellcolor{green!15}$+$0.530 & \cellcolor{green!15}0.579 \\
  & 30 & 0.735 & 0.651 & 0.787 & 0.533 & 0.694 & 0.734 & 0.486 & 0.440 & 0.085 & 0.482 & \cellcolor{green!15}$+$0.397 & \cellcolor{green!15}0.434 & \cellcolor{green!15}$+$0.282 & \cellcolor{green!15}0.308 & \cellcolor{green!15}$+$0.287 & \cellcolor{green!15}0.314 \\
  & 40 & 0.645 & 0.466 & 0.735 & 0.393 & 0.551 & 0.645 & 0.309 & 0.164 & 0.073 & 0.303 & \cellcolor{green!15}$+$0.231 & \cellcolor{green!15}0.249 & \cellcolor{green!15}$+$0.137 & \cellcolor{green!15}0.148 & \cellcolor{green!15}$+$0.143 & \cellcolor{green!15}0.154 \\
  & 50 & 0.573 & 0.288 & 0.695 & 0.242 & 0.410 & 0.573 & 0.173 & 0.064 & 0.049 & 0.168 & \cellcolor{green!15}$+$0.119 & \cellcolor{green!15}0.125 & \cellcolor{green!15}$+$0.050 & \cellcolor{green!15}0.053 & \cellcolor{green!15}$+$0.056 & \cellcolor{green!15}0.059 \\
  & 60 & 0.555 & 0.222 & 0.688 & 0.210 & 0.353 & 0.555 & 0.127 & 0.000 & 0.061 & 0.125 & \cellcolor{green!15}$+$0.064 & \cellcolor{green!15}0.068 & \cellcolor{green!15}$+$0.006 & \cellcolor{green!15}0.006 & \cellcolor{green!15}$+$0.012 & \cellcolor{green!15}0.012 \\
\midrule
\multirow{5}{*}{\makecell[l]{\textbf{Opus 4.7}\\\textbf{(med.\ think)}\\(Case C)}}
  & 5  & 0.945 & 0.947 & 0.942 & 0.892 & 0.944 & 0.945 & 0.897 & 0.813 & 0.077 & 0.891 & \cellcolor{green!15}$+$0.813 & \cellcolor{green!15}0.882 & \cellcolor{green!15}$+$0.705 & \cellcolor{green!15}0.764 & \cellcolor{green!15}$+$0.706 & \cellcolor{green!15}0.766 \\
  & 10 & 0.959 & 0.960 & 0.958 & 0.918 & 0.958 & 0.959 & 0.917 & 0.861 & 0.068 & 0.919 & \cellcolor{green!15}$+$0.851 & \cellcolor{green!15}0.913 & \cellcolor{green!15}$+$0.748 & \cellcolor{green!15}0.802 & \cellcolor{green!15}$+$0.749 & \cellcolor{green!15}0.804 \\
  & 25 & 0.935 & 0.956 & 0.875 & 0.838 & 0.889 & 0.894 & 0.759 & 0.655 & 0.035 & 0.791 & \cellcolor{green!15}$+$0.756 & \cellcolor{green!15}0.783 & \cellcolor{green!15}$+$0.557 & \cellcolor{green!15}0.577 & \cellcolor{green!15}$+$0.567 & \cellcolor{green!15}0.587 \\
  & 50 & 0.527 & 0.620 & 0.372 & 0.134 & 0.467 & 0.555 & 0.200 & 0.060 & 0.040 & 0.218 & \cellcolor{green!15}$+$0.178 & \cellcolor{green!15}0.185 & \cellcolor{green!15}$+$0.083 & \cellcolor{green!15}0.087 & \cellcolor{green!15}$+$0.090 & \cellcolor{green!15}0.094 \\
  & 60 & 0.530 & 0.656 & 0.256 & 0.157 & 0.382 & 0.549 & 0.156 & 0.000 & 0.050 & 0.146 & \cellcolor{green!15}$+$0.096 & \cellcolor{green!15}0.101 & \cellcolor{green!15}$+$0.028 & \cellcolor{green!15}0.029 & \cellcolor{green!15}$+$0.034 & \cellcolor{green!15}0.036 \\
\midrule
\multirow{7}{*}{\makecell[l]{\textbf{DeepSeek-}\\\textbf{Reasoner}\\(Case C)}}
  & 5  & 0.986 & 0.986 & 0.987 & 0.973 & 0.986 & 0.986 & 0.973 & 0.973 & 0.079 & 0.973 & \cellcolor{green!15}$+$0.893 & \cellcolor{green!15}0.970 & \cellcolor{green!15}$+$0.813 & \cellcolor{green!15}0.883 & \cellcolor{green!15}$+$0.812 & \cellcolor{green!15}0.882 \\
  & 10 & 0.876 & 0.865 & 0.886 & 0.774 & 0.872 & 0.879 & 0.780 & 0.780 & 0.067 & 0.761 & \cellcolor{green!15}$+$0.694 & \cellcolor{green!15}0.744 & \cellcolor{green!15}$+$0.574 & \cellcolor{green!15}0.616 & \cellcolor{green!15}$+$0.577 & \cellcolor{green!15}0.619 \\
  & 25 & 0.626 & 0.454 & 0.716 & 0.349 & 0.541 & 0.633 & 0.289 & 0.144 & 0.088 & 0.293 & \cellcolor{green!15}$+$0.205 & \cellcolor{green!15}0.225 & \cellcolor{green!15}$+$0.110 & \cellcolor{green!15}0.120 & \cellcolor{green!15}$+$0.115 & \cellcolor{green!15}0.126 \\
  & 30 & 0.515 & 0.154 & 0.660 & 0.075 & 0.288 & 0.519 & 0.073 & 0.010 & 0.077 & 0.083 & \cellcolor{green!15}$+$0.006 & \cellcolor{green!15}0.006 & $-$0.041 & 0.000 & $-$0.036 & 0.000 \\
  & 40 & 0.524 & 0.183 & 0.664 & 0.072 & 0.317 & 0.520 & 0.093 & 0.031 & 0.069 & 0.101 & \cellcolor{green!15}$+$0.032 & \cellcolor{green!15}0.034 & $-$0.021 & 0.000 & $-$0.015 & 0.000 \\
  & 50 & 0.486 & 0.194 & 0.622 & $-$0.011 & 0.324 & 0.496 & 0.090 & 0.000 & 0.051 & 0.105 & \cellcolor{green!15}$+$0.054 & \cellcolor{green!15}0.057 & \cellcolor{green!15}$+$0.001 & \cellcolor{green!15}0.001 & \cellcolor{green!15}$+$0.006 & \cellcolor{green!15}0.006 \\
  & 60 & 0.509 & 0.160 & 0.653 & 0.049 & 0.294 & 0.513 & 0.086 & 0.010 & 0.057 & 0.086 & \cellcolor{green!15}$+$0.029 & \cellcolor{green!15}0.031 & $-$0.018 & 0.000 & $-$0.013 & 0.000 \\
\bottomrule
\multicolumn{18}{l}{Positive $\RMC^{\pm}$ entries indicate model discrimination above the estimated Case-C side-channel ceiling.}
\end{tabular}%
}
\end{table*}

\paragraph{Case D diagnostic: separability-only lower bound (3-SAT).}
Case D removes both ingredients that make Cases A--C conservative relative to raw balanced discrimination: it does not subtract any shortcut ceiling, and it does not apply the SAT-side witness min-gate.  The reported quantity is therefore just the factorized separability lower bound itself,
\[
\RMC_D(G)=\LCB_{\alpha/2}(\hat\theta_P(G))\cdot\LCB_{\alpha/2}(\hat\theta_N(G)),
\]
with $\alpha=0.05$ and Bonferroni split $\alpha/2$ across the two arms.
Because $\shortcutCeiling=0$, the point-estimate Case-D quantity is simply $\hat s$, and the normalized score would be numerically identical to the reported value; we therefore omit redundant $\mathrm{nRMC}$ columns.

\Cref{tab:com_pern_caseD} should be read as a separability-only upper envelope.  It answers a narrower question than Cases B and C: how much balanced SAT/UNSAT discrimination can be certified from category outcomes alone, before any shortcut charge is imposed?  Under this view, all three certified commercial 3-SAT models remain positive throughout their tested ranges.  GPT-5 declines from a pooled CP lower bound of $0.929$ at $N=5$ to $0.067$ at $N=60$; Opus 4.7 (medium thinking) declines from $0.782$ at $N=5$ to $0.078$ at $N=60$; and DeepSeek-Reasoner declines from $0.892$ at $N=5$ to $0.039$ at $N=60$.  The contrast with Cases B and C is the point: when Case D stays positive but the shortcut-adjusted certificate fails, the model still shows confidence-bounded balanced discrimination, but not enough to exceed the shortcut account under test.

\begin{table*}[ht]
\centering
\caption{Per-$N$ \textbf{Case D} separability-only diagnostics for the three certified commercial reasoning models on 3-SAT.
Case D removes both shortcut subtraction and witness gating.  The reported quantity is therefore the factorized separability statistic itself: the point-estimate column is $\hat{s}=\hat\theta_P\hat\theta_N$, RMC(CP) is the pooled Clopper--Pearson factorized lower bound, and RMC(Wilson) is the pooled Wilson factorized lower bound.
Because $\shortcutCeiling=0$ in Case D, RMC(no LCB)$=\hat{s}$ and the normalized score is numerically identical to the reported value, so redundant $\mathrm{nRMC}$ columns are omitted.
\colorbox{green!15}{Green cells}: positive separability lower bound under that estimator.}
\label{tab:com_pern_caseD}
\renewcommand{\arraystretch}{1.12}
\resizebox{\textwidth}{!}{%
\begin{tabular}{p{2.2cm}r rrrr rr r rr}
\toprule
Model & $N$ & Acc & F1$_{\mathrm{SAT}}$ & F1$_{\mathrm{UNSAT}}$ & MCC & ADR & ADR$^{+w}$ & $\hat{s}$ & RMC(CP) & RMC(Wilson) \\
\midrule
\multirow{7}{*}{\makecell[l]{\textbf{GPT-5}\\(Case D)}}
  & 5  & 1.000 & 1.000 & 1.000 & 1.000 & 1.000 & 1.000 & 1.000 & \cellcolor{green!15}$+$0.929 & \cellcolor{green!15}$+$0.927 \\
  & 10 & 0.990 & 0.991 & 0.989 & 0.980 & 0.989 & 0.989 & 0.982 & \cellcolor{green!15}$+$0.897 & \cellcolor{green!15}$+$0.897 \\
  & 25 & 0.864 & 0.847 & 0.877 & 0.745 & 0.736 & 0.673 & 0.734 & \cellcolor{green!15}$+$0.612 & \cellcolor{green!15}$+$0.615 \\
  & 30 & 0.735 & 0.651 & 0.787 & 0.533 & 0.486 & 0.440 & 0.482 & \cellcolor{green!15}$+$0.367 & \cellcolor{green!15}$+$0.372 \\
  & 40 & 0.645 & 0.466 & 0.735 & 0.393 & 0.309 & 0.164 & 0.303 & \cellcolor{green!15}$+$0.210 & \cellcolor{green!15}$+$0.216 \\
  & 50 & 0.573 & 0.288 & 0.695 & 0.242 & 0.173 & 0.064 & 0.168 & \cellcolor{green!15}$+$0.099 & \cellcolor{green!15}$+$0.105 \\
  & 60 & 0.555 & 0.222 & 0.688 & 0.210 & 0.127 & 0.000 & 0.125 & \cellcolor{green!15}$+$0.067 & \cellcolor{green!15}$+$0.072 \\
\midrule
\multirow{5}{*}{\makecell[l]{\textbf{Opus 4.7}\\\textbf{(med.\ think)}\\(Case D)}}
  & 5  & 0.945 & 0.947 & 0.942 & 0.892 & 0.897 & 0.813 & 0.891 & \cellcolor{green!15}$+$0.782 & \cellcolor{green!15}$+$0.784 \\
  & 10 & 0.959 & 0.960 & 0.958 & 0.918 & 0.917 & 0.861 & 0.919 & \cellcolor{green!15}$+$0.816 & \cellcolor{green!15}$+$0.817 \\
  & 25 & 0.935 & 0.956 & 0.875 & 0.838 & 0.759 & 0.655 & 0.791 & \cellcolor{green!15}$+$0.592 & \cellcolor{green!15}$+$0.602 \\
  & 50 & 0.527 & 0.620 & 0.372 & 0.134 & 0.200 & 0.060 & 0.218 & \cellcolor{green!15}$+$0.124 & \cellcolor{green!15}$+$0.131 \\
  & 60 & 0.530 & 0.656 & 0.256 & 0.157 & 0.156 & 0.000 & 0.146 & \cellcolor{green!15}$+$0.078 & \cellcolor{green!15}$+$0.084 \\
\midrule
\multirow{7}{*}{\makecell[l]{\textbf{DeepSeek-}\\\textbf{Reasoner}\\(Case D)}}
  & 5  & 0.986 & 0.986 & 0.987 & 0.973 & 0.973 & 0.973 & 0.973 & \cellcolor{green!15}$+$0.892 & \cellcolor{green!15}$+$0.892 \\
  & 10 & 0.876 & 0.865 & 0.886 & 0.774 & 0.780 & 0.780 & 0.761 & \cellcolor{green!15}$+$0.641 & \cellcolor{green!15}$+$0.644 \\
  & 25 & 0.626 & 0.454 & 0.716 & 0.349 & 0.289 & 0.144 & 0.293 & \cellcolor{green!15}$+$0.197 & \cellcolor{green!15}$+$0.203 \\
  & 30 & 0.515 & 0.154 & 0.660 & 0.075 & 0.073 & 0.010 & 0.083 & \cellcolor{green!15}$+$0.036 & \cellcolor{green!15}$+$0.041 \\
  & 40 & 0.524 & 0.183 & 0.664 & 0.072 & 0.093 & 0.031 & 0.101 & \cellcolor{green!15}$+$0.048 & \cellcolor{green!15}$+$0.053 \\
  & 50 & 0.486 & 0.194 & 0.622 & $-$0.011 & 0.090 & 0.000 & 0.105 & \cellcolor{green!15}$+$0.052 & \cellcolor{green!15}$+$0.057 \\
  & 60 & 0.509 & 0.160 & 0.653 & 0.049 & 0.086 & 0.010 & 0.086 & \cellcolor{green!15}$+$0.039 & \cellcolor{green!15}$+$0.044 \\
\bottomrule
\multicolumn{11}{l}{\colorbox{green!15}{\phantom{X}} = positive separability lower bound under that estimator.}
\end{tabular}%
}
\end{table*}

\subsection{RQ2 continued: Open-source model certification on 2-SAT}
\label{sec:open}

We evaluate 13 open-source/open-weight configurations on 2-SAT instances across $N\in\{3,5,8,10,15,20,25,30,40,50\}$ variables; the task-assignment rationale and the cross-family comparability caveat are given in \Cref{sec:setup_exp}.
The same fundamental split as in the commercial domain applies: 10 non-thinking configurations exhibit label-bias failure at every $N$, while 3 thinking-mode configurations certify at small $N$.

\paragraph{Non-thinking models: label-bias failure at all $N$.}
Ten configurations---LLaMA-3.1-8B and 70B, Gemma-4-31B, Qwen2.5-32B, Qwen3-30B-A3B (with and without thinking), DeepSeek-R1-Distill-32B, and others---predict almost exclusively SAT or UNSAT for every formula regardless of $N$.
Instance accuracies cluster at $0.499$--$0.560$ across all $N$ values, and $\RMCSAT^{\pm}$ sits in $[-0.249, -0.150]$.
Enabling thinking mode on the Qwen3-30B-A3B-Instruct architecture yields no improvement: its thinking variant ($-0.207$) matches its non-thinking counterpart ($-0.217$), supporting the interpretation that the all-SAT bias comes from model behavior in this family rather than prompt format alone.
DeepSeek-R1-Distill-32B is a partial exception with $\theta_P=0.174$ and $\theta_N=0.884$---it has some discriminative tendency---but its $\LCB_{\mathrm{fact}}$ never exceeds the $0.25$ shortcut ceiling at any tested $N$.

\paragraph{Open-weight reasoning-mode models: per-$N$ certification.}
The current local 2-SAT open-weight reasoning-mode runs cover six model families: Qwen3-14B, Qwen3-32B, QwQ-32B, OpenAI gpt-oss-20b, NVIDIA Nemotron-H-47B-Reasoning-128K, and NVIDIA Llama-Nemotron-Super-49B v1.5.
Among them, the three Qwen-family models remain the only configurations with sustained certification beyond the smallest-$N$ regime.
Under the latest local rerun, gpt-oss-20b now certifies through $N=20$ under the Case-A pooled-CP criterion and fails from $N=25$ onward; Llama-Nemotron-Super-49B v1.5 in NVIDIA's reasoning-on mode certifies through $N=10$ and then fails from $N=15$ onward; Nemotron-H-47B-Reasoning-128K in reasoning-on mode does not clear the Case-A pooled-CP gate at any tested $N$.
\Cref{tab:os_pern} reports the corresponding per-$N$ metrics in the same format as \Cref{tab:com_pern}: per-class F1 and three global-pooled RMC estimators (RMC(no LCB), RMC(CP), RMC(Wilson)).

\paragraph{Compact notation for the open-weight 2-SAT per-$N$ tables.}
To avoid repeating long header explanations across \Cref{tab:os_pern,tab:os_pern_caseB,tab:os_pern_caseC,tab:os_pern_caseD}, we use a unified compact notation.
`F1$_S$` and `F1$_U$` denote the SAT and UNSAT classwise F1 scores.
`RMC$_0$` denotes RMC(no LCB), `RMC$_{CP}$` the global-pooled Clopper--Pearson margin, and `RMC$_W$` the global-pooled Wilson margin; `n`-prefixed columns are the corresponding normalized scores.
In Cases B and C, `$\shortcutCeiling$` is the per-$N$ mean estimated ceiling and `$\hat{s}$` is the corresponding point-estimate separability.
Green cells indicate positive certified margins under that estimator.

\begin{table*}[t]
\centering
\caption{Open-weight 2-SAT per-$N$ metrics under Case A.
$\mathrm{nRMC}=\max(0,\RMC^{\pm})/0.75$.
$N^*=$ largest $N$ with positive global-pooled RMC(CP) under the Case A ceiling.
Only NVIDIA \emph{reasoning-on} runs are reported; non-reasoning-on NVIDIA outputs are omitted from this table.}
\label{tab:os_pern}
\renewcommand{\arraystretch}{1.12}
\resizebox{\textwidth}{!}{%
\begin{tabular}{p{2.2cm}r rrrr rr rr rr rr rr}
\toprule
\multirow{2}{*}{Model} & \multirow{2}{*}{$N$}
  & \multirow{2}{*}{Acc}
  & \multicolumn{2}{c}{F1 (per class)}
  & \multirow{2}{*}{MCC}
  & \multirow{2}{*}{G-mean}
  & \multirow{2}{*}{BalAcc}
  & \multirow{2}{*}{ADR} & \multirow{2}{*}{ADR$^{+w}$}
  & \multicolumn{2}{c}{RMC(no LCB)}
  & \multicolumn{2}{c}{RMC(CP)}
  & \multicolumn{2}{c}{RMC(Wilson)} \\
\cmidrule(lr){4-5}\cmidrule(lr){11-12}\cmidrule(lr){13-14}\cmidrule(lr){15-16}
 & & & F1$_{\mathrm{SAT}}$ & F1$_{\mathrm{UNSAT}}$ & MCC & & & & &
  $\RMC$ & $\mathrm{nRMC}$
  & $\RMC^{\pm}$ & $\mathrm{nRMC}$
  & $\RMC^{\pm}$ & $\mathrm{nRMC}$ \\
\midrule
\multirow{11}{*}{\makecell[l]{\textbf{Qwen3-14B}\\\textbf{(think)}\\($N^*=20$)}}
  & 3  & 0.990 & 0.990 & 0.990 & 0.981 & 0.990 & 0.990 & 0.981 & 0.981 & \cellcolor{green!15}$+$0.731 & \cellcolor{green!15}0.974 & \cellcolor{green!15}$+$0.672 & \cellcolor{green!15}0.896 & \cellcolor{green!15}$+$0.671 & \cellcolor{green!15}0.895 \\
  & 5  & 0.974 & 0.973 & 0.975 & 0.949 & 0.974 & 0.974 & 0.948 & 0.942 & \cellcolor{green!15}$+$0.698 & \cellcolor{green!15}0.931 & \cellcolor{green!15}$+$0.629 & \cellcolor{green!15}0.839 & \cellcolor{green!15}$+$0.629 & \cellcolor{green!15}0.839 \\
  & 8  & 0.974 & 0.973 & 0.975 & 0.949 & 0.974 & 0.974 & 0.948 & 0.942 & \cellcolor{green!15}$+$0.698 & \cellcolor{green!15}0.931 & \cellcolor{green!15}$+$0.629 & \cellcolor{green!15}0.839 & \cellcolor{green!15}$+$0.629 & \cellcolor{green!15}0.839 \\
  & 10 & 0.942 & 0.938 & 0.945 & 0.889 & 0.940 & 0.942 & 0.883 & 0.877 & \cellcolor{green!15}$+$0.633 & \cellcolor{green!15}0.844 & \cellcolor{green!15}$+$0.552 & \cellcolor{green!15}0.736 & \cellcolor{green!15}$+$0.553 & \cellcolor{green!15}0.737 \\
  & 15 & 0.844 & 0.815 & 0.865 & 0.724 & 0.830 & 0.844 & 0.688 & 0.656 & \cellcolor{green!15}$+$0.438 & \cellcolor{green!15}0.584 & \cellcolor{green!15}$+$0.344 & \cellcolor{green!15}0.459 & \cellcolor{green!15}$+$0.346 & \cellcolor{green!15}0.462 \\
  & 20 & 0.669 & 0.519 & 0.748 & 0.432 & 0.592 & 0.669 & 0.338 & 0.286 & \cellcolor{green!15}$+$0.100 & \cellcolor{green!15}0.134 & \cellcolor{green!15}$+$0.016 & \cellcolor{green!15}0.021 & \cellcolor{green!15}$+$0.020 & \cellcolor{green!15}0.026 \\
  & 25 & 0.604 & 0.371 & 0.711 & 0.309 & 0.477 & 0.604 & 0.214 & 0.169 & $-$0.022 & 0.000 & $-$0.092 & 0.000 & $-$0.087 & 0.000 \\
  & 30 & 0.516 & 0.280 & 0.636 & 0.043 & 0.399 & 0.516 & 0.162 & 0.058 & $-$0.091 & 0.000 & $-$0.149 & 0.000 & $-$0.145 & 0.000 \\
  & 40 & 0.486 & 0.330 & 0.583 & -0.020 & 0.427 & 0.491 & 0.196 & 0.000 & $-$0.068 & 0.000 & $-$0.131 & 0.000 & $-$0.128 & 0.000 \\
  & 50 & 0.483 & 0.315 & 0.584 & -0.040 & 0.416 & 0.483 & 0.154 & 0.000 & $-$0.077 & 0.000 & $-$0.140 & 0.000 & $-$0.136 & 0.000 \\
  & 60 & 0.556 & 0.482 & 0.612 & 0.117 & 0.537 & 0.556 & 0.294 & 0.000 & \cellcolor{green!15}$+$0.039 & \cellcolor{green!15}0.051 & $-$0.046 & 0.000 & $-$0.042 & 0.000 \\
\midrule
\multirow{11}{*}{\makecell[l]{\textbf{Qwen3-32B}\\\textbf{(think)}\\($N^*=15$)}}
  & 3  & 0.979 & 0.979 & 0.980 & 0.959 & 0.979 & 0.979 & 0.959 & 0.959 & \cellcolor{green!15}$+$0.709 & \cellcolor{green!15}0.945 & \cellcolor{green!15}$+$0.629 & \cellcolor{green!15}0.839 & \cellcolor{green!15}$+$0.629 & \cellcolor{green!15}0.839 \\
  & 5  & 0.955 & 0.952 & 0.957 & 0.913 & 0.953 & 0.955 & 0.909 & 0.909 & \cellcolor{green!15}$+$0.659 & \cellcolor{green!15}0.879 & \cellcolor{green!15}$+$0.568 & \cellcolor{green!15}0.757 & \cellcolor{green!15}$+$0.569 & \cellcolor{green!15}0.758 \\
  & 8  & 0.926 & 0.920 & 0.931 & 0.861 & 0.923 & 0.926 & 0.851 & 0.843 & \cellcolor{green!15}$+$0.601 & \cellcolor{green!15}0.802 & \cellcolor{green!15}$+$0.502 & \cellcolor{green!15}0.669 & \cellcolor{green!15}$+$0.503 & \cellcolor{green!15}0.671 \\
  & 10 & 0.905 & 0.895 & 0.913 & 0.825 & 0.900 & 0.905 & 0.810 & 0.777 & \cellcolor{green!15}$+$0.560 & \cellcolor{green!15}0.747 & \cellcolor{green!15}$+$0.457 & \cellcolor{green!15}0.609 & \cellcolor{green!15}$+$0.458 & \cellcolor{green!15}0.611 \\
  & 15 & 0.748 & 0.663 & 0.799 & 0.574 & 0.704 & 0.748 & 0.496 & 0.430 & \cellcolor{green!15}$+$0.246 & \cellcolor{green!15}0.328 & \cellcolor{green!15}$+$0.142 & \cellcolor{green!15}0.189 & \cellcolor{green!15}$+$0.146 & \cellcolor{green!15}0.194 \\
  & 20 & 0.624 & 0.397 & 0.727 & 0.376 & 0.498 & 0.624 & 0.248 & 0.207 & $-$0.002 & 0.000 & $-$0.081 & 0.000 & $-$0.076 & 0.000 \\
  & 25 & 0.568 & 0.246 & 0.698 & 0.254 & 0.375 & 0.567 & 0.142 & 0.100 & $-$0.110 & 0.000 & $-$0.169 & 0.000 & $-$0.164 & 0.000 \\
  & 30 & 0.574 & 0.270 & 0.700 & 0.270 & 0.395 & 0.574 & 0.157 & 0.083 & $-$0.094 & 0.000 & $-$0.157 & 0.000 & $-$0.152 & 0.000 \\
  & 40 & 0.525 & 0.173 & 0.667 & 0.094 & 0.307 & 0.525 & 0.099 & 0.000 & $-$0.156 & 0.000 & $-$0.203 & 0.000 & $-$0.198 & 0.000 \\
  & 50 & 0.471 & 0.274 & 0.584 & -0.074 & 0.385 & 0.469 & 0.151 & 0.000 & $-$0.102 & 0.000 & $-$0.163 & 0.000 & $-$0.159 & 0.000 \\
  & 60 & 0.485 & 0.267 & 0.603 & -0.014 & 0.383 & 0.494 & 0.167 & 0.000 & $-$0.103 & 0.000 & $-$0.165 & 0.000 & $-$0.161 & 0.000 \\
\midrule
\multirow{11}{*}{\makecell[l]{\textbf{QwQ-32B}\\\textbf{(think)}\\($N^*=15$)}}
  & 3  & 1.000 & 1.000 & 1.000 & 1.000 & 1.000 & 1.000 & 1.000 & 1.000 & \cellcolor{green!15}$+$0.750 & \cellcolor{green!15}1.000 & \cellcolor{green!15}$+$0.678 & \cellcolor{green!15}0.904 & \cellcolor{green!15}$+$0.677 & \cellcolor{green!15}0.902 \\
  & 5  & 0.990 & 0.990 & 0.990 & 0.980 & 0.990 & 0.990 & 0.980 & 0.980 & \cellcolor{green!15}$+$0.730 & \cellcolor{green!15}0.973 & \cellcolor{green!15}$+$0.645 & \cellcolor{green!15}0.860 & \cellcolor{green!15}$+$0.645 & \cellcolor{green!15}0.859 \\
  & 8  & 0.980 & 0.980 & 0.980 & 0.960 & 0.980 & 0.980 & 0.960 & 0.960 & \cellcolor{green!15}$+$0.710 & \cellcolor{green!15}0.947 & \cellcolor{green!15}$+$0.614 & \cellcolor{green!15}0.818 & \cellcolor{green!15}$+$0.614 & \cellcolor{green!15}0.819 \\
  & 10 & 0.949 & 0.947 & 0.952 & 0.904 & 0.948 & 0.949 & 0.899 & 0.869 & \cellcolor{green!15}$+$0.649 & \cellcolor{green!15}0.865 & \cellcolor{green!15}$+$0.542 & \cellcolor{green!15}0.723 & \cellcolor{green!15}$+$0.543 & \cellcolor{green!15}0.724 \\
  & 15 & 0.803 & 0.755 & 0.835 & 0.659 & 0.778 & 0.803 & 0.606 & 0.556 & \cellcolor{green!15}$+$0.356 & \cellcolor{green!15}0.475 & \cellcolor{green!15}$+$0.234 & \cellcolor{green!15}0.313 & \cellcolor{green!15}$+$0.239 & \cellcolor{green!15}0.318 \\
  & 20 & 0.641 & 0.441 & 0.736 & 0.406 & 0.532 & 0.641 & 0.283 & 0.242 & \cellcolor{green!15}$+$0.033 & \cellcolor{green!15}0.044 & $-$0.060 & 0.000 & $-$0.054 & 0.000 \\
  & 25 & 0.576 & 0.276 & 0.700 & 0.270 & 0.400 & 0.576 & 0.162 & 0.111 & $-$0.090 & 0.000 & $-$0.160 & 0.000 & $-$0.154 & 0.000 \\
  & 30 & 0.515 & 0.059 & 0.673 & 0.124 & 0.174 & 0.515 & 0.030 & 0.000 & $-$0.220 & 0.000 & $-$0.244 & 0.000 & $-$0.240 & 0.000 \\
  & 40 & 0.487 & 0.093 & 0.642 & 0.043 & 0.221 & 0.509 & 0.033 & 0.011 & $-$0.201 & 0.000 & $-$0.235 & 0.000 & $-$0.230 & 0.000 \\
  & 50 & 0.506 & 0.044 & 0.667 & 0.044 & 0.150 & 0.506 & 0.023 & 0.000 & $-$0.228 & 0.000 & $-$0.247 & 0.000 & $-$0.244 & 0.000 \\
  & 60 & 0.500 & -- & 0.667 & -- & 0.000 & 0.500 & 0.000 & 0.000 & $-$0.250 & 0.000 & $-$0.250 & 0.000 & $-$0.250 & 0.000 \\
\midrule
\multirow{11}{*}{\makecell[l]{\textbf{gpt-oss-20b}\\\textbf{(reason)}\\($N^*=20$)}}
  & 3  & 1.000 & 1.000 & 1.000 & 1.000 & 1.000 & 1.000 & 1.000 & 1.000 & \cellcolor{green!15}$+$0.750 & \cellcolor{green!15}1.000 & \cellcolor{green!15}$+$0.550 & \cellcolor{green!15}0.733 & \cellcolor{green!15}$+$0.552 & \cellcolor{green!15}0.736 \\
  & 5  & 0.970 & 0.970 & 0.970 & 0.939 & 0.970 & 0.970 & 0.939 & 0.939 & \cellcolor{green!15}$+$0.690 & \cellcolor{green!15}0.920 & \cellcolor{green!15}$+$0.460 & \cellcolor{green!15}0.613 & \cellcolor{green!15}$+$0.467 & \cellcolor{green!15}0.623 \\
  & 8  & 0.924 & 0.918 & 0.930 & 0.858 & 0.921 & 0.924 & 0.848 & 0.818 & \cellcolor{green!15}$+$0.598 & \cellcolor{green!15}0.798 & \cellcolor{green!15}$+$0.359 & \cellcolor{green!15}0.479 & \cellcolor{green!15}$+$0.369 & \cellcolor{green!15}0.492 \\
  & 10 & 0.894 & 0.896 & 0.892 & 0.788 & 0.894 & 0.894 & 0.818 & 0.758 & \cellcolor{green!15}$+$0.549 & \cellcolor{green!15}0.732 & \cellcolor{green!15}$+$0.293 & \cellcolor{green!15}0.391 & \cellcolor{green!15}$+$0.305 & \cellcolor{green!15}0.407 \\
  & 15 & 0.848 & 0.833 & 0.861 & 0.709 & 0.844 & 0.848 & 0.697 & 0.636 & \cellcolor{green!15}$+$0.462 & \cellcolor{green!15}0.616 & \cellcolor{green!15}$+$0.211 & \cellcolor{green!15}0.281 & \cellcolor{green!15}$+$0.224 & \cellcolor{green!15}0.299 \\
  & 20 & 0.742 & 0.679 & 0.785 & 0.528 & 0.716 & 0.742 & 0.515 & 0.424 & \cellcolor{green!15}$+$0.262 & \cellcolor{green!15}0.350 & \cellcolor{green!15}$+$0.040 & \cellcolor{green!15}0.053 & \cellcolor{green!15}$+$0.055 & \cellcolor{green!15}0.074 \\
  & 25 & 0.652 & 0.596 & 0.693 & 0.315 & 0.637 & 0.652 & 0.364 & 0.242 & \cellcolor{green!15}$+$0.156 & \cellcolor{green!15}0.208 & $-$0.045 & 0.000 & $-$0.031 & 0.000 \\
  & 30 & 0.576 & 0.333 & 0.689 & 0.221 & 0.446 & 0.576 & 0.182 & 0.091 & $-$0.051 & 0.000 & $-$0.178 & 0.000 & $-$0.164 & 0.000 \\
  & 40 & 0.492 & 0.250 & 0.615 & 0.115 & 0.374 & 0.537 & 0.154 & 0.000 & $-$0.110 & 0.000 & $-$0.212 & 0.000 & $-$0.200 & 0.000 \\
  & 50 & 0.500 & 0.083 & 0.656 & 0.000 & 0.208 & 0.500 & 0.045 & 0.000 & $-$0.207 & 0.000 & $-$0.249 & 0.000 & $-$0.244 & 0.000 \\
  & 60 & 0.523 & 0.160 & 0.667 & 0.090 & 0.295 & 0.523 & 0.091 & 0.000 & $-$0.163 & 0.000 & $-$0.241 & 0.000 & $-$0.230 & 0.000 \\
\midrule
\multirow{11}{*}{\makecell[l]{\textbf{Nemotron-H-47B}\\\textbf{Reasoning-128K}\\\textbf{(r-on)}\\($N^*=0$)}}
  & 3  & 0.525 & 0.663 & 0.200 & 0.087 & 0.332 & 0.525 & 0.068 & 0.038 & $-$0.140 & 0.000 & $-$0.160 & 0.000 & $-$0.159 & 0.000 \\
  & 5  & 0.513 & 0.670 & 0.071 & 0.082 & 0.191 & 0.513 & 0.030 & 0.002 & $-$0.213 & 0.000 & $-$0.225 & 0.000 & $-$0.224 & 0.000 \\
  & 8  & 0.500 & 0.666 & 0.002 & 0.000 & 0.032 & 0.500 & 0.001 & 0.000 & $-$0.249 & 0.000 & $-$0.250 & 0.000 & $-$0.250 & 0.000 \\
  & 10 & 0.497 & 0.664 & 0.002 & -0.049 & 0.032 & 0.497 & 0.001 & 0.000 & $-$0.249 & 0.000 & $-$0.250 & 0.000 & $-$0.250 & 0.000 \\
  & 15 & 0.512 & 0.669 & 0.072 & 0.072 & 0.194 & 0.512 & 0.026 & 0.000 & $-$0.213 & 0.000 & $-$0.224 & 0.000 & $-$0.223 & 0.000 \\
  & 20 & 0.511 & 0.666 & 0.091 & 0.053 & 0.219 & 0.510 & 0.034 & 0.000 & $-$0.202 & 0.000 & $-$0.215 & 0.000 & $-$0.214 & 0.000 \\
  & 25 & 0.520 & 0.663 & 0.167 & 0.074 & 0.301 & 0.520 & 0.068 & 0.000 & $-$0.159 & 0.000 & $-$0.178 & 0.000 & $-$0.177 & 0.000 \\
  & 30 & 0.499 & 0.660 & 0.047 & -0.007 & 0.155 & 0.499 & 0.020 & 0.000 & $-$0.226 & 0.000 & $-$0.235 & 0.000 & $-$0.234 & 0.000 \\
  & 40 & 0.518 & 0.670 & 0.105 & 0.094 & 0.236 & 0.518 & 0.052 & 0.000 & $-$0.194 & 0.000 & $-$0.209 & 0.000 & $-$0.208 & 0.000 \\
  & 50 & 0.509 & 0.661 & 0.110 & 0.038 & 0.241 & 0.509 & 0.048 & 0.000 & $-$0.192 & 0.000 & $-$0.206 & 0.000 & $-$0.205 & 0.000 \\
  & 60 & 0.521 & 0.666 & 0.156 & 0.085 & 0.291 & 0.521 & 0.073 & 0.000 & $-$0.165 & 0.000 & $-$0.183 & 0.000 & $-$0.182 & 0.000 \\
\midrule
\multirow{11}{*}{\makecell[l]{\textbf{Llama-Nemotron}\\\textbf{Super-49B v1.5}\\\textbf{(r-on)}\\($N^*=10$)}}
  & 3  & 1.000 & 1.000 & 1.000 & 1.000 & 1.000 & 1.000 & 1.000 & 0.955 & \cellcolor{green!15}$+$0.750 & \cellcolor{green!15}1.000 & \cellcolor{green!15}$+$0.465 & \cellcolor{green!15}0.620 & \cellcolor{green!15}$+$0.475 & \cellcolor{green!15}0.633 \\
  & 5  & 0.977 & 0.978 & 0.977 & 0.956 & 0.977 & 0.977 & 0.955 & 0.955 & \cellcolor{green!15}$+$0.705 & \cellcolor{green!15}0.939 & \cellcolor{green!15}$+$0.402 & \cellcolor{green!15}0.537 & \cellcolor{green!15}$+$0.416 & \cellcolor{green!15}0.554 \\
  & 8  & 0.955 & 0.952 & 0.957 & 0.913 & 0.953 & 0.955 & 0.909 & 0.909 & \cellcolor{green!15}$+$0.659 & \cellcolor{green!15}0.879 & \cellcolor{green!15}$+$0.349 & \cellcolor{green!15}0.465 & \cellcolor{green!15}$+$0.365 & \cellcolor{green!15}0.486 \\
  & 10 & 0.864 & 0.842 & 0.880 & 0.756 & 0.853 & 0.864 & 0.727 & 0.727 & \cellcolor{green!15}$+$0.477 & \cellcolor{green!15}0.636 & \cellcolor{green!15}$+$0.171 & \cellcolor{green!15}0.228 & \cellcolor{green!15}$+$0.191 & \cellcolor{green!15}0.255 \\
  & 15 & 0.727 & 0.625 & 0.786 & 0.542 & 0.674 & 0.727 & 0.455 & 0.318 & \cellcolor{green!15}$+$0.205 & \cellcolor{green!15}0.273 & $-$0.044 & 0.000 & $-$0.021 & 0.000 \\
  & 20 & 0.750 & 0.667 & 0.800 & 0.577 & 0.707 & 0.750 & 0.500 & 0.455 & \cellcolor{green!15}$+$0.250 & \cellcolor{green!15}0.333 & $-$0.011 & 0.000 & \cellcolor{green!15}$+$0.012 & \cellcolor{green!15}0.015 \\
  & 25 & 0.727 & 0.625 & 0.786 & 0.542 & 0.674 & 0.727 & 0.455 & 0.136 & \cellcolor{green!15}$+$0.205 & \cellcolor{green!15}0.273 & $-$0.044 & 0.000 & $-$0.021 & 0.000 \\
  & 30 & 0.614 & 0.485 & 0.691 & 0.262 & 0.560 & 0.614 & 0.273 & 0.045 & \cellcolor{green!15}$+$0.064 & \cellcolor{green!15}0.085 & $-$0.138 & 0.000 & $-$0.118 & 0.000 \\
  & 40 & 0.591 & 0.357 & 0.700 & 0.265 & 0.466 & 0.591 & 0.182 & 0.000 & $-$0.033 & 0.000 & $-$0.190 & 0.000 & $-$0.171 & 0.000 \\
  & 50 & 0.568 & 0.457 & 0.642 & 0.149 & 0.530 & 0.568 & 0.318 & 0.000 & \cellcolor{green!15}$+$0.031 & \cellcolor{green!15}0.041 & $-$0.156 & 0.000 & $-$0.138 & 0.000 \\
  & 60 & 0.400 & 0.222 & 0.512 & -0.024 & 0.340 & 0.491 & 0.077 & 0.000 & $-$0.135 & 0.000 & $-$0.234 & 0.000 & $-$0.223 & 0.000 \\
\bottomrule
\multicolumn{16}{l}{\colorbox{green!15}{\phantom{X}} = that estimator's $\RMC^{\pm}>0$.}
\end{tabular}%
}
\end{table*}

\begin{table*}[t]
\centering
\caption{Conditional correctness of the UNSAT-clause-subset follow-up on 2-SAT open-source models.
For each $(\text{model},N)$ cell, the denominator includes only follow-up outputs whose parsed result is \texttt{UNSAT\_CLAUSE\_SUBSET\_FOUND}; therefore this is a conditional accuracy of the subset itself rather than an ADR-style end-to-end metric.
Cells are shown as ``correct/found = rate''.
These values are generated by [compute\_unsat\_clause\_subset\_found\_accuracy.py](/scratch/c00590656/loni/Group_Evaluation/SAT_Group_Evaluation/analysis/open_source_LLM_2SAT_analysis/unsat_clause_subset/compute_unsat_clause_subset_found_accuracy.py) and written to [unsat\_clause\_subset\_found\_accuracy.by\_N.csv](/scratch/c00590656/loni/Group_Evaluation/SAT_Group_Evaluation/analysis/open_source_LLM_2SAT_analysis/unsat_clause_subset/unsat_clause_subset_found_accuracy.by_N.csv).}
\label{tab:os_unsat_subset_conditional}
\renewcommand{\arraystretch}{1.08}
\resizebox{\textwidth}{!}{%
\begin{tabular}{p{3.3cm}ccccccccccc}
\toprule
Model & $N{=}3$ & $N{=}5$ & $N{=}8$ & $N{=}10$ & $N{=}15$ & $N{=}20$ & $N{=}25$ & $N{=}30$ & $N{=}40$ & $N{=}50$ & $N{=}60$ \\
\midrule
\makecell[l]{Qwen3-14B\\(think)}
& 24/44 = 0.545
& 20/44 = 0.455
& 8/43 = 0.186
& 20/43 = 0.465
& 9/47 = 0.191
& 0/36 = 0.000
& 1/44 = 0.023
& 1/39 = 0.026
& 4/44 = 0.091
& 0/37 = 0.000
& 0/28 = 0.000 \\
\makecell[l]{QwQ-32B\\(think)}
& 5/10 = 0.500
& 3/8 = 0.375
& 4/10 = 0.400
& 6/11 = 0.545
& 5/16 = 0.312
& 4/19 = 0.211
& 1/22 = 0.045
& 1/25 = 0.040
& 0/17 = 0.000
& 1/14 = 0.071
& 0/9 = 0.000 \\
\makecell[l]{Qwen3-30B-A3B\\Instruct-2507\\(think)}
& 70/5773 = 0.012
& --
& 8/898 = 0.009
& --
& 20/327 = 0.061
& 7/109 = 0.064
& --
& 5/48 = 0.104
& --
& --
& -- \\
\bottomrule
\end{tabular}%
}
\end{table*}

\paragraph{UNSAT-subset follow-up accuracy conditioned on subset emission.}
The conditional verification rates in \Cref{tab:os_unsat_subset_conditional} isolate a different question from ADR: when a model already commits to \texttt{UNSAT\_CLAUSE\_SUBSET\_FOUND}, how often is the returned clause subset actually UNSAT?
Under this conditional metric, Qwen3-14B and QwQ-32B are both around $0.45$--$0.55$ at small $N$ but collapse rapidly as $N$ grows, while Qwen3-30B-A3B-Instruct remains low even when conditioned on subset emission itself.
Qwen3-32B is not included in this subtable because the current workspace does not contain a completed UNSAT-clause-subset follow-up output directory for that model.

\paragraph{Finding 4: Qwen3-14B achieves the highest open-weight 2-SAT construction-level boundary ($N^*=20$).}
Qwen3-14B holds strong margins through $N=15$ (RMC(CP) $=+0.344$, RMC(Wilson) $=+0.346$, RMC(no LCB) $=+0.438$, F1$_{\mathrm{SAT}}=0.815$, F1$_{\mathrm{UNSAT}}=0.865$) and retains a marginal construction-level certificate at $N=20$ (RMC(CP) $+0.015$, RMC(Wilson) $+0.019$, RMC(no LCB) $+0.100$), giving $N^*=20$.
F1$_{\mathrm{SAT}}$ nevertheless drops sharply to $0.519$ at $N=20$.  Under the recalibrated residual Case B ceiling, the model still keeps a small positive certificate at $N=20$ (RMC(CP) $+0.041$), before failing at $N=25$; thus the Case B boundary now matches the construction-level boundary at $N=20$ (\Cref{tab:os_pern_caseB}).
Because 2-SAT and 3-SAT belong to different complexity regimes, this boundary is not comparable to the commercial 3-SAT boundaries.

\paragraph{Finding 5: QwQ-32B achieves the strongest GP margins at small $N$ among open-source models.}
At $N=3$, QwQ-32B achieves RMC(CP) $+0.678$, RMC(no LCB) $+0.750$ (theoretical maximum, nRMC$=1.000$), RMC(Wilson) $+0.677$, F1$_{\mathrm{SAT}}=1.000$, F1$_{\mathrm{UNSAT}}=1.000$---numerically similar to GPT-5's small-$N$ 3-SAT margins, although the two task families are not directly comparable.
Its construction-level boundary is $N^*=15$ (RMC(CP) $+0.234$ at $N=15$, negative at $N=20$; RMC(CP) $+0.542$ and RMC(Wilson) $+0.543$ at $N=10$), and its calibrated Case B boundary is also $N=15$ (\Cref{tab:os_pern_caseB}).

\paragraph{Finding 6: Qwen3-32B has the weakest GP margins among the three certified open-source models but consistent behavior across estimators.}
Qwen3-32B achieves $N^*=15$: at $N=15$, RMC(CP) drops to $+0.142$ before going fully negative at $N=20$; at $N=10$ its margins are RMC(CP) $+0.457$, RMC(Wilson) $+0.458$, RMC(no LCB) $+0.560$, with F1$_{\mathrm{SAT}}=0.895$ and F1$_{\mathrm{UNSAT}}=0.913$.
The close agreement between RMC(CP) and RMC(Wilson) confirms that CP and Wilson LCBs are interchangeable at the group sizes used here.
Under the recalibrated residual Case B ceiling, the boundary remains at $N=15$, with Case B RMC(CP) still positive at $N=15$ ($+0.245$) and turning negative only at $N=20$ (\Cref{tab:os_pern_caseB}).

\paragraph{Additional open-weight baselines: short-range certification and pilot sensitivity.}
The newly added open-weight reasoning baselines clarify that the sustained Qwen-family certificates are not a generic consequence of enabling reasoning mode.
gpt-oss-20b certifies only through $N=8$ under Case A: its pooled CP margin is $+0.095$ at $N=8$, already turns negative at $N=10$ ($-0.017$), and remains negative thereafter, although the Wilson pooled estimate stays barely positive at $N=10$ ($+0.024$).
Nemotron-H-47B-Reasoning-128K never clears the Case-A pooled CP gate, with its best small-$N$ margin still negative ($-0.101$ at $N=3$), reflecting strong SAT-positive but weak UNSAT discrimination.
Llama-Nemotron-Super-49B v1.5 is based on a very small local pilot: it is positive at $N=3$ (Case-A pooled CP $+0.261$) and negative at $N=5$ under CP ($-0.058$), though the corresponding Wilson row remains slightly positive ($+0.017$).

\paragraph{Case B sensitivity: CNF side-channel adversary battery (2-SAT).}
Applying the same final residual KS adversary battery as in \Cref{sec:rq2}, with permutation calibration for the max-over-features bias, gives calibrated Case B ceilings for the 2-SAT setting.
Across the six open-weight reasoning-mode runs with current local 2-SAT data, the latest per-$N$ mean residual $\shortcutCeiling$ values range from $0.024$ to $0.313$.
Under the updated Case B definition, these are residual shortcut contributions above the Case-A chance floor, with $\hatshortcutCeiling=\bigl((1+\hatshortcutDisc)/2\bigr)^2-0.25$.  They correspond to calibrated residual distinguishability $\hatshortcutDisc$ between roughly $0.19$ and $0.49$.  The 2-SAT minimal-edit construction therefore leaves measurable residual surface leakage, which is precisely why the calibrated Case B certificate, rather than the Case A idealization, is the primary certificate.
\Cref{tab:os_pern_caseB} reports the Case B results.

Under this calibrated Case B ceiling, the global-pooled CP margin remains positive for Qwen3-14B through $N=20$, for Qwen3-32B through $N=15$, for QwQ-32B through $N=15$, for gpt-oss-20b through $N=20$, and for Llama-Nemotron-Super-49B v1.5 (reasoning-on) through $N=15$; Nemotron-H-47B-Reasoning-128K (reasoning-on) remains negative throughout.
These Case B boundaries pair with the construction-level boundaries $N^*=20$, $15$, and $15$ reported in \Cref{tab:os_pern} for the three Qwen-family models, so Qwen3-14B now retains its certificate through the same tested boundary under Case B, while Qwen3-32B and QwQ-32B still tighten relative to the easier small-$N$ regime.
At larger $N$, the models fail because the statistical lower bound on within-group separability drops below the estimated shortcut ceiling.

\begin{table*}[ht]
\centering
\caption{Open-weight 2-SAT per-$N$ metrics under Case B (estimated $\shortcutCeiling$).
Values come from the latest local Case-B run with the final residual CNF side-channel feature set and permutation calibration.
$\hat{s}$ is the point-estimate separability and $\mathrm{nRMC}=\max(0,\RMC^{\pm})/(1-\shortcutCeiling)$.}
\label{tab:os_pern_caseB}
\renewcommand{\arraystretch}{1.12}
\resizebox{\textwidth}{!}{%
\begin{tabular}{p{2.2cm}r rrrr rr rr rr rr rr rr}
\toprule
\multirow{2}{*}{Model} & \multirow{2}{*}{$N$}
  & \multirow{2}{*}{Acc}
  & \multicolumn{2}{c}{F1 (per class)}
  & \multirow{2}{*}{MCC}
  & \multirow{2}{*}{G-mean}
  & \multirow{2}{*}{BalAcc}
  & \multirow{2}{*}{ADR} & \multirow{2}{*}{ADR$^{+w}$}
  & \multirow{2}{*}{$\shortcutCeiling$}
  & \multirow{2}{*}{$\hat{s}$}
  & \multicolumn{2}{c}{RMC(no LCB)}
  & \multicolumn{2}{c}{RMC(CP)}
  & \multicolumn{2}{c}{RMC(Wilson)} \\
\cmidrule(lr){4-5}\cmidrule(lr){13-14}\cmidrule(lr){15-16}\cmidrule(lr){17-18}
 & & & F1$_{\mathrm{SAT}}$ & F1$_{\mathrm{UNSAT}}$ & MCC & & & & & & &
  $\RMC$ & $\mathrm{nRMC}$
  & $\RMC^{\pm}$ & $\mathrm{nRMC}$
  & $\RMC^{\pm}$ & $\mathrm{nRMC}$ \\
\midrule
\multirow{11}{*}{\makecell[l]{\textbf{Qwen3-14B}\\\textbf{(think)}\\($N^*=20$)}}
  & 3  & 0.990 & 0.990 & 0.990 & 0.981 & 0.990 & 0.990 & 0.981 & 0.981 & 0.162 & 0.981 & \cellcolor{green!15}$+$0.819 & \cellcolor{green!15}0.977 & \cellcolor{green!15}$+$0.760 & \cellcolor{green!15}0.907 & \cellcolor{green!15}$+$0.760 & \cellcolor{green!15}0.906 \\
  & 5  & 0.974 & 0.973 & 0.975 & 0.949 & 0.974 & 0.974 & 0.948 & 0.942 & 0.129 & 0.948 & \cellcolor{green!15}$+$0.819 & \cellcolor{green!15}0.940 & \cellcolor{green!15}$+$0.750 & \cellcolor{green!15}0.861 & \cellcolor{green!15}$+$0.750 & \cellcolor{green!15}0.861 \\
  & 8  & 0.974 & 0.973 & 0.975 & 0.949 & 0.974 & 0.974 & 0.948 & 0.942 & 0.167 & 0.948 & \cellcolor{green!15}$+$0.781 & \cellcolor{green!15}0.938 & \cellcolor{green!15}$+$0.712 & \cellcolor{green!15}0.855 & \cellcolor{green!15}$+$0.712 & \cellcolor{green!15}0.855 \\
  & 10 & 0.942 & 0.938 & 0.945 & 0.889 & 0.940 & 0.942 & 0.883 & 0.877 & 0.171 & 0.883 & \cellcolor{green!15}$+$0.712 & \cellcolor{green!15}0.859 & \cellcolor{green!15}$+$0.631 & \cellcolor{green!15}0.761 & \cellcolor{green!15}$+$0.632 & \cellcolor{green!15}0.762 \\
  & 15 & 0.844 & 0.815 & 0.865 & 0.724 & 0.830 & 0.844 & 0.688 & 0.656 & 0.125 & 0.688 & \cellcolor{green!15}$+$0.563 & \cellcolor{green!15}0.644 & \cellcolor{green!15}$+$0.469 & \cellcolor{green!15}0.536 & \cellcolor{green!15}$+$0.471 & \cellcolor{green!15}0.539 \\
  & 20 & 0.669 & 0.519 & 0.748 & 0.432 & 0.592 & 0.669 & 0.338 & 0.286 & 0.225 & 0.350 & \cellcolor{green!15}$+$0.125 & \cellcolor{green!15}0.161 & \cellcolor{green!15}$+$0.041 & \cellcolor{green!15}0.053 & \cellcolor{green!15}$+$0.045 & \cellcolor{green!15}0.058 \\
  & 25 & 0.604 & 0.371 & 0.711 & 0.309 & 0.477 & 0.604 & 0.214 & 0.169 & 0.206 & 0.228 & \cellcolor{green!15}$+$0.022 & \cellcolor{green!15}0.028 & $-$0.047 & 0.000 & $-$0.043 & 0.000 \\
  & 30 & 0.516 & 0.280 & 0.636 & 0.043 & 0.399 & 0.516 & 0.162 & 0.058 & 0.227 & 0.159 & $-$0.068 & 0.000 & $-$0.126 & 0.000 & $-$0.122 & 0.000 \\
  & 40 & 0.486 & 0.330 & 0.583 & -0.020 & 0.427 & 0.491 & 0.196 & 0.000 & 0.280 & 0.182 & $-$0.098 & 0.000 & $-$0.162 & 0.000 & $-$0.158 & 0.000 \\
  & 50 & 0.483 & 0.315 & 0.584 & -0.040 & 0.416 & 0.483 & 0.154 & 0.000 & 0.244 & 0.173 & $-$0.071 & 0.000 & $-$0.134 & 0.000 & $-$0.130 & 0.000 \\
  & 60 & 0.556 & 0.482 & 0.612 & 0.117 & 0.537 & 0.556 & 0.294 & 0.000 & 0.276 & 0.289 & \cellcolor{green!15}$+$0.013 & \cellcolor{green!15}0.017 & $-$0.072 & 0.000 & $-$0.068 & 0.000 \\
\midrule
\multirow{11}{*}{\makecell[l]{\textbf{Qwen3-32B}\\\textbf{(think)}\\($N^*=15$)}}
  & 3  & 0.979 & 0.979 & 0.980 & 0.959 & 0.979 & 0.979 & 0.959 & 0.959 & 0.167 & 0.959 & \cellcolor{green!15}$+$0.792 & \cellcolor{green!15}0.950 & \cellcolor{green!15}$+$0.712 & \cellcolor{green!15}0.855 & \cellcolor{green!15}$+$0.712 & \cellcolor{green!15}0.855 \\
  & 5  & 0.955 & 0.952 & 0.957 & 0.913 & 0.953 & 0.955 & 0.909 & 0.909 & 0.107 & 0.909 & \cellcolor{green!15}$+$0.802 & \cellcolor{green!15}0.898 & \cellcolor{green!15}$+$0.711 & \cellcolor{green!15}0.796 & \cellcolor{green!15}$+$0.712 & \cellcolor{green!15}0.797 \\
  & 8  & 0.926 & 0.920 & 0.931 & 0.861 & 0.923 & 0.926 & 0.851 & 0.843 & 0.163 & 0.851 & \cellcolor{green!15}$+$0.688 & \cellcolor{green!15}0.822 & \cellcolor{green!15}$+$0.589 & \cellcolor{green!15}0.704 & \cellcolor{green!15}$+$0.590 & \cellcolor{green!15}0.705 \\
  & 10 & 0.905 & 0.895 & 0.913 & 0.825 & 0.900 & 0.905 & 0.810 & 0.777 & 0.168 & 0.810 & \cellcolor{green!15}$+$0.642 & \cellcolor{green!15}0.771 & \cellcolor{green!15}$+$0.538 & \cellcolor{green!15}0.647 & \cellcolor{green!15}$+$0.540 & \cellcolor{green!15}0.649 \\
  & 15 & 0.748 & 0.663 & 0.799 & 0.574 & 0.704 & 0.748 & 0.496 & 0.430 & 0.147 & 0.496 & \cellcolor{green!15}$+$0.349 & \cellcolor{green!15}0.409 & \cellcolor{green!15}$+$0.245 & \cellcolor{green!15}0.287 & \cellcolor{green!15}$+$0.249 & \cellcolor{green!15}0.292 \\
  & 20 & 0.624 & 0.397 & 0.727 & 0.376 & 0.498 & 0.624 & 0.248 & 0.207 & 0.198 & 0.248 & \cellcolor{green!15}$+$0.050 & \cellcolor{green!15}0.063 & $-$0.029 & 0.000 & $-$0.024 & 0.000 \\
  & 25 & 0.568 & 0.246 & 0.698 & 0.254 & 0.375 & 0.567 & 0.142 & 0.100 & 0.178 & 0.140 & $-$0.038 & 0.000 & $-$0.098 & 0.000 & $-$0.092 & 0.000 \\
  & 30 & 0.574 & 0.270 & 0.700 & 0.270 & 0.395 & 0.574 & 0.157 & 0.083 & 0.233 & 0.156 & $-$0.077 & 0.000 & $-$0.140 & 0.000 & $-$0.135 & 0.000 \\
  & 40 & 0.525 & 0.173 & 0.667 & 0.094 & 0.307 & 0.525 & 0.099 & 0.000 & 0.304 & 0.094 & $-$0.209 & 0.000 & $-$0.257 & 0.000 & $-$0.252 & 0.000 \\
  & 50 & 0.471 & 0.274 & 0.584 & -0.074 & 0.385 & 0.469 & 0.151 & 0.000 & 0.234 & 0.148 & $-$0.086 & 0.000 & $-$0.148 & 0.000 & $-$0.144 & 0.000 \\
  & 60 & 0.485 & 0.267 & 0.603 & -0.014 & 0.383 & 0.494 & 0.167 & 0.000 & 0.238 & 0.147 & $-$0.091 & 0.000 & $-$0.153 & 0.000 & $-$0.148 & 0.000 \\
\midrule
\multirow{11}{*}{\makecell[l]{\textbf{QwQ-32B}\\\textbf{(think)}\\($N^*=15$)}}
  & 3  & 1.000 & 1.000 & 1.000 & 1.000 & 1.000 & 1.000 & 1.000 & 1.000 & 0.142 & 1.000 & \cellcolor{green!15}$+$0.858 & \cellcolor{green!15}1.000 & \cellcolor{green!15}$+$0.786 & \cellcolor{green!15}0.916 & \cellcolor{green!15}$+$0.785 & \cellcolor{green!15}0.915 \\
  & 5  & 0.990 & 0.990 & 0.990 & 0.980 & 0.990 & 0.990 & 0.980 & 0.980 & 0.118 & 0.980 & \cellcolor{green!15}$+$0.862 & \cellcolor{green!15}0.977 & \cellcolor{green!15}$+$0.777 & \cellcolor{green!15}0.881 & \cellcolor{green!15}$+$0.777 & \cellcolor{green!15}0.880 \\
  & 8  & 0.980 & 0.980 & 0.980 & 0.960 & 0.980 & 0.980 & 0.960 & 0.960 & 0.156 & 0.960 & \cellcolor{green!15}$+$0.804 & \cellcolor{green!15}0.952 & \cellcolor{green!15}$+$0.708 & \cellcolor{green!15}0.839 & \cellcolor{green!15}$+$0.708 & \cellcolor{green!15}0.839 \\
  & 10 & 0.949 & 0.947 & 0.952 & 0.904 & 0.948 & 0.949 & 0.899 & 0.869 & 0.157 & 0.899 & \cellcolor{green!15}$+$0.742 & \cellcolor{green!15}0.880 & \cellcolor{green!15}$+$0.635 & \cellcolor{green!15}0.753 & \cellcolor{green!15}$+$0.636 & \cellcolor{green!15}0.755 \\
  & 15 & 0.803 & 0.755 & 0.835 & 0.659 & 0.778 & 0.803 & 0.606 & 0.556 & 0.161 & 0.606 & \cellcolor{green!15}$+$0.445 & \cellcolor{green!15}0.530 & \cellcolor{green!15}$+$0.323 & \cellcolor{green!15}0.385 & \cellcolor{green!15}$+$0.327 & \cellcolor{green!15}0.390 \\
  & 20 & 0.641 & 0.441 & 0.736 & 0.406 & 0.532 & 0.641 & 0.283 & 0.242 & 0.208 & 0.283 & \cellcolor{green!15}$+$0.074 & \cellcolor{green!15}0.094 & $-$0.019 & 0.000 & $-$0.013 & 0.000 \\
  & 25 & 0.576 & 0.276 & 0.700 & 0.270 & 0.400 & 0.576 & 0.162 & 0.111 & 0.194 & 0.160 & $-$0.034 & 0.000 & $-$0.104 & 0.000 & $-$0.098 & 0.000 \\
  & 30 & 0.515 & 0.059 & 0.673 & 0.124 & 0.174 & 0.515 & 0.030 & 0.000 & 0.223 & 0.030 & $-$0.193 & 0.000 & $-$0.217 & 0.000 & $-$0.213 & 0.000 \\
  & 40 & 0.487 & 0.093 & 0.642 & 0.043 & 0.221 & 0.509 & 0.033 & 0.011 & 0.286 & 0.049 & $-$0.237 & 0.000 & $-$0.271 & 0.000 & $-$0.266 & 0.000 \\
  & 50 & 0.506 & 0.044 & 0.667 & 0.044 & 0.150 & 0.506 & 0.023 & 0.000 & 0.232 & 0.022 & $-$0.210 & 0.000 & $-$0.230 & 0.000 & $-$0.227 & 0.000 \\
  & 60 & 0.500 & -- & 0.667 & -- & 0.000 & 0.500 & 0.000 & 0.000 & 0.272 & 0.000 & $-$0.272 & 0.000 & $-$0.272 & 0.000 & $-$0.272 & 0.000 \\
\midrule
\multirow{11}{*}{\makecell[l]{\textbf{gpt-oss-20b}\\\textbf{(reason)}\\($N^*=20$)}}
  & 3  & 1.000 & 1.000 & 1.000 & 1.000 & 1.000 & 1.000 & 1.000 & 1.000 & 0.109 & 1.000 & \cellcolor{green!15}$+$0.891 & \cellcolor{green!15}1.000 & \cellcolor{green!15}$+$0.691 & \cellcolor{green!15}0.775 & \cellcolor{green!15}$+$0.694 & \cellcolor{green!15}0.778 \\
  & 5  & 0.970 & 0.970 & 0.970 & 0.939 & 0.970 & 0.970 & 0.939 & 0.939 & 0.067 & 0.940 & \cellcolor{green!15}$+$0.873 & \cellcolor{green!15}0.936 & \cellcolor{green!15}$+$0.642 & \cellcolor{green!15}0.689 & \cellcolor{green!15}$+$0.650 & \cellcolor{green!15}0.697 \\
  & 8  & 0.924 & 0.918 & 0.930 & 0.858 & 0.921 & 0.924 & 0.848 & 0.818 & 0.126 & 0.848 & \cellcolor{green!15}$+$0.722 & \cellcolor{green!15}0.827 & \cellcolor{green!15}$+$0.483 & \cellcolor{green!15}0.553 & \cellcolor{green!15}$+$0.493 & \cellcolor{green!15}0.564 \\
  & 10 & 0.894 & 0.896 & 0.892 & 0.788 & 0.894 & 0.894 & 0.818 & 0.758 & 0.209 & 0.799 & \cellcolor{green!15}$+$0.590 & \cellcolor{green!15}0.746 & \cellcolor{green!15}$+$0.335 & \cellcolor{green!15}0.423 & \cellcolor{green!15}$+$0.347 & \cellcolor{green!15}0.438 \\
  & 15 & 0.848 & 0.833 & 0.861 & 0.709 & 0.844 & 0.848 & 0.697 & 0.636 & 0.175 & 0.712 & \cellcolor{green!15}$+$0.537 & \cellcolor{green!15}0.651 & \cellcolor{green!15}$+$0.286 & \cellcolor{green!15}0.346 & \cellcolor{green!15}$+$0.299 & \cellcolor{green!15}0.363 \\
  & 20 & 0.742 & 0.679 & 0.785 & 0.528 & 0.716 & 0.742 & 0.515 & 0.424 & 0.242 & 0.512 & \cellcolor{green!15}$+$0.271 & \cellcolor{green!15}0.357 & \cellcolor{green!15}$+$0.048 & \cellcolor{green!15}0.064 & \cellcolor{green!15}$+$0.064 & \cellcolor{green!15}0.084 \\
  & 25 & 0.652 & 0.596 & 0.693 & 0.315 & 0.637 & 0.652 & 0.364 & 0.242 & 0.302 & 0.406 & \cellcolor{green!15}$+$0.104 & \cellcolor{green!15}0.149 & $-$0.097 & 0.000 & $-$0.082 & 0.000 \\
  & 30 & 0.576 & 0.333 & 0.689 & 0.221 & 0.446 & 0.576 & 0.182 & 0.091 & 0.260 & 0.199 & $-$0.061 & 0.000 & $-$0.189 & 0.000 & $-$0.175 & 0.000 \\
  & 40 & 0.492 & 0.250 & 0.615 & 0.115 & 0.374 & 0.537 & 0.154 & 0.000 & 0.252 & 0.140 & $-$0.113 & 0.000 & $-$0.214 & 0.000 & $-$0.202 & 0.000 \\
  & 50 & 0.500 & 0.083 & 0.656 & 0.000 & 0.208 & 0.500 & 0.045 & 0.000 & 0.247 & 0.043 & $-$0.204 & 0.000 & $-$0.246 & 0.000 & $-$0.241 & 0.000 \\
  & 60 & 0.523 & 0.160 & 0.667 & 0.090 & 0.295 & 0.523 & 0.091 & 0.000 & 0.313 & 0.087 & $-$0.226 & 0.000 & $-$0.304 & 0.000 & $-$0.293 & 0.000 \\
\midrule
\multirow{11}{*}{\makecell[l]{\textbf{Nemotron-H-47B}\\\textbf{Reasoning-128K}\\\textbf{(r-on)}\\($N^*=0$)}}
  & 3  & 0.525 & 0.663 & 0.200 & 0.087 & 0.332 & 0.525 & 0.068 & 0.038 & 0.141 & 0.110 & $-$0.030 & 0.000 & $-$0.051 & 0.000 & $-$0.050 & 0.000 \\
  & 5  & 0.513 & 0.670 & 0.071 & 0.082 & 0.191 & 0.513 & 0.030 & 0.002 & 0.151 & 0.037 & $-$0.115 & 0.000 & $-$0.126 & 0.000 & $-$0.125 & 0.000 \\
  & 8  & 0.500 & 0.666 & 0.002 & 0.000 & 0.032 & 0.500 & 0.001 & 0.000 & 0.159 & 0.001 & $-$0.158 & 0.000 & $-$0.159 & 0.000 & $-$0.159 & 0.000 \\
  & 10 & 0.497 & 0.664 & 0.002 & -0.049 & 0.032 & 0.497 & 0.001 & 0.000 & 0.160 & 0.001 & $-$0.159 & 0.000 & $-$0.160 & 0.000 & $-$0.159 & 0.000 \\
  & 15 & 0.512 & 0.669 & 0.072 & 0.072 & 0.194 & 0.512 & 0.026 & 0.000 & 0.192 & 0.037 & $-$0.154 & 0.000 & $-$0.166 & 0.000 & $-$0.165 & 0.000 \\
  & 20 & 0.511 & 0.666 & 0.091 & 0.053 & 0.219 & 0.510 & 0.034 & 0.000 & 0.217 & 0.048 & $-$0.169 & 0.000 & $-$0.182 & 0.000 & $-$0.181 & 0.000 \\
  & 25 & 0.520 & 0.663 & 0.167 & 0.074 & 0.301 & 0.520 & 0.068 & 0.000 & 0.212 & 0.091 & $-$0.121 & 0.000 & $-$0.140 & 0.000 & $-$0.139 & 0.000 \\
  & 30 & 0.499 & 0.660 & 0.047 & -0.007 & 0.155 & 0.499 & 0.020 & 0.000 & 0.225 & 0.024 & $-$0.201 & 0.000 & $-$0.210 & 0.000 & $-$0.209 & 0.000 \\
  & 40 & 0.518 & 0.670 & 0.105 & 0.094 & 0.236 & 0.518 & 0.052 & 0.000 & 0.245 & 0.056 & $-$0.189 & 0.000 & $-$0.204 & 0.000 & $-$0.203 & 0.000 \\
  & 50 & 0.509 & 0.661 & 0.110 & 0.038 & 0.241 & 0.509 & 0.048 & 0.000 & 0.242 & 0.058 & $-$0.183 & 0.000 & $-$0.198 & 0.000 & $-$0.197 & 0.000 \\
  & 60 & 0.521 & 0.666 & 0.156 & 0.085 & 0.291 & 0.521 & 0.073 & 0.000 & 0.238 & 0.085 & $-$0.154 & 0.000 & $-$0.171 & 0.000 & $-$0.171 & 0.000 \\
\midrule
\multirow{11}{*}{\makecell[l]{\textbf{Llama-Nemotron}\\\textbf{Super-49B v1.5}\\\textbf{(r-on)}\\($N^*=15$)}}
  & 3  & 1.000 & 1.000 & 1.000 & 1.000 & 1.000 & 1.000 & 1.000 & 0.955 & 0.024 & 1.000 & \cellcolor{green!15}$+$0.976 & \cellcolor{green!15}1.000 & \cellcolor{green!15}$+$0.691 & \cellcolor{green!15}0.708 & \cellcolor{green!15}$+$0.701 & \cellcolor{green!15}0.718 \\
  & 5  & 0.977 & 0.978 & 0.977 & 0.956 & 0.977 & 0.977 & 0.955 & 0.955 & 0.101 & 0.955 & \cellcolor{green!15}$+$0.853 & \cellcolor{green!15}0.949 & \cellcolor{green!15}$+$0.551 & \cellcolor{green!15}0.613 & \cellcolor{green!15}$+$0.565 & \cellcolor{green!15}0.628 \\
  & 8  & 0.955 & 0.952 & 0.957 & 0.913 & 0.953 & 0.955 & 0.909 & 0.909 & 0.139 & 0.909 & \cellcolor{green!15}$+$0.770 & \cellcolor{green!15}0.894 & \cellcolor{green!15}$+$0.460 & \cellcolor{green!15}0.534 & \cellcolor{green!15}$+$0.475 & \cellcolor{green!15}0.552 \\
  & 10 & 0.864 & 0.842 & 0.880 & 0.756 & 0.853 & 0.864 & 0.727 & 0.727 & 0.174 & 0.727 & \cellcolor{green!15}$+$0.554 & \cellcolor{green!15}0.670 & \cellcolor{green!15}$+$0.247 & \cellcolor{green!15}0.299 & \cellcolor{green!15}$+$0.268 & \cellcolor{green!15}0.324 \\
  & 15 & 0.727 & 0.625 & 0.786 & 0.542 & 0.674 & 0.727 & 0.455 & 0.318 & 0.127 & 0.455 & \cellcolor{green!15}$+$0.327 & \cellcolor{green!15}0.375 & \cellcolor{green!15}$+$0.079 & \cellcolor{green!15}0.091 & \cellcolor{green!15}$+$0.102 & \cellcolor{green!15}0.117 \\
  & 20 & 0.750 & 0.667 & 0.800 & 0.577 & 0.707 & 0.750 & 0.500 & 0.455 & 0.279 & 0.500 & \cellcolor{green!15}$+$0.221 & \cellcolor{green!15}0.307 & $-$0.040 & 0.000 & $-$0.017 & 0.000 \\
  & 25 & 0.727 & 0.625 & 0.786 & 0.542 & 0.674 & 0.727 & 0.455 & 0.136 & 0.313 & 0.455 & \cellcolor{green!15}$+$0.142 & \cellcolor{green!15}0.206 & $-$0.107 & 0.000 & $-$0.084 & 0.000 \\
  & 30 & 0.614 & 0.485 & 0.691 & 0.262 & 0.560 & 0.614 & 0.273 & 0.045 & 0.217 & 0.314 & \cellcolor{green!15}$+$0.097 & \cellcolor{green!15}0.124 & $-$0.105 & 0.000 & $-$0.085 & 0.000 \\
  & 40 & 0.591 & 0.357 & 0.700 & 0.265 & 0.466 & 0.591 & 0.182 & 0.000 & 0.313 & 0.217 & $-$0.096 & 0.000 & $-$0.253 & 0.000 & $-$0.234 & 0.000 \\
  & 50 & 0.568 & 0.457 & 0.642 & 0.149 & 0.530 & 0.568 & 0.318 & 0.000 & 0.247 & 0.281 & \cellcolor{green!15}$+$0.034 & \cellcolor{green!15}0.045 & $-$0.153 & 0.000 & $-$0.135 & 0.000 \\
  & 60 & 0.400 & 0.222 & 0.512 & -0.024 & 0.340 & 0.491 & 0.077 & 0.000 & 0.163 & 0.115 & $-$0.048 & 0.000 & $-$0.147 & 0.000 & $-$0.136 & 0.000 \\
\bottomrule
\multicolumn{18}{l}{\colorbox{green!15}{\phantom{X}} = estimator clears the estimated Case B side-channel ceiling.}
\end{tabular}%
}
\end{table*}

\paragraph{Case C sensitivity: complexity-debiased side-channel ceiling (2-SAT).}
Case C applies the same feature-level Frisch--Waugh--Lovell residualization as in the commercial 3-SAT domain: each feature in the final Case-B residual battery is first regressed on the raw linear covariates (num\_vars, num\_clauses, num\_literals, file\_size\_bytes) via pooled per-feature ridge, the per-instance residual replaces the raw feature value, and the same max-KS adversary plus permutation calibration are applied; the final ceiling transform remains
$\hatshortcutCeiling=\bigl((1+\hatshortcutDisc)/2\bigr)^2-0.25$.
Thus Case C is stricter than Case A but typically less punitive than raw Case B when the measured leakage mostly tracks formula growth rather than shortcut-only content.

As in the commercial 3-SAT setting, the reported Case C $\hat\shortcutDisc$ on 2-SAT is the 50/50 blend of the permutation-calibrated residualized estimate and the Case A construction-level reference $\delta=0$:
\[
\hat\shortcutDisc^{(\text{C})}
=
\tfrac{1}{2}\max\!\bigl(0,\hat\shortcutDisc_{\text{obs,resid}}-\hat b_{\text{perm,resid}}\bigr).
\]

\Cref{tab:os_pern_caseC} shows that with this blend, the per-$N$ mean residual $\shortcutCeiling$ on 2-SAT sits in the $0.023$--$0.136$ range across the current local open-weight runs, well below the raw Case B values while still high enough to bite at larger $N$ once $\hat{s}$ also drops.
Under the global-pooled CP estimator with this ceiling, Qwen3-14B remains Case C certified through $N=60$ ($+0.076$), Qwen3-32B remains certified through $N=20$ ($+0.042$) and fails at $N=25$ ($-0.046$), QwQ-32B remains certified through $N=20$ ($+0.066$) and fails at $N=25$ ($-0.040$), gpt-oss-20b remains certified through $N=25$ ($+0.078$) and fails at $N=30$ ($-0.055$), Llama-Nemotron-Super-49B v1.5 (reasoning-on) remains certified through all reported $N$ with a borderline positive pooled-CP margin at $N=60$ ($+0.001$), and Nemotron-H-47B-Reasoning-128K (reasoning-on) remains negative throughout on pooled CP.
The Wilson factorized variant gives essentially the same boundaries.
Compared with the raw Case B boundaries on 2-SAT, the Case C certificate is one tested step later for all three open-source models; relative to the Case A construction-level $N^{*}$, Case C tightens or matches but does not collapse, which is the behavior expected from partialing-out generic structural complexity rather than the entire side-channel battery and then taking a 50/50 blend with the construction-level idealization.

\begin{table*}[ht]
\centering
\caption{Open-weight 2-SAT per-$N$ metrics under Case C (complexity-debiased estimated $\shortcutCeiling$).
Values come from the latest local Case-C run with feature-level residualization, permutation calibration, and the default 50/50 blend with the Case-A reference.
$\hat{s}$ is the point-estimate separability and $\mathrm{nRMC}=\max(0,\RMC^{\pm})/(1-\shortcutCeiling)$.}
\label{tab:os_pern_caseC}
\renewcommand{\arraystretch}{1.12}
\resizebox{\textwidth}{!}{%
\begin{tabular}{p{2.2cm}r rrrr rr rr rr rr rr rr}
\toprule
\multirow{2}{*}{Model} & \multirow{2}{*}{$N$}
  & \multirow{2}{*}{Acc}
  & \multicolumn{2}{c}{F1 (per class)}
  & \multirow{2}{*}{MCC}
  & \multirow{2}{*}{G-mean}
  & \multirow{2}{*}{BalAcc}
  & \multirow{2}{*}{ADR} & \multirow{2}{*}{ADR$^{+w}$}
  & \multirow{2}{*}{$\shortcutCeiling$}
  & \multirow{2}{*}{$\hat{s}$}
  & \multicolumn{2}{c}{RMC(no LCB)}
  & \multicolumn{2}{c}{RMC(CP)}
  & \multicolumn{2}{c}{RMC(Wilson)} \\
\cmidrule(lr){4-5}\cmidrule(lr){13-14}\cmidrule(lr){15-16}\cmidrule(lr){17-18}
 & & & F1$_{\mathrm{SAT}}$ & F1$_{\mathrm{UNSAT}}$ & MCC & & & & & & &
  $\RMC$ & $\mathrm{nRMC}$
  & $\RMC^{\pm}$ & $\mathrm{nRMC}$
  & $\RMC^{\pm}$ & $\mathrm{nRMC}$ \\
\midrule
\multirow{11}{*}{\makecell[l]{\textbf{Qwen3-14B}\\\textbf{(think)}\\($N^*=60$)}}
  & 3  & 0.990 & 0.990 & 0.990 & 0.981 & 0.990 & 0.990 & 0.981 & 0.981 & 0.102 & 0.981 & \cellcolor{green!15}$+$0.879 & \cellcolor{green!15}0.978 & \cellcolor{green!15}$+$0.820 & \cellcolor{green!15}0.913 & \cellcolor{green!15}$+$0.819 & \cellcolor{green!15}0.912 \\
  & 5  & 0.974 & 0.973 & 0.975 & 0.949 & 0.974 & 0.974 & 0.948 & 0.942 & 0.115 & 0.948 & \cellcolor{green!15}$+$0.833 & \cellcolor{green!15}0.941 & \cellcolor{green!15}$+$0.764 & \cellcolor{green!15}0.863 & \cellcolor{green!15}$+$0.764 & \cellcolor{green!15}0.863 \\
  & 8  & 0.974 & 0.973 & 0.975 & 0.949 & 0.974 & 0.974 & 0.948 & 0.942 & 0.129 & 0.948 & \cellcolor{green!15}$+$0.819 & \cellcolor{green!15}0.940 & \cellcolor{green!15}$+$0.750 & \cellcolor{green!15}0.861 & \cellcolor{green!15}$+$0.750 & \cellcolor{green!15}0.861 \\
  & 10 & 0.942 & 0.938 & 0.945 & 0.889 & 0.940 & 0.942 & 0.883 & 0.877 & 0.127 & 0.883 & \cellcolor{green!15}$+$0.757 & \cellcolor{green!15}0.866 & \cellcolor{green!15}$+$0.676 & \cellcolor{green!15}0.773 & \cellcolor{green!15}$+$0.676 & \cellcolor{green!15}0.774 \\
  & 15 & 0.844 & 0.815 & 0.865 & 0.724 & 0.830 & 0.844 & 0.688 & 0.656 & 0.127 & 0.688 & \cellcolor{green!15}$+$0.562 & \cellcolor{green!15}0.643 & \cellcolor{green!15}$+$0.468 & \cellcolor{green!15}0.536 & \cellcolor{green!15}$+$0.470 & \cellcolor{green!15}0.538 \\
  & 20 & 0.669 & 0.519 & 0.748 & 0.432 & 0.592 & 0.669 & 0.338 & 0.286 & 0.123 & 0.350 & \cellcolor{green!15}$+$0.228 & \cellcolor{green!15}0.259 & \cellcolor{green!15}$+$0.143 & \cellcolor{green!15}0.163 & \cellcolor{green!15}$+$0.147 & \cellcolor{green!15}0.168 \\
  & 25 & 0.604 & 0.371 & 0.711 & 0.309 & 0.477 & 0.604 & 0.214 & 0.169 & 0.129 & 0.228 & \cellcolor{green!15}$+$0.099 & \cellcolor{green!15}0.114 & \cellcolor{green!15}$+$0.030 & \cellcolor{green!15}0.034 & \cellcolor{green!15}$+$0.034 & \cellcolor{green!15}0.039 \\
  & 30 & 0.516 & 0.280 & 0.636 & 0.043 & 0.399 & 0.516 & 0.162 & 0.058 & 0.123 & 0.159 & \cellcolor{green!15}$+$0.036 & \cellcolor{green!15}0.041 & $-$0.022 & 0.000 & $-$0.018 & 0.000 \\
  & 40 & 0.486 & 0.330 & 0.583 & -0.020 & 0.427 & 0.491 & 0.196 & 0.000 & 0.120 & 0.182 & \cellcolor{green!15}$+$0.062 & \cellcolor{green!15}0.071 & $-$0.002 & 0.000 & \cellcolor{green!15}$+$0.002 & \cellcolor{green!15}0.002 \\
  & 50 & 0.483 & 0.315 & 0.584 & -0.040 & 0.416 & 0.483 & 0.154 & 0.000 & 0.127 & 0.173 & \cellcolor{green!15}$+$0.046 & \cellcolor{green!15}0.053 & $-$0.016 & 0.000 & $-$0.013 & 0.000 \\
  & 60 & 0.556 & 0.482 & 0.612 & 0.117 & 0.537 & 0.556 & 0.294 & 0.000 & 0.129 & 0.289 & \cellcolor{green!15}$+$0.160 & \cellcolor{green!15}0.183 & \cellcolor{green!15}$+$0.076 & \cellcolor{green!15}0.087 & \cellcolor{green!15}$+$0.079 & \cellcolor{green!15}0.091 \\
\midrule
\multirow{11}{*}{\makecell[l]{\textbf{Qwen3-32B}\\\textbf{(think)}\\($N^*=20$)}}
  & 3  & 0.979 & 0.979 & 0.980 & 0.959 & 0.979 & 0.979 & 0.959 & 0.959 & 0.107 & 0.959 & \cellcolor{green!15}$+$0.851 & \cellcolor{green!15}0.954 & \cellcolor{green!15}$+$0.772 & \cellcolor{green!15}0.864 & \cellcolor{green!15}$+$0.772 & \cellcolor{green!15}0.864 \\
  & 5  & 0.955 & 0.952 & 0.957 & 0.913 & 0.953 & 0.955 & 0.909 & 0.909 & 0.114 & 0.909 & \cellcolor{green!15}$+$0.795 & \cellcolor{green!15}0.897 & \cellcolor{green!15}$+$0.703 & \cellcolor{green!15}0.794 & \cellcolor{green!15}$+$0.704 & \cellcolor{green!15}0.795 \\
  & 8  & 0.926 & 0.920 & 0.931 & 0.861 & 0.923 & 0.926 & 0.851 & 0.843 & 0.132 & 0.851 & \cellcolor{green!15}$+$0.720 & \cellcolor{green!15}0.829 & \cellcolor{green!15}$+$0.620 & \cellcolor{green!15}0.714 & \cellcolor{green!15}$+$0.621 & \cellcolor{green!15}0.716 \\
  & 10 & 0.905 & 0.895 & 0.913 & 0.825 & 0.900 & 0.905 & 0.810 & 0.777 & 0.127 & 0.810 & \cellcolor{green!15}$+$0.683 & \cellcolor{green!15}0.782 & \cellcolor{green!15}$+$0.580 & \cellcolor{green!15}0.664 & \cellcolor{green!15}$+$0.582 & \cellcolor{green!15}0.666 \\
  & 15 & 0.748 & 0.663 & 0.799 & 0.574 & 0.704 & 0.748 & 0.496 & 0.430 & 0.122 & 0.496 & \cellcolor{green!15}$+$0.374 & \cellcolor{green!15}0.426 & \cellcolor{green!15}$+$0.270 & \cellcolor{green!15}0.307 & \cellcolor{green!15}$+$0.274 & \cellcolor{green!15}0.312 \\
  & 20 & 0.624 & 0.397 & 0.727 & 0.376 & 0.498 & 0.624 & 0.248 & 0.207 & 0.127 & 0.248 & \cellcolor{green!15}$+$0.121 & \cellcolor{green!15}0.139 & \cellcolor{green!15}$+$0.042 & \cellcolor{green!15}0.048 & \cellcolor{green!15}$+$0.047 & \cellcolor{green!15}0.054 \\
  & 25 & 0.568 & 0.246 & 0.698 & 0.254 & 0.375 & 0.567 & 0.142 & 0.100 & 0.127 & 0.140 & \cellcolor{green!15}$+$0.014 & \cellcolor{green!15}0.016 & $-$0.046 & 0.000 & $-$0.040 & 0.000 \\
  & 30 & 0.574 & 0.270 & 0.700 & 0.270 & 0.395 & 0.574 & 0.157 & 0.083 & 0.127 & 0.156 & \cellcolor{green!15}$+$0.029 & \cellcolor{green!15}0.033 & $-$0.034 & 0.000 & $-$0.028 & 0.000 \\
  & 40 & 0.525 & 0.173 & 0.667 & 0.094 & 0.307 & 0.525 & 0.099 & 0.000 & 0.124 & 0.094 & $-$0.030 & 0.000 & $-$0.077 & 0.000 & $-$0.072 & 0.000 \\
  & 50 & 0.471 & 0.274 & 0.584 & -0.074 & 0.385 & 0.469 & 0.151 & 0.000 & 0.125 & 0.148 & \cellcolor{green!15}$+$0.023 & \cellcolor{green!15}0.027 & $-$0.038 & 0.000 & $-$0.034 & 0.000 \\
  & 60 & 0.485 & 0.267 & 0.603 & -0.014 & 0.383 & 0.494 & 0.167 & 0.000 & 0.122 & 0.147 & \cellcolor{green!15}$+$0.025 & \cellcolor{green!15}0.029 & $-$0.037 & 0.000 & $-$0.032 & 0.000 \\
\midrule
\multirow{11}{*}{\makecell[l]{\textbf{QwQ-32B}\\\textbf{(think)}\\($N^*=20$)}}
  & 3  & 1.000 & 1.000 & 1.000 & 1.000 & 1.000 & 1.000 & 1.000 & 1.000 & 0.097 & 1.000 & \cellcolor{green!15}$+$0.903 & \cellcolor{green!15}1.000 & \cellcolor{green!15}$+$0.831 & \cellcolor{green!15}0.920 & \cellcolor{green!15}$+$0.829 & \cellcolor{green!15}0.919 \\
  & 5  & 0.990 & 0.990 & 0.990 & 0.980 & 0.990 & 0.990 & 0.980 & 0.980 & 0.112 & 0.980 & \cellcolor{green!15}$+$0.868 & \cellcolor{green!15}0.977 & \cellcolor{green!15}$+$0.783 & \cellcolor{green!15}0.882 & \cellcolor{green!15}$+$0.783 & \cellcolor{green!15}0.881 \\
  & 8  & 0.980 & 0.980 & 0.980 & 0.960 & 0.980 & 0.980 & 0.960 & 0.960 & 0.130 & 0.960 & \cellcolor{green!15}$+$0.830 & \cellcolor{green!15}0.954 & \cellcolor{green!15}$+$0.734 & \cellcolor{green!15}0.843 & \cellcolor{green!15}$+$0.735 & \cellcolor{green!15}0.844 \\
  & 10 & 0.949 & 0.947 & 0.952 & 0.904 & 0.948 & 0.949 & 0.899 & 0.869 & 0.127 & 0.899 & \cellcolor{green!15}$+$0.772 & \cellcolor{green!15}0.884 & \cellcolor{green!15}$+$0.665 & \cellcolor{green!15}0.762 & \cellcolor{green!15}$+$0.667 & \cellcolor{green!15}0.763 \\
  & 15 & 0.803 & 0.755 & 0.835 & 0.659 & 0.778 & 0.803 & 0.606 & 0.556 & 0.127 & 0.606 & \cellcolor{green!15}$+$0.480 & \cellcolor{green!15}0.549 & \cellcolor{green!15}$+$0.358 & \cellcolor{green!15}0.410 & \cellcolor{green!15}$+$0.362 & \cellcolor{green!15}0.415 \\
  & 20 & 0.641 & 0.441 & 0.736 & 0.406 & 0.532 & 0.641 & 0.283 & 0.242 & 0.124 & 0.283 & \cellcolor{green!15}$+$0.159 & \cellcolor{green!15}0.182 & \cellcolor{green!15}$+$0.066 & \cellcolor{green!15}0.075 & \cellcolor{green!15}$+$0.072 & \cellcolor{green!15}0.083 \\
  & 25 & 0.576 & 0.276 & 0.700 & 0.270 & 0.400 & 0.576 & 0.162 & 0.111 & 0.130 & 0.160 & \cellcolor{green!15}$+$0.030 & \cellcolor{green!15}0.035 & $-$0.040 & 0.000 & $-$0.033 & 0.000 \\
  & 30 & 0.515 & 0.059 & 0.673 & 0.124 & 0.174 & 0.515 & 0.030 & 0.000 & 0.127 & 0.030 & $-$0.096 & 0.000 & $-$0.120 & 0.000 & $-$0.117 & 0.000 \\
  & 40 & 0.487 & 0.093 & 0.642 & 0.043 & 0.221 & 0.509 & 0.033 & 0.011 & 0.118 & 0.049 & $-$0.069 & 0.000 & $-$0.103 & 0.000 & $-$0.098 & 0.000 \\
  & 50 & 0.506 & 0.044 & 0.667 & 0.044 & 0.150 & 0.506 & 0.023 & 0.000 & 0.127 & 0.022 & $-$0.104 & 0.000 & $-$0.124 & 0.000 & $-$0.121 & 0.000 \\
  & 60 & 0.500 & -- & 0.667 & -- & 0.000 & 0.500 & 0.000 & 0.000 & 0.127 & 0.000 & $-$0.127 & 0.000 & $-$0.127 & 0.000 & $-$0.127 & 0.000 \\
\midrule
\multirow{11}{*}{\makecell[l]{\textbf{gpt-oss-20b}\\\textbf{(reason)}\\($N^*=25$)}}
  & 3  & 1.000 & 1.000 & 1.000 & 1.000 & 1.000 & 1.000 & 1.000 & 1.000 & 0.091 & 1.000 & \cellcolor{green!15}$+$0.909 & \cellcolor{green!15}1.000 & \cellcolor{green!15}$+$0.709 & \cellcolor{green!15}0.780 & \cellcolor{green!15}$+$0.711 & \cellcolor{green!15}0.783 \\
  & 5  & 0.970 & 0.970 & 0.970 & 0.939 & 0.970 & 0.970 & 0.939 & 0.939 & 0.066 & 0.940 & \cellcolor{green!15}$+$0.874 & \cellcolor{green!15}0.936 & \cellcolor{green!15}$+$0.643 & \cellcolor{green!15}0.689 & \cellcolor{green!15}$+$0.651 & \cellcolor{green!15}0.697 \\
  & 8  & 0.924 & 0.918 & 0.930 & 0.858 & 0.921 & 0.924 & 0.848 & 0.818 & 0.136 & 0.848 & \cellcolor{green!15}$+$0.712 & \cellcolor{green!15}0.825 & \cellcolor{green!15}$+$0.473 & \cellcolor{green!15}0.547 & \cellcolor{green!15}$+$0.483 & \cellcolor{green!15}0.559 \\
  & 10 & 0.894 & 0.896 & 0.892 & 0.788 & 0.894 & 0.894 & 0.818 & 0.758 & 0.127 & 0.799 & \cellcolor{green!15}$+$0.672 & \cellcolor{green!15}0.770 & \cellcolor{green!15}$+$0.417 & \cellcolor{green!15}0.477 & \cellcolor{green!15}$+$0.429 & \cellcolor{green!15}0.491 \\
  & 15 & 0.848 & 0.833 & 0.861 & 0.709 & 0.844 & 0.848 & 0.697 & 0.636 & 0.127 & 0.712 & \cellcolor{green!15}$+$0.585 & \cellcolor{green!15}0.670 & \cellcolor{green!15}$+$0.334 & \cellcolor{green!15}0.382 & \cellcolor{green!15}$+$0.348 & \cellcolor{green!15}0.398 \\
  & 20 & 0.742 & 0.679 & 0.785 & 0.528 & 0.716 & 0.742 & 0.515 & 0.424 & 0.127 & 0.512 & \cellcolor{green!15}$+$0.386 & \cellcolor{green!15}0.442 & \cellcolor{green!15}$+$0.163 & \cellcolor{green!15}0.187 & \cellcolor{green!15}$+$0.179 & \cellcolor{green!15}0.205 \\
  & 25 & 0.652 & 0.596 & 0.693 & 0.315 & 0.637 & 0.652 & 0.364 & 0.242 & 0.127 & 0.406 & \cellcolor{green!15}$+$0.279 & \cellcolor{green!15}0.320 & \cellcolor{green!15}$+$0.078 & \cellcolor{green!15}0.090 & \cellcolor{green!15}$+$0.093 & \cellcolor{green!15}0.106 \\
  & 30 & 0.576 & 0.333 & 0.689 & 0.221 & 0.446 & 0.576 & 0.182 & 0.091 & 0.127 & 0.199 & \cellcolor{green!15}$+$0.073 & \cellcolor{green!15}0.083 & $-$0.055 & 0.000 & $-$0.041 & 0.000 \\
  & 40 & 0.492 & 0.250 & 0.615 & 0.115 & 0.374 & 0.537 & 0.154 & 0.000 & 0.107 & 0.140 & \cellcolor{green!15}$+$0.033 & \cellcolor{green!15}0.037 & $-$0.068 & 0.000 & $-$0.056 & 0.000 \\
  & 50 & 0.500 & 0.083 & 0.656 & 0.000 & 0.208 & 0.500 & 0.045 & 0.000 & 0.127 & 0.043 & $-$0.083 & 0.000 & $-$0.126 & 0.000 & $-$0.120 & 0.000 \\
  & 60 & 0.523 & 0.160 & 0.667 & 0.090 & 0.295 & 0.523 & 0.091 & 0.000 & 0.100 & 0.087 & $-$0.013 & 0.000 & $-$0.091 & 0.000 & $-$0.080 & 0.000 \\
\midrule
\multirow{11}{*}{\makecell[l]{\textbf{Nemotron-H-47B}\\\textbf{Reasoning-128K}\\\textbf{(r-on)}\\($N^*=0$)}}
  & 3  & 0.525 & 0.663 & 0.200 & 0.087 & 0.332 & 0.525 & 0.068 & 0.038 & 0.104 & 0.110 & \cellcolor{green!15}$+$0.006 & \cellcolor{green!15}0.007 & $-$0.014 & 0.000 & $-$0.013 & 0.000 \\
  & 5  & 0.513 & 0.670 & 0.071 & 0.082 & 0.191 & 0.513 & 0.030 & 0.002 & 0.121 & 0.037 & $-$0.085 & 0.000 & $-$0.096 & 0.000 & $-$0.095 & 0.000 \\
  & 8  & 0.500 & 0.666 & 0.002 & 0.000 & 0.032 & 0.500 & 0.001 & 0.000 & 0.126 & 0.001 & $-$0.125 & 0.000 & $-$0.126 & 0.000 & $-$0.126 & 0.000 \\
  & 10 & 0.497 & 0.664 & 0.002 & -0.049 & 0.032 & 0.497 & 0.001 & 0.000 & 0.126 & 0.001 & $-$0.125 & 0.000 & $-$0.126 & 0.000 & $-$0.126 & 0.000 \\
  & 15 & 0.512 & 0.669 & 0.072 & 0.072 & 0.194 & 0.512 & 0.026 & 0.000 & 0.122 & 0.037 & $-$0.085 & 0.000 & $-$0.096 & 0.000 & $-$0.095 & 0.000 \\
  & 20 & 0.511 & 0.666 & 0.091 & 0.053 & 0.219 & 0.510 & 0.034 & 0.000 & 0.125 & 0.048 & $-$0.077 & 0.000 & $-$0.091 & 0.000 & $-$0.090 & 0.000 \\
  & 25 & 0.520 & 0.663 & 0.167 & 0.074 & 0.301 & 0.520 & 0.068 & 0.000 & 0.125 & 0.091 & $-$0.034 & 0.000 & $-$0.053 & 0.000 & $-$0.052 & 0.000 \\
  & 30 & 0.499 & 0.660 & 0.047 & -0.007 & 0.155 & 0.499 & 0.020 & 0.000 & 0.128 & 0.024 & $-$0.104 & 0.000 & $-$0.113 & 0.000 & $-$0.112 & 0.000 \\
  & 40 & 0.518 & 0.670 & 0.105 & 0.094 & 0.236 & 0.518 & 0.052 & 0.000 & 0.126 & 0.056 & $-$0.070 & 0.000 & $-$0.084 & 0.000 & $-$0.083 & 0.000 \\
  & 50 & 0.509 & 0.661 & 0.110 & 0.038 & 0.241 & 0.509 & 0.048 & 0.000 & 0.127 & 0.058 & $-$0.068 & 0.000 & $-$0.083 & 0.000 & $-$0.082 & 0.000 \\
  & 60 & 0.521 & 0.666 & 0.156 & 0.085 & 0.291 & 0.521 & 0.073 & 0.000 & 0.127 & 0.085 & $-$0.042 & 0.000 & $-$0.060 & 0.000 & $-$0.059 & 0.000 \\
\midrule
\multirow{11}{*}{\makecell[l]{\textbf{Llama-Nemotron}\\\textbf{Super-49B v1.5}\\\textbf{(r-on)}\\($N^*=25$)}}
  & 3  & 1.000 & 1.000 & 1.000 & 1.000 & 1.000 & 1.000 & 1.000 & 0.955 & 0.087 & 1.000 & \cellcolor{green!15}$+$0.913 & \cellcolor{green!15}1.000 & \cellcolor{green!15}$+$0.628 & \cellcolor{green!15}0.688 & \cellcolor{green!15}$+$0.638 & \cellcolor{green!15}0.699 \\
  & 5  & 0.977 & 0.978 & 0.977 & 0.956 & 0.977 & 0.977 & 0.955 & 0.955 & 0.100 & 0.955 & \cellcolor{green!15}$+$0.855 & \cellcolor{green!15}0.950 & \cellcolor{green!15}$+$0.553 & \cellcolor{green!15}0.614 & \cellcolor{green!15}$+$0.566 & \cellcolor{green!15}0.629 \\
  & 8  & 0.955 & 0.952 & 0.957 & 0.913 & 0.953 & 0.955 & 0.909 & 0.909 & 0.141 & 0.909 & \cellcolor{green!15}$+$0.768 & \cellcolor{green!15}0.894 & \cellcolor{green!15}$+$0.458 & \cellcolor{green!15}0.533 & \cellcolor{green!15}$+$0.474 & \cellcolor{green!15}0.551 \\
  & 10 & 0.864 & 0.842 & 0.880 & 0.756 & 0.853 & 0.864 & 0.727 & 0.727 & 0.127 & 0.727 & \cellcolor{green!15}$+$0.601 & \cellcolor{green!15}0.688 & \cellcolor{green!15}$+$0.294 & \cellcolor{green!15}0.337 & \cellcolor{green!15}$+$0.315 & \cellcolor{green!15}0.360 \\
  & 15 & 0.727 & 0.625 & 0.786 & 0.542 & 0.674 & 0.727 & 0.455 & 0.318 & 0.127 & 0.455 & \cellcolor{green!15}$+$0.328 & \cellcolor{green!15}0.376 & \cellcolor{green!15}$+$0.080 & \cellcolor{green!15}0.091 & \cellcolor{green!15}$+$0.103 & \cellcolor{green!15}0.118 \\
  & 20 & 0.750 & 0.667 & 0.800 & 0.577 & 0.707 & 0.750 & 0.500 & 0.455 & 0.127 & 0.500 & \cellcolor{green!15}$+$0.373 & \cellcolor{green!15}0.428 & \cellcolor{green!15}$+$0.112 & \cellcolor{green!15}0.128 & \cellcolor{green!15}$+$0.135 & \cellcolor{green!15}0.155 \\
  & 25 & 0.727 & 0.625 & 0.786 & 0.542 & 0.674 & 0.727 & 0.455 & 0.136 & 0.127 & 0.455 & \cellcolor{green!15}$+$0.328 & \cellcolor{green!15}0.376 & \cellcolor{green!15}$+$0.080 & \cellcolor{green!15}0.091 & \cellcolor{green!15}$+$0.103 & \cellcolor{green!15}0.118 \\
  & 30 & 0.614 & 0.485 & 0.691 & 0.262 & 0.560 & 0.614 & 0.273 & 0.045 & 0.127 & 0.314 & \cellcolor{green!15}$+$0.188 & \cellcolor{green!15}0.215 & $-$0.015 & 0.000 & \cellcolor{green!15}$+$0.005 & \cellcolor{green!15}0.006 \\
  & 40 & 0.591 & 0.357 & 0.700 & 0.265 & 0.466 & 0.591 & 0.182 & 0.000 & 0.127 & 0.217 & \cellcolor{green!15}$+$0.090 & \cellcolor{green!15}0.103 & $-$0.066 & 0.000 & $-$0.047 & 0.000 \\
  & 50 & 0.568 & 0.457 & 0.642 & 0.149 & 0.530 & 0.568 & 0.318 & 0.000 & 0.127 & 0.281 & \cellcolor{green!15}$+$0.154 & \cellcolor{green!15}0.177 & $-$0.033 & 0.000 & $-$0.015 & 0.000 \\
  & 60 & 0.400 & 0.222 & 0.512 & -0.024 & 0.340 & 0.491 & 0.077 & 0.000 & 0.023 & 0.115 & \cellcolor{green!15}$+$0.092 & \cellcolor{green!15}0.094 & $-$0.007 & 0.000 & \cellcolor{green!15}$+$0.004 & \cellcolor{green!15}0.004 \\
\bottomrule
\multicolumn{18}{l}{\colorbox{green!15}{\phantom{X}} = estimator clears the estimated Case C side-channel ceiling.}
\end{tabular}%
}
\end{table*}

\paragraph{Case D diagnostic: separability-only lower bound (2-SAT).}
Case D has the same meaning on 2-SAT as on commercial 3-SAT: it removes shortcut subtraction and witness gating and reports the factorized separability statistic itself.  The resulting tables therefore isolate balanced SAT/UNSAT discrimination without asking whether that discrimination exceeds any explicit shortcut account.

\Cref{tab:os_pern_caseD} shows a more differentiated picture once the added OpenAI and NVIDIA baselines are included.  The three certified Qwen-family thinking models remain Case D positive throughout the tested 2-SAT range: Qwen3-14B declines from pooled CP $0.922$ at $N=3$ to $0.110$ at $N=50$; Qwen3-32B declines from $0.879$ to $0.087$ over the same range; and QwQ-32B declines from $0.928$ to $0.003$.
By contrast, gpt-oss-20b remains Case-D positive through $N=20$ and turns negative at $N=25$ under pooled CP, Nemotron-H-47B-Reasoning-128K (reasoning-on) remains negative throughout on pooled CP, and Llama-Nemotron-Super-49B v1.5 (reasoning-on) remains positive through $N=10$ under pooled CP before crossing negative at $N=15$.
This ranking differs from the Case B and Case C boundary view in a predictable way: Case D rewards balanced discrimination alone, whereas Cases B and C additionally ask whether that discrimination survives an explicit shortcut account.  The widening gap at larger $N$ is therefore diagnostic of shortcut-adjusted fragility rather than of total collapse of separability.

\begin{table*}[ht]
\centering
\caption{Open-weight 2-SAT per-$N$ diagnostics under Case D (separability only).
Case D removes shortcut subtraction and witness gating, so the table reports the factorized separability point estimate $\hat{s}$ together with the pooled CP and Wilson lower bounds.}
\label{tab:os_pern_caseD}
\renewcommand{\arraystretch}{1.12}
\resizebox{\textwidth}{!}{%
\begin{tabular}{p{2.2cm}r rrrr rr rr r rr}
\toprule
\multirow{2}{*}{Model} & \multirow{2}{*}{$N$}
  & \multirow{2}{*}{Acc}
  & \multicolumn{2}{c}{F1 (per class)}
  & \multirow{2}{*}{MCC}
  & \multirow{2}{*}{G-mean}
  & \multirow{2}{*}{BalAcc}
  & \multirow{2}{*}{ADR} & \multirow{2}{*}{ADR$^{+w}$}
  & \multirow{2}{*}{$\hat{s}$}
  & \multirow{2}{*}{RMC(CP)}
  & \multirow{2}{*}{RMC(Wilson)} \\
\cmidrule(lr){4-5}
 & & & F1$_{\mathrm{SAT}}$ & F1$_{\mathrm{UNSAT}}$ & MCC & & & & & & & \\
\midrule
\multirow{11}{*}{\makecell[l]{\textbf{Qwen3-14B}\\\textbf{(think)}\\($N^*=20$)}}
  & 3  & 0.990 & 0.990 & 0.990 & 0.981 & 0.990 & 0.990 & 0.981 & 0.981 & 0.981 & \cellcolor{green!15}$+$0.672 & \cellcolor{green!15}$+$0.671 \\
  & 5  & 0.974 & 0.973 & 0.975 & 0.949 & 0.974 & 0.974 & 0.948 & 0.942 & 0.948 & \cellcolor{green!15}$+$0.629 & \cellcolor{green!15}$+$0.629 \\
  & 8  & 0.974 & 0.973 & 0.975 & 0.949 & 0.974 & 0.974 & 0.948 & 0.942 & 0.948 & \cellcolor{green!15}$+$0.629 & \cellcolor{green!15}$+$0.629 \\
  & 10 & 0.942 & 0.938 & 0.945 & 0.889 & 0.940 & 0.942 & 0.883 & 0.877 & 0.883 & \cellcolor{green!15}$+$0.552 & \cellcolor{green!15}$+$0.553 \\
  & 15 & 0.844 & 0.815 & 0.865 & 0.724 & 0.830 & 0.844 & 0.688 & 0.656 & 0.688 & \cellcolor{green!15}$+$0.344 & \cellcolor{green!15}$+$0.346 \\
  & 20 & 0.669 & 0.519 & 0.748 & 0.432 & 0.592 & 0.669 & 0.338 & 0.286 & 0.350 & \cellcolor{green!15}$+$0.016 & \cellcolor{green!15}$+$0.020 \\
  & 25 & 0.604 & 0.371 & 0.711 & 0.309 & 0.477 & 0.604 & 0.214 & 0.169 & 0.228 & $-$0.092 & $-$0.087 \\
  & 30 & 0.516 & 0.280 & 0.636 & 0.043 & 0.399 & 0.516 & 0.162 & 0.058 & 0.159 & $-$0.149 & $-$0.145 \\
  & 40 & 0.486 & 0.330 & 0.583 & -0.020 & 0.427 & 0.491 & 0.196 & 0.000 & 0.182 & $-$0.131 & $-$0.128 \\
  & 50 & 0.483 & 0.315 & 0.584 & -0.040 & 0.416 & 0.483 & 0.154 & 0.000 & 0.173 & $-$0.140 & $-$0.136 \\
  & 60 & 0.556 & 0.482 & 0.612 & 0.117 & 0.537 & 0.556 & 0.294 & 0.000 & 0.289 & $-$0.046 & $-$0.042 \\
\midrule
\multirow{11}{*}{\makecell[l]{\textbf{Qwen3-32B}\\\textbf{(think)}\\($N^*=15$)}}
  & 3  & 0.979 & 0.979 & 0.980 & 0.959 & 0.979 & 0.979 & 0.959 & 0.959 & 0.959 & \cellcolor{green!15}$+$0.629 & \cellcolor{green!15}$+$0.629 \\
  & 5  & 0.955 & 0.952 & 0.957 & 0.913 & 0.953 & 0.955 & 0.909 & 0.909 & 0.909 & \cellcolor{green!15}$+$0.568 & \cellcolor{green!15}$+$0.569 \\
  & 8  & 0.926 & 0.920 & 0.931 & 0.861 & 0.923 & 0.926 & 0.851 & 0.843 & 0.851 & \cellcolor{green!15}$+$0.502 & \cellcolor{green!15}$+$0.503 \\
  & 10 & 0.905 & 0.895 & 0.913 & 0.825 & 0.900 & 0.905 & 0.810 & 0.777 & 0.810 & \cellcolor{green!15}$+$0.457 & \cellcolor{green!15}$+$0.458 \\
  & 15 & 0.748 & 0.663 & 0.799 & 0.574 & 0.704 & 0.748 & 0.496 & 0.430 & 0.496 & \cellcolor{green!15}$+$0.142 & \cellcolor{green!15}$+$0.146 \\
  & 20 & 0.624 & 0.397 & 0.727 & 0.376 & 0.498 & 0.624 & 0.248 & 0.207 & 0.248 & $-$0.081 & $-$0.076 \\
  & 25 & 0.568 & 0.246 & 0.698 & 0.254 & 0.375 & 0.567 & 0.142 & 0.100 & 0.140 & $-$0.169 & $-$0.164 \\
  & 30 & 0.574 & 0.270 & 0.700 & 0.270 & 0.395 & 0.574 & 0.157 & 0.083 & 0.156 & $-$0.157 & $-$0.152 \\
  & 40 & 0.525 & 0.173 & 0.667 & 0.094 & 0.307 & 0.525 & 0.099 & 0.000 & 0.094 & $-$0.203 & $-$0.198 \\
  & 50 & 0.471 & 0.274 & 0.584 & -0.074 & 0.385 & 0.469 & 0.151 & 0.000 & 0.148 & $-$0.163 & $-$0.159 \\
  & 60 & 0.485 & 0.267 & 0.603 & -0.014 & 0.383 & 0.494 & 0.167 & 0.000 & 0.147 & $-$0.165 & $-$0.161 \\
\midrule
\multirow{11}{*}{\makecell[l]{\textbf{QwQ-32B}\\\textbf{(think)}\\($N^*=15$)}}
  & 3  & 1.000 & 1.000 & 1.000 & 1.000 & 1.000 & 1.000 & 1.000 & 1.000 & 1.000 & \cellcolor{green!15}$+$0.678 & \cellcolor{green!15}$+$0.677 \\
  & 5  & 0.990 & 0.990 & 0.990 & 0.980 & 0.990 & 0.990 & 0.980 & 0.980 & 0.980 & \cellcolor{green!15}$+$0.645 & \cellcolor{green!15}$+$0.645 \\
  & 8  & 0.980 & 0.980 & 0.980 & 0.960 & 0.980 & 0.980 & 0.960 & 0.960 & 0.960 & \cellcolor{green!15}$+$0.614 & \cellcolor{green!15}$+$0.614 \\
  & 10 & 0.949 & 0.947 & 0.952 & 0.904 & 0.948 & 0.949 & 0.899 & 0.869 & 0.899 & \cellcolor{green!15}$+$0.542 & \cellcolor{green!15}$+$0.543 \\
  & 15 & 0.803 & 0.755 & 0.835 & 0.659 & 0.778 & 0.803 & 0.606 & 0.556 & 0.606 & \cellcolor{green!15}$+$0.234 & \cellcolor{green!15}$+$0.239 \\
  & 20 & 0.641 & 0.441 & 0.736 & 0.406 & 0.532 & 0.641 & 0.283 & 0.242 & 0.283 & $-$0.060 & $-$0.054 \\
  & 25 & 0.576 & 0.276 & 0.700 & 0.270 & 0.400 & 0.576 & 0.162 & 0.111 & 0.160 & $-$0.160 & $-$0.154 \\
  & 30 & 0.515 & 0.059 & 0.673 & 0.124 & 0.174 & 0.515 & 0.030 & 0.000 & 0.030 & $-$0.244 & $-$0.240 \\
  & 40 & 0.487 & 0.093 & 0.642 & 0.043 & 0.221 & 0.509 & 0.033 & 0.011 & 0.049 & $-$0.235 & $-$0.230 \\
  & 50 & 0.506 & 0.044 & 0.667 & 0.044 & 0.150 & 0.506 & 0.023 & 0.000 & 0.022 & $-$0.247 & $-$0.244 \\
  & 60 & 0.500 & -- & 0.667 & -- & 0.000 & 0.500 & 0.000 & 0.000 & 0.000 & $-$0.250 & $-$0.250 \\
\midrule
\multirow{11}{*}{\makecell[l]{\textbf{gpt-oss-20b}\\\textbf{(reason)}\\($N^*=20$)}}
  & 3  & 1.000 & 1.000 & 1.000 & 1.000 & 1.000 & 1.000 & 1.000 & 1.000 & 1.000 & \cellcolor{green!15}$+$0.550 & \cellcolor{green!15}$+$0.552 \\
  & 5  & 0.970 & 0.970 & 0.970 & 0.939 & 0.970 & 0.970 & 0.939 & 0.939 & 0.940 & \cellcolor{green!15}$+$0.460 & \cellcolor{green!15}$+$0.467 \\
  & 8  & 0.924 & 0.918 & 0.930 & 0.858 & 0.921 & 0.924 & 0.848 & 0.818 & 0.848 & \cellcolor{green!15}$+$0.359 & \cellcolor{green!15}$+$0.369 \\
  & 10 & 0.894 & 0.896 & 0.892 & 0.788 & 0.894 & 0.894 & 0.818 & 0.758 & 0.799 & \cellcolor{green!15}$+$0.293 & \cellcolor{green!15}$+$0.305 \\
  & 15 & 0.848 & 0.833 & 0.861 & 0.709 & 0.844 & 0.848 & 0.697 & 0.636 & 0.712 & \cellcolor{green!15}$+$0.211 & \cellcolor{green!15}$+$0.224 \\
  & 20 & 0.742 & 0.679 & 0.785 & 0.528 & 0.716 & 0.742 & 0.515 & 0.424 & 0.512 & \cellcolor{green!15}$+$0.040 & \cellcolor{green!15}$+$0.055 \\
  & 25 & 0.652 & 0.596 & 0.693 & 0.315 & 0.637 & 0.652 & 0.364 & 0.242 & 0.406 & $-$0.045 & $-$0.031 \\
  & 30 & 0.576 & 0.333 & 0.689 & 0.221 & 0.446 & 0.576 & 0.182 & 0.091 & 0.199 & $-$0.178 & $-$0.164 \\
  & 40 & 0.492 & 0.250 & 0.615 & 0.115 & 0.374 & 0.537 & 0.154 & 0.000 & 0.140 & $-$0.212 & $-$0.200 \\
  & 50 & 0.500 & 0.083 & 0.656 & 0.000 & 0.208 & 0.500 & 0.045 & 0.000 & 0.043 & $-$0.249 & $-$0.244 \\
  & 60 & 0.523 & 0.160 & 0.667 & 0.090 & 0.295 & 0.523 & 0.091 & 0.000 & 0.087 & $-$0.241 & $-$0.230 \\
\midrule
\multirow{11}{*}{\makecell[l]{\textbf{Nemotron-H-47B}\\\textbf{Reasoning-128K}\\\textbf{(r-on)}\\($N^*=0$)}}
  & 3  & 0.525 & 0.663 & 0.200 & 0.087 & 0.332 & 0.525 & 0.068 & 0.038 & 0.110 & $-$0.160 & $-$0.159 \\
  & 5  & 0.513 & 0.670 & 0.071 & 0.082 & 0.191 & 0.513 & 0.030 & 0.002 & 0.037 & $-$0.225 & $-$0.224 \\
  & 8  & 0.500 & 0.666 & 0.002 & 0.000 & 0.032 & 0.500 & 0.001 & 0.000 & 0.001 & $-$0.250 & $-$0.250 \\
  & 10 & 0.497 & 0.664 & 0.002 & -0.049 & 0.032 & 0.497 & 0.001 & 0.000 & 0.001 & $-$0.250 & $-$0.250 \\
  & 15 & 0.512 & 0.669 & 0.072 & 0.072 & 0.194 & 0.512 & 0.026 & 0.000 & 0.037 & $-$0.224 & $-$0.223 \\
  & 20 & 0.511 & 0.666 & 0.091 & 0.053 & 0.219 & 0.510 & 0.034 & 0.000 & 0.048 & $-$0.215 & $-$0.214 \\
  & 25 & 0.520 & 0.663 & 0.167 & 0.074 & 0.301 & 0.520 & 0.068 & 0.000 & 0.091 & $-$0.178 & $-$0.177 \\
  & 30 & 0.499 & 0.660 & 0.047 & -0.007 & 0.155 & 0.499 & 0.020 & 0.000 & 0.024 & $-$0.235 & $-$0.234 \\
  & 40 & 0.518 & 0.670 & 0.105 & 0.094 & 0.236 & 0.518 & 0.052 & 0.000 & 0.056 & $-$0.209 & $-$0.208 \\
  & 50 & 0.509 & 0.661 & 0.110 & 0.038 & 0.241 & 0.509 & 0.048 & 0.000 & 0.058 & $-$0.206 & $-$0.205 \\
  & 60 & 0.521 & 0.666 & 0.156 & 0.085 & 0.291 & 0.521 & 0.073 & 0.000 & 0.085 & $-$0.183 & $-$0.182 \\
\midrule
\multirow{11}{*}{\makecell[l]{\textbf{Llama-Nemotron}\\\textbf{Super-49B v1.5}\\\textbf{(r-on)}\\($N^*=10$)}}
  & 3  & 1.000 & 1.000 & 1.000 & 1.000 & 1.000 & 1.000 & 1.000 & 0.955 & 1.000 & \cellcolor{green!15}$+$0.465 & \cellcolor{green!15}$+$0.475 \\
  & 5  & 0.977 & 0.978 & 0.977 & 0.956 & 0.977 & 0.977 & 0.955 & 0.955 & 0.955 & \cellcolor{green!15}$+$0.402 & \cellcolor{green!15}$+$0.416 \\
  & 8  & 0.955 & 0.952 & 0.957 & 0.913 & 0.953 & 0.955 & 0.909 & 0.909 & 0.909 & \cellcolor{green!15}$+$0.349 & \cellcolor{green!15}$+$0.365 \\
  & 10 & 0.864 & 0.842 & 0.880 & 0.756 & 0.853 & 0.864 & 0.727 & 0.727 & 0.727 & \cellcolor{green!15}$+$0.171 & \cellcolor{green!15}$+$0.191 \\
  & 15 & 0.727 & 0.625 & 0.786 & 0.542 & 0.674 & 0.727 & 0.455 & 0.318 & 0.455 & $-$0.044 & $-$0.021 \\
  & 20 & 0.750 & 0.667 & 0.800 & 0.577 & 0.707 & 0.750 & 0.500 & 0.455 & 0.500 & $-$0.011 & \cellcolor{green!15}$+$0.012 \\
  & 25 & 0.727 & 0.625 & 0.786 & 0.542 & 0.674 & 0.727 & 0.455 & 0.136 & 0.455 & $-$0.044 & $-$0.021 \\
  & 30 & 0.614 & 0.485 & 0.691 & 0.262 & 0.560 & 0.614 & 0.273 & 0.045 & 0.314 & $-$0.138 & $-$0.118 \\
  & 40 & 0.591 & 0.357 & 0.700 & 0.265 & 0.466 & 0.591 & 0.182 & 0.000 & 0.217 & $-$0.190 & $-$0.171 \\
  & 50 & 0.568 & 0.457 & 0.642 & 0.149 & 0.530 & 0.568 & 0.318 & 0.000 & 0.281 & $-$0.156 & $-$0.138 \\
  & 60 & 0.400 & 0.222 & 0.512 & -0.024 & 0.340 & 0.491 & 0.077 & 0.000 & 0.115 & $-$0.234 & $-$0.223 \\
\bottomrule
\multicolumn{13}{l}{\colorbox{green!15}{\phantom{X}} = positive separability lower bound under that estimator.}
\end{tabular}%
}
\end{table*}

\subsection{Open-weight 3-Colouring supplementary audit}
\label{sec:os_graph_coloring}

We also reran the six open-weight models that appear in the supplementary manuscript on the latest local grouped 3-Colouring benchmark and summarize the corresponding Case A--D results here.
The current graph-coloring grid reports $N\in\{8,10,15,20,30,40,50,60\}$.
As in the SAT-domain tables above, we report only successful local outputs with parsed results; failed requests and missing outputs are excluded from every denominator.

\paragraph{Compact notation for the open-weight 3-Colouring per-$N$ tables.}
Here F1$_C$ denotes the F1 score for the \emph{3-colourable} class and F1$_N$ denotes the F1 score for the \emph{non-3-colourable} class.
ADR$^{+a}$ is the assignment-checked constructive ADR based on returned colouring witnesses.
$\hat{s}$ is the factorized separability point estimate, LCB$_{CP}$ and LCB$_W$ are the pooled Clopper--Pearson and pooled Wilson separability lower bounds, $w_{LCB}$ is the witness lower bound, and RMC$^{\pm}$ is the final signed certificate after the case-specific ceiling and witness gate.
For Case B and Case C, $\hat{\delta}$ denotes the estimated residual shortcut distinguishability and $\shortcutCeiling=\bigl((1+\hat{\delta})/2\bigr)^2-0.25$.
On the current graph-coloring pipeline, the residual side-channel audit returns non-trivial $\hat{\delta}$ for most reported cells; Case B reports the raw permutation-calibrated residual ceiling, and Case C reports the complexity-debiased 50/50 blend, so the Case-C ceiling is materially smaller than Case B.

\paragraph{Case A: fixed-ceiling certification on grouped 3-Colouring.}
Under the fixed Case-A ceiling of $0.25$, certification is markedly short-range.
Qwen3-14B, QwQ-32B, gpt-oss-20b, and Llama-Nemotron-Super-49B v1.5 (reasoning-on) are positive at $N=8$, but only gpt-oss-20b remains positive at $N=10$.
Qwen3-32B and Nemotron-H-47B-Reasoning-128K (reasoning-on) never clear the final signed Case-A gate on the current graph-coloring runs.
Thus, unlike the 2-SAT setting, the graph-coloring Case-A bottleneck is already visible at the smallest reported $N$ values.

\begin{table*}[ht]
\centering
\caption{Open-weight 3-Colouring per-$N$ metrics under Case A (fixed $\shortcutCeiling=0.25$).
$\mathrm{nRMC}=\max(0,\RMC^{\pm})/(1-\shortcutCeiling)=\max(0,\RMC^{\pm})/0.75$.
Three estimators are reported: RMC(no LCB) uses the point estimate $\hat{s}$; RMC(CP) uses the Clopper--Pearson factorized LCB with Bonferroni split; RMC(Wilson) uses the Wilson factorized LCB with Bonferroni split.
Model abbreviations: Q14B = Qwen3-14B (think); Q32B = Qwen3-32B (think); QwQ = QwQ-32B (think); gpt-oss = gpt-oss-20b (reason); NH-47B = Nemotron-H-47B-Reasoning-128K (reasoning-on); LN-49B = Llama-3.3-Nemotron-Super-49B v1.5 (reasoning-on).}
\label{tab:gc_os_pern}
\renewcommand{\arraystretch}{1.10}
\resizebox{\textwidth}{!}{%
\begin{tabular}{p{1.55cm}r rrrr rr rr r rr rr rr}
\toprule
\multirow{2}{*}{Model} & \multirow{2}{*}{$N$}
  & \multirow{2}{*}{Acc}
  & \multicolumn{2}{c}{F1 (per class)}
  & \multirow{2}{*}{MCC}
  & \multirow{2}{*}{G-mean}
  & \multirow{2}{*}{BalAcc}
  & \multirow{2}{*}{ADR} & \multirow{2}{*}{ADR$^{+a}$}
  & \multirow{2}{*}{$\hat{s}$}
  & \multicolumn{2}{c}{RMC(no LCB)}
  & \multicolumn{2}{c}{RMC(CP)}
  & \multicolumn{2}{c}{RMC(Wilson)} \\
\cmidrule(lr){4-5}\cmidrule(lr){12-13}\cmidrule(lr){14-15}\cmidrule(lr){16-17}
 & & & F1$_{C}$ & F1$_{N}$ & MCC & & & & & &
  $\RMC$ & $\mathrm{nRMC}$
  & $\RMC^{\pm}$ & $\mathrm{nRMC}$
  & $\RMC^{\pm}$ & $\mathrm{nRMC}$ \\
\midrule
\multirow{8}{*}{\makecell[l]{\textbf{Q14B}\\\textbf{(think)}}}
  & 8 & 0.844 & 0.816 & 0.865 & 0.722 & 0.830 & 0.844 & 0.688 & 0.686 & 0.689 & \cellcolor{green!15}$+$0.439 & \cellcolor{green!15}0.585 & \cellcolor{green!15}$+$0.395 & \cellcolor{green!15}0.526 & \cellcolor{green!15}$+$0.395 & \cellcolor{green!15}0.527 \\
  & 10 & 0.755 & 0.680 & 0.802 & 0.585 & 0.717 & 0.757 & 0.517 & 0.507 & 0.515 & \cellcolor{green!15}$+$0.265 & \cellcolor{green!15}0.353 & \cellcolor{green!15}$+$0.210 & \cellcolor{green!15}0.280 & \cellcolor{green!15}$+$0.211 & \cellcolor{green!15}0.281 \\
  & 15 & 0.565 & 0.455 & 0.638 & 0.144 & 0.528 & 0.566 & 0.288 & 0.116 & 0.279 & \cellcolor{green!15}$+$0.029 & \cellcolor{green!15}0.038 & $-$0.013 & 0.000 & $-$0.012 & 0.000 \\
  & 20 & 0.519 & 0.604 & 0.386 & 0.042 & 0.472 & 0.519 & 0.230 & 0.000 & 0.223 & $-$0.027 & 0.000 & $-$0.072 & 0.000 & $-$0.071 & 0.000 \\
  & 30 & 0.505 & 0.661 & 0.080 & 0.026 & 0.204 & 0.505 & 0.038 & 0.000 & 0.042 & $-$0.208 & 0.000 & $-$0.226 & 0.000 & $-$0.225 & 0.000 \\
  & 40 & 0.499 & 0.656 & 0.079 & -0.006 & 0.202 & 0.499 & 0.035 & 0.000 & 0.041 & $-$0.209 & 0.000 & $-$0.227 & 0.000 & $-$0.225 & 0.000 \\
  & 50 & 0.539 & 0.693 & 0.072 & 0.015 & 0.194 & 0.503 & 0.039 & 0.000 & 0.038 & $-$0.212 & 0.000 & $-$0.243 & 0.000 & $-$0.238 & 0.000 \\
  & 60 & -- & -- & -- & -- & -- & -- & -- & -- & -- & -- & -- & -- & -- & -- & -- \\
\midrule
\multirow{8}{*}{\makecell[l]{\textbf{Q32B}\\\textbf{(think)}}}
  & 8 & 0.816 & 0.777 & 0.843 & 0.673 & 0.797 & 0.816 & 0.634 & 0.628 & 0.635 & \cellcolor{green!15}$+$0.385 & \cellcolor{green!15}0.514 & \cellcolor{green!15}$+$0.321 & \cellcolor{green!15}0.428 & \cellcolor{green!15}$+$0.322 & \cellcolor{green!15}0.429 \\
  & 10 & 0.692 & 0.564 & 0.762 & 0.474 & 0.627 & 0.692 & 0.389 & 0.369 & 0.393 & \cellcolor{green!15}$+$0.143 & \cellcolor{green!15}0.191 & \cellcolor{green!15}$+$0.066 & \cellcolor{green!15}0.088 & \cellcolor{green!15}$+$0.069 & \cellcolor{green!15}0.092 \\
  & 15 & 0.567 & 0.349 & 0.675 & 0.193 & 0.457 & 0.571 & 0.210 & 0.129 & 0.209 & $-$0.041 & 0.000 & $-$0.088 & 0.000 & $-$0.086 & 0.000 \\
  & 20 & 0.546 & 0.467 & 0.604 & 0.100 & 0.526 & 0.547 & 0.278 & 0.052 & 0.276 & \cellcolor{green!15}$+$0.026 & \cellcolor{green!15}0.035 & $-$0.045 & 0.000 & $-$0.042 & 0.000 \\
  & 30 & 0.497 & 0.586 & 0.361 & -0.006 & 0.449 & 0.497 & 0.214 & 0.000 & 0.202 & $-$0.048 & 0.000 & $-$0.109 & 0.000 & $-$0.106 & 0.000 \\
  & 40 & 0.508 & 0.668 & 0.052 & 0.059 & 0.163 & 0.508 & 0.027 & 0.000 & 0.026 & $-$0.224 & 0.000 & $-$0.242 & 0.000 & $-$0.239 & 0.000 \\
  & 50 & 0.583 & 0.722 & 0.167 & 0.227 & 0.302 & 0.545 & 0.091 & 0.000 & 0.091 & $-$0.159 & 0.000 & $-$0.233 & 0.000 & $-$0.221 & 0.000 \\
  & 60 & -- & -- & -- & -- & -- & -- & -- & -- & -- & -- & -- & -- & -- & -- & -- \\
\midrule
\multirow{8}{*}{\makecell[l]{\textbf{QwQ}\\\textbf{(think)}}}
  & 8 & 0.917 & 0.910 & 0.924 & 0.846 & 0.913 & 0.917 & 0.834 & 0.824 & 0.834 & \cellcolor{green!15}$+$0.584 & \cellcolor{green!15}0.779 & \cellcolor{green!15}$+$0.527 & \cellcolor{green!15}0.703 & \cellcolor{green!15}$+$0.528 & \cellcolor{green!15}0.704 \\
  & 10 & 0.835 & 0.802 & 0.858 & 0.709 & 0.818 & 0.835 & 0.669 & 0.669 & 0.669 & \cellcolor{green!15}$+$0.419 & \cellcolor{green!15}0.559 & \cellcolor{green!15}$+$0.311 & \cellcolor{green!15}0.414 & \cellcolor{green!15}$+$0.314 & \cellcolor{green!15}0.418 \\
  & 15 & 0.599 & 0.343 & 0.711 & 0.325 & 0.455 & 0.601 & 0.211 & 0.201 & 0.207 & $-$0.043 & 0.000 & $-$0.090 & 0.000 & $-$0.088 & 0.000 \\
  & 20 & 0.550 & 0.180 & 0.689 & 0.228 & 0.315 & 0.550 & 0.099 & 0.091 & 0.099 & $-$0.151 & 0.000 & $-$0.199 & 0.000 & $-$0.194 & 0.000 \\
  & 30 & 0.504 & 0.189 & 0.643 & 0.013 & 0.321 & 0.504 & 0.099 & 0.008 & 0.103 & $-$0.147 & 0.000 & $-$0.197 & 0.000 & $-$0.192 & 0.000 \\
  & 40 & 0.546 & 0.525 & 0.565 & 0.097 & 0.545 & 0.548 & 0.282 & 0.000 & 0.297 & \cellcolor{green!15}$+$0.047 & \cellcolor{green!15}0.062 & $-$0.049 & 0.000 & $-$0.045 & 0.000 \\
  & 50 & 0.500 & 0.645 & 0.154 & 0.000 & 0.287 & 0.500 & 0.091 & 0.000 & 0.083 & $-$0.167 & 0.000 & $-$0.249 & 0.000 & $-$0.240 & 0.000 \\
  & 60 & -- & -- & -- & -- & -- & -- & -- & -- & -- & -- & -- & -- & -- & -- & -- \\
\midrule
\multirow{8}{*}{\makecell[l]{\textbf{gpt-oss}\\\textbf{(reason)}}}
  & 8 & 0.871 & 0.861 & 0.879 & 0.752 & 0.868 & 0.872 & 0.786 & 0.762 & 0.753 & \cellcolor{green!15}$+$0.503 & \cellcolor{green!15}0.671 & \cellcolor{green!15}$+$0.286 & \cellcolor{green!15}0.382 & \cellcolor{green!15}$+$0.296 & \cellcolor{green!15}0.394 \\
  & 10 & 0.849 & 0.831 & 0.863 & 0.723 & 0.843 & 0.852 & 0.714 & 0.714 & 0.710 & \cellcolor{green!15}$+$0.460 & \cellcolor{green!15}0.613 & \cellcolor{green!15}$+$0.250 & \cellcolor{green!15}0.334 & \cellcolor{green!15}$+$0.260 & \cellcolor{green!15}0.346 \\
  & 15 & 0.547 & 0.264 & 0.672 & 0.182 & 0.389 & 0.556 & 0.143 & 0.119 & 0.152 & $-$0.098 & 0.000 & $-$0.194 & 0.000 & $-$0.183 & 0.000 \\
  & 20 & 0.643 & 0.444 & 0.737 & 0.372 & 0.535 & 0.635 & 0.300 & 0.300 & 0.286 & \cellcolor{green!15}$+$0.036 & \cellcolor{green!15}0.048 & $-$0.109 & 0.000 & $-$0.095 & 0.000 \\
  & 30 & 0.506 & 0.050 & 0.667 & -0.059 & 0.160 & 0.489 & 0.028 & 0.000 & 0.026 & $-$0.224 & 0.000 & $-$0.249 & 0.000 & $-$0.246 & 0.000 \\
  & 40 & 0.537 & 0.194 & 0.675 & 0.076 & 0.327 & 0.522 & 0.048 & 0.000 & 0.107 & $-$0.143 & 0.000 & $-$0.231 & 0.000 & $-$0.219 & 0.000 \\
  & 50 & 0.500 & 0.162 & 0.644 & -0.072 & 0.296 & 0.476 & 0.069 & 0.000 & 0.088 & $-$0.162 & 0.000 & $-$0.235 & 0.000 & $-$0.225 & 0.000 \\
  & 60 & 0.571 & 0.308 & 0.690 & 0.197 & 0.426 & 0.566 & 0.133 & 0.000 & 0.181 & $-$0.069 & 0.000 & $-$0.191 & 0.000 & $-$0.177 & 0.000 \\
\midrule
\multirow{8}{*}{\makecell[l]{\textbf{NH-47B}\\\textbf{(r-on)}}}
  & 8 & 0.658 & 0.575 & 0.713 & 0.342 & 0.629 & 0.658 & 0.318 & 0.003 & 0.395 & \cellcolor{green!15}$+$0.145 & \cellcolor{green!15}0.193 & \cellcolor{green!15}$+$0.122 & \cellcolor{green!15}0.163 & \cellcolor{green!15}$+$0.123 & \cellcolor{green!15}0.164 \\
  & 10 & 0.629 & 0.626 & 0.632 & 0.258 & 0.629 & 0.629 & 0.272 & 0.002 & 0.396 & \cellcolor{green!15}$+$0.146 & \cellcolor{green!15}0.194 & \cellcolor{green!15}$+$0.122 & \cellcolor{green!15}0.163 & \cellcolor{green!15}$+$0.123 & \cellcolor{green!15}0.163 \\
  & 15 & 0.584 & 0.502 & 0.642 & 0.177 & 0.560 & 0.584 & 0.190 & 0.000 & 0.314 & \cellcolor{green!15}$+$0.064 & \cellcolor{green!15}0.085 & \cellcolor{green!15}$+$0.043 & \cellcolor{green!15}0.057 & \cellcolor{green!15}$+$0.043 & \cellcolor{green!15}0.058 \\
  & 20 & 0.560 & 0.422 & 0.645 & 0.137 & 0.507 & 0.560 & 0.154 & 0.000 & 0.257 & \cellcolor{green!15}$+$0.007 & \cellcolor{green!15}0.009 & $-$0.012 & 0.000 & $-$0.012 & 0.000 \\
  & 30 & 0.527 & 0.304 & 0.642 & 0.071 & 0.418 & 0.527 & 0.106 & 0.000 & 0.175 & $-$0.075 & 0.000 & $-$0.091 & 0.000 & $-$0.091 & 0.000 \\
  & 40 & 0.504 & 0.060 & 0.663 & 0.025 & 0.176 & 0.504 & 0.021 & 0.000 & 0.031 & $-$0.219 & 0.000 & $-$0.225 & 0.000 & $-$0.225 & 0.000 \\
  & 50 & 0.504 & 0.040 & 0.666 & 0.033 & 0.142 & 0.504 & 0.012 & 0.000 & 0.020 & $-$0.230 & 0.000 & $-$0.235 & 0.000 & $-$0.235 & 0.000 \\
  & 60 & 0.500 & -- & 0.667 & -- & 0.000 & 0.500 & 0.000 & 0.000 & 0.000 & $-$0.250 & 0.000 & $-$0.250 & 0.000 & $-$0.250 & 0.000 \\
\midrule
\multirow{8}{*}{\makecell[l]{\textbf{LN-49B}\\\textbf{(r-on)}}}
  & 8 & 0.943 & 0.940 & 0.946 & 0.892 & 0.941 & 0.943 & 0.886 & 0.886 & 0.886 & \cellcolor{green!15}$+$0.636 & \cellcolor{green!15}0.848 & \cellcolor{green!15}$+$0.444 & \cellcolor{green!15}0.592 & \cellcolor{green!15}$+$0.449 & \cellcolor{green!15}0.599 \\
  & 10 & 0.818 & 0.778 & 0.846 & 0.683 & 0.798 & 0.818 & 0.636 & 0.636 & 0.636 & \cellcolor{green!15}$+$0.386 & \cellcolor{green!15}0.515 & \cellcolor{green!15}$+$0.189 & \cellcolor{green!15}0.252 & \cellcolor{green!15}$+$0.199 & \cellcolor{green!15}0.266 \\
  & 15 & 0.557 & 0.235 & 0.688 & 0.210 & 0.365 & 0.557 & 0.136 & 0.136 & 0.133 & $-$0.117 & 0.000 & $-$0.204 & 0.000 & $-$0.194 & 0.000 \\
  & 20 & 0.545 & 0.167 & 0.688 & 0.218 & 0.302 & 0.545 & 0.091 & 0.091 & 0.091 & $-$0.159 & 0.000 & $-$0.227 & 0.000 & $-$0.217 & 0.000 \\
  & 30 & 0.477 & 0.148 & 0.623 & -0.072 & 0.280 & 0.477 & 0.068 & 0.000 & 0.079 & $-$0.171 & 0.000 & $-$0.232 & 0.000 & $-$0.224 & 0.000 \\
  & 40 & 0.523 & 0.400 & 0.604 & 0.050 & 0.481 & 0.523 & 0.205 & 0.000 & 0.231 & $-$0.019 & 0.000 & $-$0.144 & 0.000 & $-$0.134 & 0.000 \\
  & 50 & 0.598 & 0.598 & 0.598 & 0.196 & 0.598 & 0.598 & 0.349 & 0.000 & 0.357 & \cellcolor{green!15}$+$0.107 & \cellcolor{green!15}0.143 & $-$0.058 & 0.000 & $-$0.048 & 0.000 \\
  & 60 & 0.439 & 0.351 & 0.507 & -0.126 & 0.418 & 0.439 & 0.152 & 0.000 & 0.174 & $-$0.076 & 0.000 & $-$0.189 & 0.000 & $-$0.179 & 0.000 \\
\bottomrule
\multicolumn{17}{l}{\colorbox{green!15}{\phantom{X}} = that estimator's $\RMC^{\pm}>0$.}
\end{tabular}%
}
\end{table*}

\paragraph{Case B sensitivity: residual side-channel ceiling on grouped 3-Colouring.}
The current graph-coloring residual battery now produces non-trivial $\hat{\delta}$ on most reported cells, with per-$N$ residual ceilings $\shortcutCeiling$ typically in the $0.15$--$0.30$ range.
Consequently, the Case-B certificate charges these residual shortcuts before reporting the certified margin.
Under this Case B certificate, the residual ceiling is non-trivial across $N$, and only small-$N$ positive margins remain: QwQ-32B and gpt-oss-20b stay positive at $N=8$ and $N=10$, while Qwen3-14B and Llama-Nemotron-Super-49B v1.5 (reasoning-on) keep marginal positive certificates at $N=8$; Qwen3-32B and Nemotron-H-47B-Reasoning-128K (reasoning-on) never clear the Case-B gate.
This pattern indicates that, on the current grouped 3-Colouring runs, the primary failure mode is not a measured residual side channel but the difficulty of sustaining certified constructive evidence as $N$ grows.

\begin{table*}[ht]
\centering
\caption{Open-weight 3-Colouring per-$N$ metrics under Case B (estimated $\shortcutCeiling$).
$\mathrm{nRMC}=\max(0,\RMC^{\pm})/(1-\shortcutCeiling)$.
Three estimators are reported: RMC(no LCB) uses $\hat{s}$; RMC(CP) uses the Clopper--Pearson factorized LCB; RMC(Wilson) uses the Wilson factorized LCB.
Model abbreviations: Q14B = Qwen3-14B (think); Q32B = Qwen3-32B (think); QwQ = QwQ-32B (think); gpt-oss = gpt-oss-20b (reason); NH-47B = Nemotron-H-47B-Reasoning-128K (reasoning-on); LN-49B = Llama-3.3-Nemotron-Super-49B v1.5 (reasoning-on).}
\label{tab:gc_os_pern_caseB}
\renewcommand{\arraystretch}{1.10}
\resizebox{\textwidth}{!}{%
\begin{tabular}{p{1.55cm}r rrrr rr rr rr r rr rr rr}
\toprule
\multirow{2}{*}{Model} & \multirow{2}{*}{$N$}
  & \multirow{2}{*}{Acc}
  & \multicolumn{2}{c}{F1 (per class)}
  & \multirow{2}{*}{MCC}
  & \multirow{2}{*}{G-mean}
  & \multirow{2}{*}{BalAcc}
  & \multirow{2}{*}{ADR} & \multirow{2}{*}{ADR$^{+a}$}
  & \multirow{2}{*}{$\shortcutCeiling$}
  & \multirow{2}{*}{$\hat{\delta}$}
  & \multirow{2}{*}{$\hat{s}$}
  & \multicolumn{2}{c}{RMC(no LCB)}
  & \multicolumn{2}{c}{RMC(CP)}
  & \multicolumn{2}{c}{RMC(Wilson)} \\
\cmidrule(lr){4-5}\cmidrule(lr){14-15}\cmidrule(lr){16-17}\cmidrule(lr){18-19}
 & & & F1$_{C}$ & F1$_{N}$ & MCC & & & & & & & &
  $\RMC$ & $\mathrm{nRMC}$
  & $\RMC^{\pm}$ & $\mathrm{nRMC}$
  & $\RMC^{\pm}$ & $\mathrm{nRMC}$ \\
\midrule
\multirow{8}{*}{\makecell[l]{\textbf{Q14B}\\\textbf{(think)}}}
  & 8 & 0.844 & 0.816 & 0.865 & 0.722 & 0.830 & 0.844 & 0.688 & 0.686 & 0.116 & 0.208 & 0.689 & \cellcolor{green!15}$+$0.573 & \cellcolor{green!15}0.648 & \cellcolor{green!15}$+$0.529 & \cellcolor{green!15}0.598 & \cellcolor{green!15}$+$0.529 & \cellcolor{green!15}0.599 \\
  & 10 & 0.755 & 0.680 & 0.802 & 0.585 & 0.717 & 0.757 & 0.517 & 0.507 & 0.100 & 0.181 & 0.515 & \cellcolor{green!15}$+$0.415 & \cellcolor{green!15}0.461 & \cellcolor{green!15}$+$0.360 & \cellcolor{green!15}0.400 & \cellcolor{green!15}$+$0.361 & \cellcolor{green!15}0.401 \\
  & 15 & 0.565 & 0.455 & 0.638 & 0.144 & 0.528 & 0.566 & 0.288 & 0.116 & 0.116 & 0.208 & 0.279 & \cellcolor{green!15}$+$0.166 & \cellcolor{green!15}0.188 & \cellcolor{green!15}$+$0.124 & \cellcolor{green!15}0.141 & \cellcolor{green!15}$+$0.125 & \cellcolor{green!15}0.142 \\
  & 20 & 0.519 & 0.604 & 0.386 & 0.042 & 0.472 & 0.519 & 0.230 & 0.000 & 0.112 & 0.201 & 0.223 & \cellcolor{green!15}$+$0.106 & \cellcolor{green!15}0.120 & \cellcolor{green!15}$+$0.062 & \cellcolor{green!15}0.070 & \cellcolor{green!15}$+$0.063 & \cellcolor{green!15}0.071 \\
  & 30 & 0.505 & 0.661 & 0.080 & 0.026 & 0.204 & 0.505 & 0.038 & 0.000 & 0.124 & 0.221 & 0.042 & $-$0.090 & 0.000 & $-$0.105 & 0.000 & $-$0.104 & 0.000 \\
  & 40 & 0.499 & 0.656 & 0.079 & -0.006 & 0.202 & 0.499 & 0.035 & 0.000 & 0.110 & 0.197 & 0.041 & $-$0.069 & 0.000 & $-$0.086 & 0.000 & $-$0.085 & 0.000 \\
  & 50 & 0.539 & 0.693 & 0.072 & 0.015 & 0.194 & 0.503 & 0.039 & 0.000 & 0.049 & 0.091 & 0.038 & $-$0.011 & 0.000 & $-$0.042 & 0.000 & $-$0.037 & 0.000 \\
  & 60 & -- & -- & -- & -- & -- & -- & -- & -- & -- & -- & -- & -- & -- & -- & -- & -- & -- \\
\midrule
\multirow{8}{*}{\makecell[l]{\textbf{Q32B}\\\textbf{(think)}}}
  & 8 & 0.816 & 0.777 & 0.843 & 0.673 & 0.797 & 0.816 & 0.634 & 0.628 & 0.121 & 0.217 & 0.635 & \cellcolor{green!15}$+$0.512 & \cellcolor{green!15}0.582 & \cellcolor{green!15}$+$0.447 & \cellcolor{green!15}0.508 & \cellcolor{green!15}$+$0.448 & \cellcolor{green!15}0.510 \\
  & 10 & 0.692 & 0.564 & 0.762 & 0.474 & 0.627 & 0.692 & 0.389 & 0.369 & 0.092 & 0.167 & 0.393 & \cellcolor{green!15}$+$0.301 & \cellcolor{green!15}0.331 & \cellcolor{green!15}$+$0.224 & \cellcolor{green!15}0.247 & \cellcolor{green!15}$+$0.227 & \cellcolor{green!15}0.250 \\
  & 15 & 0.567 & 0.349 & 0.675 & 0.193 & 0.457 & 0.571 & 0.210 & 0.129 & 0.117 & 0.210 & 0.209 & \cellcolor{green!15}$+$0.095 & \cellcolor{green!15}0.107 & \cellcolor{green!15}$+$0.047 & \cellcolor{green!15}0.053 & \cellcolor{green!15}$+$0.049 & \cellcolor{green!15}0.055 \\
  & 20 & 0.546 & 0.467 & 0.604 & 0.100 & 0.526 & 0.547 & 0.278 & 0.052 & 0.100 & 0.181 & 0.276 & \cellcolor{green!15}$+$0.181 & \cellcolor{green!15}0.201 & \cellcolor{green!15}$+$0.109 & \cellcolor{green!15}0.122 & \cellcolor{green!15}$+$0.112 & \cellcolor{green!15}0.125 \\
  & 30 & 0.497 & 0.586 & 0.361 & -0.006 & 0.449 & 0.497 & 0.214 & 0.000 & 0.116 & 0.209 & 0.202 & \cellcolor{green!15}$+$0.085 & \cellcolor{green!15}0.096 & \cellcolor{green!15}$+$0.025 & \cellcolor{green!15}0.028 & \cellcolor{green!15}$+$0.027 & \cellcolor{green!15}0.031 \\
  & 40 & 0.508 & 0.668 & 0.052 & 0.059 & 0.163 & 0.508 & 0.027 & 0.000 & 0.103 & 0.187 & 0.026 & $-$0.082 & 0.000 & $-$0.097 & 0.000 & $-$0.095 & 0.000 \\
  & 50 & 0.583 & 0.722 & 0.167 & 0.227 & 0.302 & 0.545 & 0.091 & 0.000 & 0.058 & 0.106 & 0.091 & \cellcolor{green!15}$+$0.033 & \cellcolor{green!15}0.035 & $-$0.041 & 0.000 & $-$0.029 & 0.000 \\
  & 60 & -- & -- & -- & -- & -- & -- & -- & -- & -- & -- & -- & -- & -- & -- & -- & -- & -- \\
\midrule
\multirow{8}{*}{\makecell[l]{\textbf{QwQ}\\\textbf{(think)}}}
  & 8 & 0.917 & 0.910 & 0.924 & 0.846 & 0.913 & 0.917 & 0.834 & 0.824 & 0.121 & 0.217 & 0.834 & \cellcolor{green!15}$+$0.713 & \cellcolor{green!15}0.812 & \cellcolor{green!15}$+$0.656 & \cellcolor{green!15}0.747 & \cellcolor{green!15}$+$0.656 & \cellcolor{green!15}0.747 \\
  & 10 & 0.835 & 0.802 & 0.858 & 0.709 & 0.818 & 0.835 & 0.669 & 0.669 & 0.093 & 0.169 & 0.669 & \cellcolor{green!15}$+$0.569 & \cellcolor{green!15}0.627 & \cellcolor{green!15}$+$0.460 & \cellcolor{green!15}0.507 & \cellcolor{green!15}$+$0.463 & \cellcolor{green!15}0.510 \\
  & 15 & 0.599 & 0.343 & 0.711 & 0.325 & 0.455 & 0.601 & 0.211 & 0.201 & 0.116 & 0.208 & 0.207 & \cellcolor{green!15}$+$0.090 & \cellcolor{green!15}0.102 & \cellcolor{green!15}$+$0.042 & \cellcolor{green!15}0.048 & \cellcolor{green!15}$+$0.045 & \cellcolor{green!15}0.051 \\
  & 20 & 0.550 & 0.180 & 0.689 & 0.228 & 0.315 & 0.550 & 0.099 & 0.091 & 0.091 & 0.165 & 0.099 & \cellcolor{green!15}$+$0.041 & \cellcolor{green!15}0.045 & $-$0.016 & 0.000 & $-$0.011 & 0.000 \\
  & 30 & 0.504 & 0.189 & 0.643 & 0.013 & 0.321 & 0.504 & 0.099 & 0.008 & 0.120 & 0.215 & 0.103 & \cellcolor{green!15}$+$0.004 & \cellcolor{green!15}0.005 & $-$0.052 & 0.000 & $-$0.047 & 0.000 \\
  & 40 & 0.546 & 0.525 & 0.565 & 0.097 & 0.545 & 0.548 & 0.282 & 0.000 & 0.124 & 0.223 & 0.297 & \cellcolor{green!15}$+$0.189 & \cellcolor{green!15}0.216 & \cellcolor{green!15}$+$0.090 & \cellcolor{green!15}0.103 & \cellcolor{green!15}$+$0.095 & \cellcolor{green!15}0.108 \\
  & 50 & 0.500 & 0.645 & 0.154 & 0.000 & 0.287 & 0.500 & 0.091 & 0.000 & 0.048 & 0.091 & 0.083 & \cellcolor{green!15}$+$0.035 & \cellcolor{green!15}0.037 & $-$0.046 & 0.000 & $-$0.037 & 0.000 \\
  & 60 & -- & -- & -- & -- & -- & -- & -- & -- & -- & -- & -- & -- & -- & -- & -- & -- & -- \\
\midrule
\multirow{8}{*}{\makecell[l]{\textbf{gpt-oss}\\\textbf{(reason)}}}
  & 8 & 0.871 & 0.861 & 0.879 & 0.752 & 0.868 & 0.872 & 0.786 & 0.762 & 0.087 & 0.159 & 0.753 & \cellcolor{green!15}$+$0.666 & \cellcolor{green!15}0.729 & \cellcolor{green!15}$+$0.449 & \cellcolor{green!15}0.492 & \cellcolor{green!15}$+$0.458 & \cellcolor{green!15}0.502 \\
  & 10 & 0.849 & 0.831 & 0.863 & 0.723 & 0.843 & 0.852 & 0.714 & 0.714 & 0.068 & 0.125 & 0.710 & \cellcolor{green!15}$+$0.642 & \cellcolor{green!15}0.689 & \cellcolor{green!15}$+$0.433 & \cellcolor{green!15}0.464 & \cellcolor{green!15}$+$0.442 & \cellcolor{green!15}0.474 \\
  & 15 & 0.547 & 0.264 & 0.672 & 0.182 & 0.389 & 0.556 & 0.143 & 0.119 & 0.106 & 0.193 & 0.152 & \cellcolor{green!15}$+$0.045 & \cellcolor{green!15}0.051 & $-$0.051 & 0.000 & $-$0.040 & 0.000 \\
  & 20 & 0.643 & 0.444 & 0.737 & 0.372 & 0.535 & 0.635 & 0.300 & 0.300 & 0.075 & 0.136 & 0.286 & \cellcolor{green!15}$+$0.211 & \cellcolor{green!15}0.228 & \cellcolor{green!15}$+$0.066 & \cellcolor{green!15}0.072 & \cellcolor{green!15}$+$0.080 & \cellcolor{green!15}0.086 \\
  & 30 & 0.506 & 0.050 & 0.667 & -0.059 & 0.160 & 0.489 & 0.028 & 0.000 & 0.113 & 0.205 & 0.026 & $-$0.087 & 0.000 & $-$0.112 & 0.000 & $-$0.109 & 0.000 \\
  & 40 & 0.537 & 0.194 & 0.675 & 0.076 & 0.327 & 0.522 & 0.048 & 0.000 & 0.136 & 0.242 & 0.107 & $-$0.029 & 0.000 & $-$0.117 & 0.000 & $-$0.105 & 0.000 \\
  & 50 & 0.500 & 0.162 & 0.644 & -0.072 & 0.296 & 0.476 & 0.069 & 0.000 & 0.127 & 0.227 & 0.088 & $-$0.039 & 0.000 & $-$0.112 & 0.000 & $-$0.102 & 0.000 \\
  & 60 & 0.571 & 0.308 & 0.690 & 0.197 & 0.426 & 0.566 & 0.133 & 0.000 & 0.127 & 0.227 & 0.181 & \cellcolor{green!15}$+$0.055 & \cellcolor{green!15}0.063 & $-$0.068 & 0.000 & $-$0.053 & 0.000 \\
\midrule
\multirow{8}{*}{\makecell[l]{\textbf{NH-47B}\\\textbf{(r-on)}}}
  & 8 & 0.658 & 0.575 & 0.713 & 0.342 & 0.629 & 0.658 & 0.318 & 0.003 & 0.109 & 0.196 & 0.395 & \cellcolor{green!15}$+$0.286 & \cellcolor{green!15}0.321 & \cellcolor{green!15}$+$0.264 & \cellcolor{green!15}0.296 & \cellcolor{green!15}$+$0.264 & \cellcolor{green!15}0.296 \\
  & 10 & 0.629 & 0.626 & 0.632 & 0.258 & 0.629 & 0.629 & 0.272 & 0.002 & 0.104 & 0.188 & 0.396 & \cellcolor{green!15}$+$0.291 & \cellcolor{green!15}0.325 & \cellcolor{green!15}$+$0.268 & \cellcolor{green!15}0.299 & \cellcolor{green!15}$+$0.268 & \cellcolor{green!15}0.300 \\
  & 15 & 0.584 & 0.502 & 0.642 & 0.177 & 0.560 & 0.584 & 0.190 & 0.000 & 0.111 & 0.200 & 0.314 & \cellcolor{green!15}$+$0.202 & \cellcolor{green!15}0.228 & \cellcolor{green!15}$+$0.181 & \cellcolor{green!15}0.204 & \cellcolor{green!15}$+$0.182 & \cellcolor{green!15}0.204 \\
  & 20 & 0.560 & 0.422 & 0.645 & 0.137 & 0.507 & 0.560 & 0.154 & 0.000 & 0.107 & 0.193 & 0.257 & \cellcolor{green!15}$+$0.149 & \cellcolor{green!15}0.167 & \cellcolor{green!15}$+$0.130 & \cellcolor{green!15}0.146 & \cellcolor{green!15}$+$0.130 & \cellcolor{green!15}0.146 \\
  & 30 & 0.527 & 0.304 & 0.642 & 0.071 & 0.418 & 0.527 & 0.106 & 0.000 & 0.105 & 0.190 & 0.175 & \cellcolor{green!15}$+$0.070 & \cellcolor{green!15}0.078 & \cellcolor{green!15}$+$0.054 & \cellcolor{green!15}0.060 & \cellcolor{green!15}$+$0.054 & \cellcolor{green!15}0.061 \\
  & 40 & 0.504 & 0.060 & 0.663 & 0.025 & 0.176 & 0.504 & 0.021 & 0.000 & 0.105 & 0.189 & 0.031 & $-$0.074 & 0.000 & $-$0.080 & 0.000 & $-$0.080 & 0.000 \\
  & 50 & 0.504 & 0.040 & 0.666 & 0.033 & 0.142 & 0.504 & 0.012 & 0.000 & 0.109 & 0.197 & 0.020 & $-$0.089 & 0.000 & $-$0.094 & 0.000 & $-$0.094 & 0.000 \\
  & 60 & 0.500 & -- & 0.667 & -- & 0.000 & 0.500 & 0.000 & 0.000 & 0.119 & 0.212 & 0.000 & $-$0.119 & 0.000 & $-$0.119 & 0.000 & $-$0.119 & 0.000 \\
\midrule
\multirow{8}{*}{\makecell[l]{\textbf{LN-49B}\\\textbf{(r-on)}}}
  & 8 & 0.943 & 0.940 & 0.946 & 0.892 & 0.941 & 0.943 & 0.886 & 0.886 & 0.087 & 0.159 & 0.886 & \cellcolor{green!15}$+$0.776 & \cellcolor{green!15}0.851 & \cellcolor{green!15}$+$0.581 & \cellcolor{green!15}0.636 & \cellcolor{green!15}$+$0.587 & \cellcolor{green!15}0.643 \\
  & 10 & 0.818 & 0.778 & 0.846 & 0.683 & 0.798 & 0.818 & 0.636 & 0.636 & 0.068 & 0.125 & 0.636 & \cellcolor{green!15}$+$0.569 & \cellcolor{green!15}0.610 & \cellcolor{green!15}$+$0.372 & \cellcolor{green!15}0.399 & \cellcolor{green!15}$+$0.382 & \cellcolor{green!15}0.410 \\
  & 15 & 0.557 & 0.235 & 0.688 & 0.210 & 0.365 & 0.557 & 0.136 & 0.136 & 0.106 & 0.193 & 0.133 & \cellcolor{green!15}$+$0.027 & \cellcolor{green!15}0.030 & $-$0.061 & 0.000 & $-$0.050 & 0.000 \\
  & 20 & 0.545 & 0.167 & 0.688 & 0.218 & 0.302 & 0.545 & 0.091 & 0.091 & 0.075 & 0.136 & 0.091 & $-$0.007 & 0.000 & $-$0.062 & 0.000 & $-$0.054 & 0.000 \\
  & 30 & 0.477 & 0.148 & 0.623 & -0.072 & 0.280 & 0.477 & 0.068 & 0.000 & 0.113 & 0.205 & 0.079 & \cellcolor{green!15}$+$0.073 & \cellcolor{green!15}0.082 & $-$0.036 & 0.000 & $-$0.025 & 0.000 \\
  & 40 & 0.523 & 0.400 & 0.604 & 0.050 & 0.481 & 0.523 & 0.205 & 0.000 & 0.134 & 0.239 & 0.231 & \cellcolor{green!15}$+$0.060 & \cellcolor{green!15}0.069 & $-$0.050 & 0.000 & $-$0.041 & 0.000 \\
  & 50 & 0.598 & 0.598 & 0.598 & 0.196 & 0.598 & 0.598 & 0.349 & 0.000 & 0.127 & 0.228 & 0.357 & \cellcolor{green!15}$+$0.228 & \cellcolor{green!15}0.261 & \cellcolor{green!15}$+$0.063 & \cellcolor{green!15}0.073 & \cellcolor{green!15}$+$0.074 & \cellcolor{green!15}0.085 \\
  & 60 & 0.439 & 0.351 & 0.507 & -0.126 & 0.418 & 0.439 & 0.152 & 0.000 & 0.127 & 0.227 & 0.174 & \cellcolor{green!15}$+$0.079 & \cellcolor{green!15}0.091 & $-$0.048 & 0.000 & $-$0.038 & 0.000 \\
\bottomrule
\multicolumn{19}{l}{\colorbox{green!15}{\phantom{X}} = that estimator's $\RMC^{\pm}>0$.}
\end{tabular}%
}
\end{table*}

\paragraph{Case C sensitivity: complexity-debiased residual ceiling on grouped 3-Colouring.}
On the current graph-coloring runs, Case C applies the complexity-debiased residualization and 50/50 blend with the Case-A reference, yielding residual ceilings $\shortcutCeiling$ typically in the $0.06$--$0.13$ range.
Relative to raw Case B, the Case-C ceiling is materially smaller because generic structural complexity is partialled out before the shortcut audit.
Under Case C, the small-$N$ positive cells extend further: Qwen3-14B, Qwen3-32B, QwQ-32B, and Llama-Nemotron-Super-49B v1.5 (reasoning-on) keep positive margins at $N=8$ and $N=10$, gpt-oss-20b remains positive through $N=10$, and Nemotron-H-47B-Reasoning-128K (reasoning-on) remains negative throughout.

\begin{table*}[ht]
\centering
\caption{Open-weight 3-Colouring per-$N$ metrics under Case C (complexity-debiased estimated $\shortcutCeiling$).
$\mathrm{nRMC}=\max(0,\RMC^{\pm})/(1-\shortcutCeiling)$.
Three estimators are reported: RMC(no LCB) uses $\hat{s}$; RMC(CP) uses the Clopper--Pearson factorized LCB; RMC(Wilson) uses the Wilson factorized LCB.
Model abbreviations: Q14B = Qwen3-14B (think); Q32B = Qwen3-32B (think); QwQ = QwQ-32B (think); gpt-oss = gpt-oss-20b (reason); NH-47B = Nemotron-H-47B-Reasoning-128K (reasoning-on); LN-49B = Llama-3.3-Nemotron-Super-49B v1.5 (reasoning-on).}
\label{tab:gc_os_pern_caseC}
\renewcommand{\arraystretch}{1.10}
\resizebox{\textwidth}{!}{%
\begin{tabular}{p{1.55cm}r rrrr rr rr rr r rr rr rr}
\toprule
\multirow{2}{*}{Model} & \multirow{2}{*}{$N$}
  & \multirow{2}{*}{Acc}
  & \multicolumn{2}{c}{F1 (per class)}
  & \multirow{2}{*}{MCC}
  & \multirow{2}{*}{G-mean}
  & \multirow{2}{*}{BalAcc}
  & \multirow{2}{*}{ADR} & \multirow{2}{*}{ADR$^{+a}$}
  & \multirow{2}{*}{$\shortcutCeiling$}
  & \multirow{2}{*}{$\hat{\delta}$}
  & \multirow{2}{*}{$\hat{s}$}
  & \multicolumn{2}{c}{RMC(no LCB)}
  & \multicolumn{2}{c}{RMC(CP)}
  & \multicolumn{2}{c}{RMC(Wilson)} \\
\cmidrule(lr){4-5}\cmidrule(lr){14-15}\cmidrule(lr){16-17}\cmidrule(lr){18-19}
 & & & F1$_{C}$ & F1$_{N}$ & MCC & & & & & & & &
  $\RMC$ & $\mathrm{nRMC}$
  & $\RMC^{\pm}$ & $\mathrm{nRMC}$
  & $\RMC^{\pm}$ & $\mathrm{nRMC}$ \\
\midrule
\multirow{8}{*}{\makecell[l]{\textbf{Q14B}\\\textbf{(think)}}}
  & 8 & 0.844 & 0.816 & 0.865 & 0.722 & 0.830 & 0.844 & 0.688 & 0.686 & 0.116 & 0.208 & 0.689 & \cellcolor{green!15}$+$0.573 & \cellcolor{green!15}0.648 & \cellcolor{green!15}$+$0.529 & \cellcolor{green!15}0.598 & \cellcolor{green!15}$+$0.529 & \cellcolor{green!15}0.599 \\
  & 10 & 0.755 & 0.680 & 0.802 & 0.585 & 0.717 & 0.757 & 0.517 & 0.507 & 0.100 & 0.181 & 0.515 & \cellcolor{green!15}$+$0.415 & \cellcolor{green!15}0.461 & \cellcolor{green!15}$+$0.360 & \cellcolor{green!15}0.400 & \cellcolor{green!15}$+$0.361 & \cellcolor{green!15}0.401 \\
  & 15 & 0.565 & 0.455 & 0.638 & 0.144 & 0.528 & 0.566 & 0.288 & 0.116 & 0.116 & 0.208 & 0.279 & \cellcolor{green!15}$+$0.166 & \cellcolor{green!15}0.188 & \cellcolor{green!15}$+$0.124 & \cellcolor{green!15}0.141 & \cellcolor{green!15}$+$0.125 & \cellcolor{green!15}0.142 \\
  & 20 & 0.519 & 0.604 & 0.386 & 0.042 & 0.472 & 0.519 & 0.230 & 0.000 & 0.112 & 0.201 & 0.223 & \cellcolor{green!15}$+$0.106 & \cellcolor{green!15}0.120 & \cellcolor{green!15}$+$0.062 & \cellcolor{green!15}0.070 & \cellcolor{green!15}$+$0.063 & \cellcolor{green!15}0.071 \\
  & 30 & 0.505 & 0.661 & 0.080 & 0.026 & 0.204 & 0.505 & 0.038 & 0.000 & 0.124 & 0.221 & 0.042 & $-$0.090 & 0.000 & $-$0.105 & 0.000 & $-$0.104 & 0.000 \\
  & 40 & 0.499 & 0.656 & 0.079 & -0.006 & 0.202 & 0.499 & 0.035 & 0.000 & 0.110 & 0.197 & 0.041 & $-$0.069 & 0.000 & $-$0.086 & 0.000 & $-$0.085 & 0.000 \\
  & 50 & 0.539 & 0.693 & 0.072 & 0.015 & 0.194 & 0.503 & 0.039 & 0.000 & 0.049 & 0.091 & 0.038 & $-$0.011 & 0.000 & $-$0.042 & 0.000 & $-$0.037 & 0.000 \\
  & 60 & -- & -- & -- & -- & -- & -- & -- & -- & -- & -- & -- & -- & -- & -- & -- & -- & -- \\
\midrule
\multirow{8}{*}{\makecell[l]{\textbf{Q32B}\\\textbf{(think)}}}
  & 8 & 0.816 & 0.777 & 0.843 & 0.673 & 0.797 & 0.816 & 0.634 & 0.628 & 0.121 & 0.217 & 0.635 & \cellcolor{green!15}$+$0.512 & \cellcolor{green!15}0.582 & \cellcolor{green!15}$+$0.447 & \cellcolor{green!15}0.508 & \cellcolor{green!15}$+$0.448 & \cellcolor{green!15}0.510 \\
  & 10 & 0.692 & 0.564 & 0.762 & 0.474 & 0.627 & 0.692 & 0.389 & 0.369 & 0.092 & 0.167 & 0.393 & \cellcolor{green!15}$+$0.301 & \cellcolor{green!15}0.331 & \cellcolor{green!15}$+$0.224 & \cellcolor{green!15}0.247 & \cellcolor{green!15}$+$0.227 & \cellcolor{green!15}0.250 \\
  & 15 & 0.567 & 0.349 & 0.675 & 0.193 & 0.457 & 0.571 & 0.210 & 0.129 & 0.117 & 0.210 & 0.209 & \cellcolor{green!15}$+$0.095 & \cellcolor{green!15}0.107 & \cellcolor{green!15}$+$0.047 & \cellcolor{green!15}0.053 & \cellcolor{green!15}$+$0.049 & \cellcolor{green!15}0.055 \\
  & 20 & 0.546 & 0.467 & 0.604 & 0.100 & 0.526 & 0.547 & 0.278 & 0.052 & 0.100 & 0.181 & 0.276 & \cellcolor{green!15}$+$0.181 & \cellcolor{green!15}0.201 & \cellcolor{green!15}$+$0.109 & \cellcolor{green!15}0.122 & \cellcolor{green!15}$+$0.112 & \cellcolor{green!15}0.125 \\
  & 30 & 0.497 & 0.586 & 0.361 & -0.006 & 0.449 & 0.497 & 0.214 & 0.000 & 0.116 & 0.209 & 0.202 & \cellcolor{green!15}$+$0.085 & \cellcolor{green!15}0.096 & \cellcolor{green!15}$+$0.025 & \cellcolor{green!15}0.028 & \cellcolor{green!15}$+$0.027 & \cellcolor{green!15}0.031 \\
  & 40 & 0.508 & 0.668 & 0.052 & 0.059 & 0.163 & 0.508 & 0.027 & 0.000 & 0.103 & 0.187 & 0.026 & $-$0.082 & 0.000 & $-$0.097 & 0.000 & $-$0.095 & 0.000 \\
  & 50 & 0.583 & 0.722 & 0.167 & 0.227 & 0.302 & 0.545 & 0.091 & 0.000 & 0.058 & 0.106 & 0.091 & \cellcolor{green!15}$+$0.033 & \cellcolor{green!15}0.035 & $-$0.041 & 0.000 & $-$0.029 & 0.000 \\
  & 60 & -- & -- & -- & -- & -- & -- & -- & -- & -- & -- & -- & -- & -- & -- & -- & -- & -- \\
\midrule
\multirow{8}{*}{\makecell[l]{\textbf{QwQ}\\\textbf{(think)}}}
  & 8 & 0.917 & 0.910 & 0.924 & 0.846 & 0.913 & 0.917 & 0.834 & 0.824 & 0.121 & 0.217 & 0.834 & \cellcolor{green!15}$+$0.713 & \cellcolor{green!15}0.812 & \cellcolor{green!15}$+$0.656 & \cellcolor{green!15}0.747 & \cellcolor{green!15}$+$0.656 & \cellcolor{green!15}0.747 \\
  & 10 & 0.835 & 0.802 & 0.858 & 0.709 & 0.818 & 0.835 & 0.669 & 0.669 & 0.093 & 0.169 & 0.669 & \cellcolor{green!15}$+$0.569 & \cellcolor{green!15}0.627 & \cellcolor{green!15}$+$0.460 & \cellcolor{green!15}0.507 & \cellcolor{green!15}$+$0.463 & \cellcolor{green!15}0.510 \\
  & 15 & 0.599 & 0.343 & 0.711 & 0.325 & 0.455 & 0.601 & 0.211 & 0.201 & 0.116 & 0.208 & 0.207 & \cellcolor{green!15}$+$0.090 & \cellcolor{green!15}0.102 & \cellcolor{green!15}$+$0.042 & \cellcolor{green!15}0.048 & \cellcolor{green!15}$+$0.045 & \cellcolor{green!15}0.051 \\
  & 20 & 0.550 & 0.180 & 0.689 & 0.228 & 0.315 & 0.550 & 0.099 & 0.091 & 0.091 & 0.165 & 0.099 & \cellcolor{green!15}$+$0.041 & \cellcolor{green!15}0.045 & $-$0.016 & 0.000 & $-$0.011 & 0.000 \\
  & 30 & 0.504 & 0.189 & 0.643 & 0.013 & 0.321 & 0.504 & 0.099 & 0.008 & 0.120 & 0.215 & 0.103 & \cellcolor{green!15}$+$0.004 & \cellcolor{green!15}0.005 & $-$0.052 & 0.000 & $-$0.047 & 0.000 \\
  & 40 & 0.546 & 0.525 & 0.565 & 0.097 & 0.545 & 0.548 & 0.282 & 0.000 & 0.124 & 0.223 & 0.297 & \cellcolor{green!15}$+$0.189 & \cellcolor{green!15}0.216 & \cellcolor{green!15}$+$0.090 & \cellcolor{green!15}0.103 & \cellcolor{green!15}$+$0.095 & \cellcolor{green!15}0.108 \\
  & 50 & 0.500 & 0.645 & 0.154 & 0.000 & 0.287 & 0.500 & 0.091 & 0.000 & 0.048 & 0.091 & 0.083 & \cellcolor{green!15}$+$0.035 & \cellcolor{green!15}0.037 & $-$0.046 & 0.000 & $-$0.037 & 0.000 \\
  & 60 & -- & -- & -- & -- & -- & -- & -- & -- & -- & -- & -- & -- & -- & -- & -- & -- & -- \\
\midrule
\multirow{8}{*}{\makecell[l]{\textbf{gpt-oss}\\\textbf{(reason)}}}
  & 8 & 0.871 & 0.861 & 0.879 & 0.752 & 0.868 & 0.872 & 0.786 & 0.762 & 0.087 & 0.159 & 0.753 & \cellcolor{green!15}$+$0.666 & \cellcolor{green!15}0.729 & \cellcolor{green!15}$+$0.449 & \cellcolor{green!15}0.492 & \cellcolor{green!15}$+$0.458 & \cellcolor{green!15}0.502 \\
  & 10 & 0.849 & 0.831 & 0.863 & 0.723 & 0.843 & 0.852 & 0.714 & 0.714 & 0.068 & 0.125 & 0.710 & \cellcolor{green!15}$+$0.642 & \cellcolor{green!15}0.689 & \cellcolor{green!15}$+$0.433 & \cellcolor{green!15}0.464 & \cellcolor{green!15}$+$0.442 & \cellcolor{green!15}0.474 \\
  & 15 & 0.547 & 0.264 & 0.672 & 0.182 & 0.389 & 0.556 & 0.143 & 0.119 & 0.106 & 0.193 & 0.152 & \cellcolor{green!15}$+$0.045 & \cellcolor{green!15}0.051 & $-$0.051 & 0.000 & $-$0.040 & 0.000 \\
  & 20 & 0.643 & 0.444 & 0.737 & 0.372 & 0.535 & 0.635 & 0.300 & 0.300 & 0.075 & 0.136 & 0.286 & \cellcolor{green!15}$+$0.211 & \cellcolor{green!15}0.228 & \cellcolor{green!15}$+$0.066 & \cellcolor{green!15}0.072 & \cellcolor{green!15}$+$0.080 & \cellcolor{green!15}0.086 \\
  & 30 & 0.506 & 0.050 & 0.667 & -0.059 & 0.160 & 0.489 & 0.028 & 0.000 & 0.113 & 0.205 & 0.026 & $-$0.087 & 0.000 & $-$0.112 & 0.000 & $-$0.109 & 0.000 \\
  & 40 & 0.537 & 0.194 & 0.675 & 0.076 & 0.327 & 0.522 & 0.048 & 0.000 & 0.136 & 0.242 & 0.107 & $-$0.029 & 0.000 & $-$0.117 & 0.000 & $-$0.105 & 0.000 \\
  & 50 & 0.500 & 0.162 & 0.644 & -0.072 & 0.296 & 0.476 & 0.069 & 0.000 & 0.127 & 0.227 & 0.088 & $-$0.039 & 0.000 & $-$0.112 & 0.000 & $-$0.102 & 0.000 \\
  & 60 & 0.571 & 0.308 & 0.690 & 0.197 & 0.426 & 0.566 & 0.133 & 0.000 & 0.127 & 0.227 & 0.181 & \cellcolor{green!15}$+$0.055 & \cellcolor{green!15}0.063 & $-$0.068 & 0.000 & $-$0.053 & 0.000 \\
\midrule
\multirow{8}{*}{\makecell[l]{\textbf{NH-47B}\\\textbf{(r-on)}}}
  & 8 & 0.658 & 0.575 & 0.713 & 0.342 & 0.629 & 0.658 & 0.318 & 0.003 & 0.109 & 0.196 & 0.395 & \cellcolor{green!15}$+$0.286 & \cellcolor{green!15}0.321 & \cellcolor{green!15}$+$0.264 & \cellcolor{green!15}0.296 & \cellcolor{green!15}$+$0.264 & \cellcolor{green!15}0.296 \\
  & 10 & 0.629 & 0.626 & 0.632 & 0.258 & 0.629 & 0.629 & 0.272 & 0.002 & 0.104 & 0.188 & 0.396 & \cellcolor{green!15}$+$0.291 & \cellcolor{green!15}0.325 & \cellcolor{green!15}$+$0.268 & \cellcolor{green!15}0.299 & \cellcolor{green!15}$+$0.268 & \cellcolor{green!15}0.300 \\
  & 15 & 0.584 & 0.502 & 0.642 & 0.177 & 0.560 & 0.584 & 0.190 & 0.000 & 0.111 & 0.200 & 0.314 & \cellcolor{green!15}$+$0.202 & \cellcolor{green!15}0.228 & \cellcolor{green!15}$+$0.181 & \cellcolor{green!15}0.204 & \cellcolor{green!15}$+$0.182 & \cellcolor{green!15}0.204 \\
  & 20 & 0.560 & 0.422 & 0.645 & 0.137 & 0.507 & 0.560 & 0.154 & 0.000 & 0.107 & 0.193 & 0.257 & \cellcolor{green!15}$+$0.149 & \cellcolor{green!15}0.167 & \cellcolor{green!15}$+$0.130 & \cellcolor{green!15}0.146 & \cellcolor{green!15}$+$0.130 & \cellcolor{green!15}0.146 \\
  & 30 & 0.527 & 0.304 & 0.642 & 0.071 & 0.418 & 0.527 & 0.106 & 0.000 & 0.105 & 0.190 & 0.175 & \cellcolor{green!15}$+$0.070 & \cellcolor{green!15}0.078 & \cellcolor{green!15}$+$0.054 & \cellcolor{green!15}0.060 & \cellcolor{green!15}$+$0.054 & \cellcolor{green!15}0.061 \\
  & 40 & 0.504 & 0.060 & 0.663 & 0.025 & 0.176 & 0.504 & 0.021 & 0.000 & 0.105 & 0.189 & 0.031 & $-$0.074 & 0.000 & $-$0.080 & 0.000 & $-$0.080 & 0.000 \\
  & 50 & 0.504 & 0.040 & 0.666 & 0.033 & 0.142 & 0.504 & 0.012 & 0.000 & 0.109 & 0.197 & 0.020 & $-$0.089 & 0.000 & $-$0.094 & 0.000 & $-$0.094 & 0.000 \\
  & 60 & 0.500 & -- & 0.667 & -- & 0.000 & 0.500 & 0.000 & 0.000 & 0.119 & 0.212 & 0.000 & $-$0.119 & 0.000 & $-$0.119 & 0.000 & $-$0.119 & 0.000 \\
\midrule
\multirow{8}{*}{\makecell[l]{\textbf{LN-49B}\\\textbf{(r-on)}}}
  & 8 & 0.943 & 0.940 & 0.946 & 0.892 & 0.941 & 0.943 & 0.886 & 0.886 & 0.087 & 0.159 & 0.886 & \cellcolor{green!15}$+$0.776 & \cellcolor{green!15}0.851 & \cellcolor{green!15}$+$0.581 & \cellcolor{green!15}0.636 & \cellcolor{green!15}$+$0.587 & \cellcolor{green!15}0.643 \\
  & 10 & 0.818 & 0.778 & 0.846 & 0.683 & 0.798 & 0.818 & 0.636 & 0.636 & 0.068 & 0.125 & 0.636 & \cellcolor{green!15}$+$0.569 & \cellcolor{green!15}0.610 & \cellcolor{green!15}$+$0.372 & \cellcolor{green!15}0.399 & \cellcolor{green!15}$+$0.382 & \cellcolor{green!15}0.410 \\
  & 15 & 0.557 & 0.235 & 0.688 & 0.210 & 0.365 & 0.557 & 0.136 & 0.136 & 0.106 & 0.193 & 0.133 & \cellcolor{green!15}$+$0.027 & \cellcolor{green!15}0.030 & $-$0.061 & 0.000 & $-$0.050 & 0.000 \\
  & 20 & 0.545 & 0.167 & 0.688 & 0.218 & 0.302 & 0.545 & 0.091 & 0.091 & 0.075 & 0.136 & 0.091 & $-$0.007 & 0.000 & $-$0.062 & 0.000 & $-$0.054 & 0.000 \\
  & 30 & 0.477 & 0.148 & 0.623 & -0.072 & 0.280 & 0.477 & 0.068 & 0.000 & 0.113 & 0.205 & 0.079 & \cellcolor{green!15}$+$0.073 & \cellcolor{green!15}0.082 & $-$0.036 & 0.000 & $-$0.025 & 0.000 \\
  & 40 & 0.523 & 0.400 & 0.604 & 0.050 & 0.481 & 0.523 & 0.205 & 0.000 & 0.134 & 0.239 & 0.231 & \cellcolor{green!15}$+$0.060 & \cellcolor{green!15}0.069 & $-$0.050 & 0.000 & $-$0.041 & 0.000 \\
  & 50 & 0.598 & 0.598 & 0.598 & 0.196 & 0.598 & 0.598 & 0.349 & 0.000 & 0.127 & 0.228 & 0.357 & \cellcolor{green!15}$+$0.228 & \cellcolor{green!15}0.261 & \cellcolor{green!15}$+$0.063 & \cellcolor{green!15}0.073 & \cellcolor{green!15}$+$0.074 & \cellcolor{green!15}0.085 \\
  & 60 & 0.439 & 0.351 & 0.507 & -0.126 & 0.418 & 0.439 & 0.152 & 0.000 & 0.127 & 0.227 & 0.174 & \cellcolor{green!15}$+$0.079 & \cellcolor{green!15}0.091 & $-$0.048 & 0.000 & $-$0.038 & 0.000 \\
\bottomrule
\multicolumn{19}{l}{\colorbox{green!15}{\phantom{X}} = that estimator's $\RMC^{\pm}>0$.}
\end{tabular}%
}
\end{table*}

\paragraph{Case D diagnostic: separability-only lower bound on grouped 3-Colouring.}
Case D removes the witness gate as well as shortcut subtraction, and therefore shows how much balanced colourable/non-colourable discrimination remains before constructive verification is imposed.
Under this diagnostic, the picture broadens substantially: gpt-oss-20b and Llama-Nemotron-Super-49B v1.5 (reasoning-on) remain positive throughout the full reported range, Qwen3-14B, Qwen3-32B, and QwQ-32B remain positive through $N=50$, and Nemotron-H-47B-Reasoning-128K (reasoning-on) remains weakly positive through $N=50$ before reaching $0$ at $N=60$.
The contrast between Case D and Cases A--C shows that, on the current grouped 3-Colouring benchmark, the principal bottleneck is constructive witness certification rather than complete collapse of balanced class separability.

\begin{table*}[ht]
\centering
\caption{Open-weight 3-Colouring per-$N$ diagnostics under Case D (separability only).
Model abbreviations: Q14B = Qwen3-14B (think); Q32B = Qwen3-32B (think); QwQ = QwQ-32B (think); gpt-oss = gpt-oss-20b (reason); NH-47B = Nemotron-H-47B-Reasoning-128K (reasoning-on); LN-49B = Llama-3.3-Nemotron-Super-49B v1.5 (reasoning-on).}
\label{tab:gc_os_pern_caseD}
\renewcommand{\arraystretch}{1.10}
\resizebox{\textwidth}{!}{%
\begin{tabular}{p{1.55cm}r rrrr rr rr rrr r}
\toprule
\multirow{2}{*}{Model} & \multirow{2}{*}{$N$}
  & \multirow{2}{*}{Acc}
  & \multicolumn{2}{c}{F1 (per class)}
  & \multirow{2}{*}{MCC}
  & \multirow{2}{*}{G-mean}
  & \multirow{2}{*}{BalAcc}
  & \multirow{2}{*}{ADR} & \multirow{2}{*}{ADR$^{+a}$}
  & \multirow{2}{*}{$\hat{s}$}
  & \multicolumn{2}{c}{Separability LCB}
  & \multirow{2}{*}{$\RMC_D$} \\
\cmidrule(lr){4-5}\cmidrule(lr){12-13}
 & & & F1$_{C}$ & F1$_{N}$ & MCC & & & & & &
  LCB$_{CP}$ & LCB$_{W}$ & \\
\midrule
\multirow{8}{*}{\makecell[l]{\textbf{Q14B}\\\textbf{(think)}}}
  & 8 & 0.844 & 0.816 & 0.865 & 0.722 & 0.830 & 0.844 & 0.688 & 0.686 & 0.689 & 0.645 & 0.645 & \cellcolor{green!15}$+$0.395 \\
  & 10 & 0.755 & 0.680 & 0.802 & 0.585 & 0.717 & 0.757 & 0.517 & 0.507 & 0.515 & 0.460 & 0.461 & \cellcolor{green!15}$+$0.210 \\
  & 15 & 0.565 & 0.455 & 0.638 & 0.144 & 0.528 & 0.566 & 0.288 & 0.116 & 0.279 & 0.237 & 0.238 & $-$0.013 \\
  & 20 & 0.519 & 0.604 & 0.386 & 0.042 & 0.472 & 0.519 & 0.230 & 0.000 & 0.223 & 0.178 & 0.179 & $-$0.072 \\
  & 30 & 0.505 & 0.661 & 0.080 & 0.026 & 0.204 & 0.505 & 0.038 & 0.000 & 0.042 & 0.024 & 0.025 & $-$0.226 \\
  & 40 & 0.499 & 0.656 & 0.079 & -0.006 & 0.202 & 0.499 & 0.035 & 0.000 & 0.041 & 0.023 & 0.025 & $-$0.227 \\
  & 50 & 0.539 & 0.693 & 0.072 & 0.015 & 0.194 & 0.503 & 0.039 & 0.000 & 0.038 & 0.007 & 0.012 & $-$0.243 \\
  & 60 & -- & -- & -- & -- & -- & -- & -- & -- & -- & -- & -- & -- \\
\midrule
\multirow{8}{*}{\makecell[l]{\textbf{Q32B}\\\textbf{(think)}}}
  & 8 & 0.816 & 0.777 & 0.843 & 0.673 & 0.797 & 0.816 & 0.634 & 0.628 & 0.635 & 0.571 & 0.572 & \cellcolor{green!15}$+$0.321 \\
  & 10 & 0.692 & 0.564 & 0.762 & 0.474 & 0.627 & 0.692 & 0.389 & 0.369 & 0.393 & 0.316 & 0.319 & \cellcolor{green!15}$+$0.066 \\
  & 15 & 0.567 & 0.349 & 0.675 & 0.193 & 0.457 & 0.571 & 0.210 & 0.129 & 0.209 & 0.162 & 0.164 & $-$0.088 \\
  & 20 & 0.546 & 0.467 & 0.604 & 0.100 & 0.526 & 0.547 & 0.278 & 0.052 & 0.276 & 0.205 & 0.208 & $-$0.045 \\
  & 30 & 0.497 & 0.586 & 0.361 & -0.006 & 0.449 & 0.497 & 0.214 & 0.000 & 0.202 & 0.141 & 0.144 & $-$0.109 \\
  & 40 & 0.508 & 0.668 & 0.052 & 0.059 & 0.163 & 0.508 & 0.027 & 0.000 & 0.026 & 0.008 & 0.011 & $-$0.242 \\
  & 50 & 0.583 & 0.722 & 0.167 & 0.227 & 0.302 & 0.545 & 0.091 & 0.000 & 0.091 & 0.017 & 0.029 & $-$0.233 \\
  & 60 & -- & -- & -- & -- & -- & -- & -- & -- & -- & -- & -- & -- \\
\midrule
\multirow{8}{*}{\makecell[l]{\textbf{QwQ}\\\textbf{(think)}}}
  & 8 & 0.917 & 0.910 & 0.924 & 0.846 & 0.913 & 0.917 & 0.834 & 0.824 & 0.834 & 0.777 & 0.778 & \cellcolor{green!15}$+$0.527 \\
  & 10 & 0.835 & 0.802 & 0.858 & 0.709 & 0.818 & 0.835 & 0.669 & 0.669 & 0.669 & 0.561 & 0.564 & \cellcolor{green!15}$+$0.311 \\
  & 15 & 0.599 & 0.343 & 0.711 & 0.325 & 0.455 & 0.601 & 0.211 & 0.201 & 0.207 & 0.160 & 0.162 & $-$0.090 \\
  & 20 & 0.550 & 0.180 & 0.689 & 0.228 & 0.315 & 0.550 & 0.099 & 0.091 & 0.099 & 0.051 & 0.056 & $-$0.199 \\
  & 30 & 0.504 & 0.189 & 0.643 & 0.013 & 0.321 & 0.504 & 0.099 & 0.008 & 0.103 & 0.053 & 0.058 & $-$0.197 \\
  & 40 & 0.546 & 0.525 & 0.565 & 0.097 & 0.545 & 0.548 & 0.282 & 0.000 & 0.297 & 0.201 & 0.205 & $-$0.049 \\
  & 50 & 0.500 & 0.645 & 0.154 & 0.000 & 0.287 & 0.500 & 0.091 & 0.000 & 0.083 & 0.001 & 0.010 & $-$0.249 \\
  & 60 & -- & -- & -- & -- & -- & -- & -- & -- & -- & -- & -- & -- \\
\midrule
\multirow{8}{*}{\makecell[l]{\textbf{gpt-oss}\\\textbf{(reason)}}}
  & 8 & 0.871 & 0.861 & 0.879 & 0.752 & 0.868 & 0.872 & 0.786 & 0.762 & 0.753 & 0.536 & 0.546 & \cellcolor{green!15}$+$0.286 \\
  & 10 & 0.849 & 0.831 & 0.863 & 0.723 & 0.843 & 0.852 & 0.714 & 0.714 & 0.710 & 0.500 & 0.510 & \cellcolor{green!15}$+$0.250 \\
  & 15 & 0.547 & 0.264 & 0.672 & 0.182 & 0.389 & 0.556 & 0.143 & 0.119 & 0.152 & 0.056 & 0.067 & $-$0.194 \\
  & 20 & 0.643 & 0.444 & 0.737 & 0.372 & 0.535 & 0.635 & 0.300 & 0.300 & 0.286 & 0.141 & 0.155 & $-$0.109 \\
  & 30 & 0.506 & 0.050 & 0.667 & -0.059 & 0.160 & 0.489 & 0.028 & 0.000 & 0.026 & 0.001 & 0.004 & $-$0.249 \\
  & 40 & 0.537 & 0.194 & 0.675 & 0.076 & 0.327 & 0.522 & 0.048 & 0.000 & 0.107 & 0.019 & 0.031 & $-$0.231 \\
  & 50 & 0.500 & 0.162 & 0.644 & -0.072 & 0.296 & 0.476 & 0.069 & 0.000 & 0.088 & 0.015 & 0.025 & $-$0.235 \\
  & 60 & 0.571 & 0.308 & 0.690 & 0.197 & 0.426 & 0.566 & 0.133 & 0.000 & 0.181 & 0.059 & 0.073 & $-$0.191 \\
\midrule
\multirow{8}{*}{\makecell[l]{\textbf{NH-47B}\\\textbf{(r-on)}}}
  & 8 & 0.658 & 0.575 & 0.713 & 0.342 & 0.629 & 0.658 & 0.318 & 0.003 & 0.395 & 0.372 & 0.373 & \cellcolor{green!15}$+$0.122 \\
  & 10 & 0.629 & 0.626 & 0.632 & 0.258 & 0.629 & 0.629 & 0.272 & 0.002 & 0.396 & 0.372 & 0.373 & \cellcolor{green!15}$+$0.122 \\
  & 15 & 0.584 & 0.502 & 0.642 & 0.177 & 0.560 & 0.584 & 0.190 & 0.000 & 0.314 & 0.293 & 0.293 & \cellcolor{green!15}$+$0.043 \\
  & 20 & 0.560 & 0.422 & 0.645 & 0.137 & 0.507 & 0.560 & 0.154 & 0.000 & 0.257 & 0.238 & 0.238 & $-$0.012 \\
  & 30 & 0.527 & 0.304 & 0.642 & 0.071 & 0.418 & 0.527 & 0.106 & 0.000 & 0.175 & 0.159 & 0.159 & $-$0.091 \\
  & 40 & 0.504 & 0.060 & 0.663 & 0.025 & 0.176 & 0.504 & 0.021 & 0.000 & 0.031 & 0.025 & 0.025 & $-$0.225 \\
  & 50 & 0.504 & 0.040 & 0.666 & 0.033 & 0.142 & 0.504 & 0.012 & 0.000 & 0.020 & 0.015 & 0.015 & $-$0.235 \\
  & 60 & 0.500 & -- & 0.667 & -- & 0.000 & 0.500 & 0.000 & 0.000 & 0.000 & 0.000 & 0.000 & $-$0.250 \\
\midrule
\multirow{8}{*}{\makecell[l]{\textbf{LN-49B}\\\textbf{(r-on)}}}
  & 8 & 0.943 & 0.940 & 0.946 & 0.892 & 0.941 & 0.943 & 0.886 & 0.886 & 0.886 & 0.694 & 0.699 & \cellcolor{green!15}$+$0.444 \\
  & 10 & 0.818 & 0.778 & 0.846 & 0.683 & 0.798 & 0.818 & 0.636 & 0.636 & 0.636 & 0.439 & 0.449 & \cellcolor{green!15}$+$0.189 \\
  & 15 & 0.557 & 0.235 & 0.688 & 0.210 & 0.365 & 0.557 & 0.136 & 0.136 & 0.133 & 0.046 & 0.056 & $-$0.204 \\
  & 20 & 0.545 & 0.167 & 0.688 & 0.218 & 0.302 & 0.545 & 0.091 & 0.091 & 0.091 & 0.023 & 0.033 & $-$0.227 \\
  & 30 & 0.477 & 0.148 & 0.623 & -0.072 & 0.280 & 0.477 & 0.068 & 0.000 & 0.079 & 0.018 & 0.026 & $-$0.232 \\
  & 40 & 0.523 & 0.400 & 0.604 & 0.050 & 0.481 & 0.523 & 0.205 & 0.000 & 0.231 & 0.106 & 0.116 & $-$0.144 \\
  & 50 & 0.598 & 0.598 & 0.598 & 0.196 & 0.598 & 0.598 & 0.349 & 0.000 & 0.357 & 0.192 & 0.202 & $-$0.058 \\
  & 60 & 0.439 & 0.351 & 0.507 & -0.126 & 0.418 & 0.439 & 0.152 & 0.000 & 0.174 & 0.061 & 0.071 & $-$0.189 \\
\bottomrule
\multicolumn{14}{l}{\colorbox{green!15}{\phantom{X}} = $\RMC_D>0$.}
\end{tabular}%
}
\end{table*}

\subsection{RQ2 synthesis: Certified capability boundaries}
\label{sec:phase}

\Cref{tab:com_pern,tab:os_pern} already present the full per-$N$ breakdown.
Here we highlight the cross-model patterns.

\paragraph{Finding 7: Sharp phase transition is universal across all certified models.}
All six certified models transition from certified to uncertified within one or two tested $N$ steps under the Case A pooled CP criterion that defines $N^*$:
GPT-5 between $N=30$ ($+0.117$) and $N=40$ ($-0.040$);
Qwen3-14B between $N=20$ ($+0.015$) and $N=25$ ($-0.086$);
Qwen3-32B and QwQ-32B between $N=15$ and $N=20$;
Opus 4.7 between $N=25$ and the next tested size $N=50$, and DeepSeek-Reasoner between $N=10$ and $N=25$ (the tested grid is coarser for these two models).
Under the recalibrated residual Case B ceiling, GPT-5 still shows an earlier Case B cliff ($N=25\to30$), whereas the three certified open-source 2-SAT models now lose their Case B certificates on the same tested step as their Case A pooled-CP boundaries: Qwen3-14B at $N=20\to25$ and Qwen3-32B/QwQ-32B at $N=15\to20$.
Under the construction-level-blended Case C ceiling (half of the permutation-calibrated residualized max-KS estimate, with the other half anchored at Case A's $\delta=0$), the residual $\shortcutCeiling$ lands in the $0.035$--$0.091$ range on 3-SAT and $0.10$--$0.13$ on 2-SAT, materially smaller than raw Case B but high enough to bite at large $N$.
Under the global-pooled CP estimator, GPT-5 and Opus 4.7 (medium thinking) keep a positive Case C margin at every tested $N$; DeepSeek-Reasoner stays certified through $N=25$ ($+0.110$) and fails at $N=30$ ($-0.041$); Qwen3-14B holds through $N=25$ ($+0.030$) and fails at $N=30$; Qwen3-32B and QwQ-32B hold through $N=20$ ($+0.042$ and $+0.066$) and fail at $N=25$.
This transition is invisible in standard metrics: GPT-5's accuracy at $N=40$ is still $0.645$ and ADR is $0.309$ (above chance), yet RMC(CP) is already $-0.040$; at $N=30$, accuracy is still $0.735$ while the calibrated Case B certificate has already failed.

\paragraph{Finding 8: RMC is slightly below ADR$^{+w}$ and the gap grows with problem size.}
Across the certified SAT models, RMC is generally below ADR$^{+w}$, which is the expected relationship if RMC discounts shortcut-driven, prior-driven, and uncertain category success instead of reproducing the witness-aware pair rate.  The gap also tends to widen as $N$ increases.  This pattern is informative.  As formulas become longer, the model's understanding of the CNF structure weakens on both sides of the contrast.  The SAT side is partially audited by assignment checking through ADR$^{+w}$, but the UNSAT side is harder: a correct UNSAT label gives weaker evidence of identified contradiction or refutation structure.  RMC penalizes this loss through the factorized positive/negative separability term and the lower-confidence bound.  Thus, an increasing ADR$^{+w}$--RMC gap at larger $N$ is consistent with a growing number of category-level successes that fall below the stability, balance, and shortcut-robustness required for certified reasoning evidence.

\paragraph{Finding 9: DeepSeek-Reasoner has the steepest capability cliff among commercial models.}
DeepSeek-Reasoner achieves CP $+0.642$ at $N=5$ (close to GPT-5's $+0.679$) but drops sharply to $+0.391$ at $N=10$ and goes negative by $N=25$ ($-0.053$).
GPT-5 maintains CP $+0.647$ at $N=10$ and $+0.362$ at $N=25$, giving a much wider certified range within the commercial 3-SAT family; among the open-weight 2-SAT models, Qwen3-14B similarly holds $+0.344$ at $N=15$, though 2-SAT and 3-SAT boundaries are not directly comparable.
The RMC(no LCB) estimator is deceptively optimistic: DeepSeek-Reasoner's RMC(no LCB) at $N=25$ is still $+0.043$ while both RMC(CP) ($-0.053$) and RMC(Wilson) ($-0.047$) already flag the model as uncertified, underscoring the risk of reporting raw separability without statistical discounting.

\subsection{RQ1 synthesis: Why label metrics fail mechanistically}
\label{sec:bias}

The decomposition of $\shat = \thetaP\times\thetaN$ reveals the mechanism behind each model's failure or success:

\begin{itemize}[leftmargin=*, itemsep=2pt]
\item \textbf{All-SAT bias} (GPT-4o-mini, GPT-5.4-mini, Claude Haiku 4.5, Opus 4.7 std/direct/Sonnet, Gemma-4-31B, LLaMA-3.1, Qwen2.5-32B, Qwen3-30B): $\thetaP\approx 1.0$, $\thetaN\approx 0.0$.
These models predict SAT for almost every instance.
They achieve $\approx50\%$ instance accuracy on balanced data, ADR $\approx 0.0$, and $\shat\approx 0.0$.
The RMC discrimination margin is essentially at the floor ($-0.250$).
This pattern is consistent across both commercial and open-source families.

\item \textbf{All-UNSAT bias} (DeepSeek-Chat): $\thetaP=0.074$, $\thetaN=0.958$.
The symmetric pattern: predicts UNSAT almost universally.
ADR $= 0.058$ appears non-trivial, but $\LCB_{\mathrm{fact}}(\shat)=0.009$ is far below the $0.25$ ceiling.
The witness rate is $0.000$ (no valid assignments for any of the few SAT predictions).

\item \textbf{Partial discriminators, not certified} (DeepSeek-R1-Distill-32B): $\thetaP\approx0.17$, $\thetaN\approx0.88$.
A mixed pattern with UNSAT lean; $\LCB_{\mathrm{fact}}<0.25$ at all $N$ due to weak positive-side accuracy.

\item \textbf{Certified partial discriminators} (GPT-5, Opus med.\ think, DeepSeek-Reasoner, Qwen3-14B, Qwen3-32B, QwQ-32B): $\thetaP\in[0.35, 0.95]$, $\thetaN\in[0.61, 1.00]$.
These models genuinely discriminate at small $N$ but degrade as $N$ grows, consistent with hitting the limits of their SAT-solving capability.
The certified open-source models (Qwen3/QwQ) tend toward stronger UNSAT detection ($\thetaN\approx 1$) with weaker positive detection compared to GPT-5.
\end{itemize}

The decomposition highlights the value of the product statistic $\shat=\thetaP\thetaN$: even DeepSeek-Chat ($\thetaN=0.96$) scores $\shat=0.059$ because $\thetaP=0.07$; even Qwen3-32B ($\thetaN=0.958$) loses certification at $N=20$ because $\thetaP$ falls below the level needed to push $\LCB_{\mathrm{fact}}$ above $0.25$.
Any metric based solely on per-class accuracy (precision/recall) would miss this compensatory degradation.

\subsection{The Shortcut-Ceiling Spectrum in Practice}
\label{sec:spectrum}

\Cref{tab:spectrum} demonstrates the shortcut-ceiling spectrum empirically on a representative shortcut model (GPT-4o-mini, always-SAT).

\begin{table}[ht]
\centering
\caption{Same model (GPT-4o-mini, always-SAT bias), different evaluation protocols.
The shortcut ceiling governs what a shortcut model can score on each protocol.
Only within-group separability collapses the ceiling to $0.25$ by construction.}
\label{tab:spectrum}
\renewcommand{\arraystretch}{1.2}
\small
\begin{tabularx}{\columnwidth}{p{1.8cm}p{1.6cm}p{2.1cm}X}
\toprule
Protocol & GPT-4o-mini score & Shortcut ceiling & Gap from ceiling \\
\midrule
Instance accuracy & 0.498 & $\rho\approx0.50$ (balanced) & $\approx0.000$ \\
F1 (SAT-positive) & 0.666 & $\ge2\pi/(1+\pi)=0.667$ & $\approx0.000$ \\
ADR (pair-level) & 0.000 & $\le0.25$ (group structure) & $\approx0.000$ \\
Within-group sep.\ ($\shat$) & 0.000 & $0.25$ & $-0.250$ \\
$\RMCSAT^{\pm}$ & $-0.250$ & $0$ (target) & $-0.250$ \\
\bottomrule
\end{tabularx}
\end{table}

The table shows that accuracy ($0.498$) and F1 ($0.666$) are both near their respective shortcut ceilings---the model is ``as good as it should be'' given only a label-bias shortcut.
The RMC framework makes this explicit: the model is certified to contribute zero certified ability ($\RMCSAT^{+}=0.000$).

\subsection{RQ3: Generalization to vulnerability detection}
\label{sec:vd_results}

\Cref{tab:vd} presents global-aggregate $\RMCVD$ results for representative VD model configurations, covering all three certification regimes observed across 49 model/scenario pairs.


\begin{table*}[t]
\centering
\caption{Vulnerability detection RMC results (global aggregate factorized LCB, $\alpha=0.05$, $\shortcutDisc=0$, no witness gate).
$\theta_P$ = fraction of vulnerable instances correctly predicted vulnerable;
$\theta_N$ = fraction of non-vulnerable instances correctly predicted non-vulnerable;
$\LCB_{\mathrm{fact}}$ = global factorized CP lower bound;
Disc.\ = discrimination margin ($= \LCB_{\mathrm{fact}} - 0.25$);
ADR = Accurate Differentiation Rate;
$\RMCVD^{\pm}$ = signed VD RMC.
Scen.: 1p/g = 1,055 pairs (1 pair/group); 10p/g = 10,550 pairs (10 pairs/group, augmented).
Shaded rows: certified ($\RMCVD^{\pm}>0$).}
\label{tab:vd}
\renewcommand{\arraystretch}{1.15}
\resizebox{\textwidth}{!}{%
\begin{tabular}{llrrrrrr}
\toprule
Model & Scen. & $\theta_P$ & $\theta_N$ & $\LCB_{\mathrm{fact}}$ & Disc.\ margin & ADR & $\RMCVD^{\pm}$ \\
\midrule
\multicolumn{8}{l}{\textit{Balanced training (\texttt{add\_after\_ratio}, \texttt{add\_after\_percent\_70}, \texttt{add\_after\_10epoch}): all 33 pairs certified}} \\
\rowcolor{green!10}\texttt{add\_after\_ratio\_2\_v04} & 10p/g & 0.753 & 0.613 & 0.450 & $+$0.200 & 0.405 & $\mathbf{+0.200}$ \\
\rowcolor{green!10}\texttt{add\_after\_ratio\_1}     & 10p/g & 0.768 & 0.598 & 0.447 & $+$0.197 & 0.406 & $\mathbf{+0.197}$ \\
\rowcolor{green!10}\texttt{add\_after\_ratio\_1\_v01}& 1p/g  & 0.754 & 0.645 & 0.447 & $+$0.197 & 0.447 & $\mathbf{+0.197}$ \\
\rowcolor{green!10}\texttt{add\_after\_percent\_70}  & 10p/g & 0.822 & 0.512 & 0.409 & $+$0.159 & 0.373 & $\mathbf{+0.159}$ \\
\rowcolor{green!10}\texttt{add\_after\_10epoch}      & 1p/g  & 0.718 & 0.609 & 0.399 & $+$0.149 & 0.359 & $\mathbf{+0.149}$ \\
\rowcolor{green!10}\texttt{add\_after\_ratio\_1\_v03}& 1p/g  & 0.852 & 0.418 & 0.322 & $+$0.072 & 0.287 & $\mathbf{+0.072}$ \\
\midrule
\multicolumn{8}{l}{\textit{Imbalanced training (\texttt{add\_after\_percent\_50}): not certified (marginal)}} \\
\texttt{add\_after\_percent\_50} & 10p/g & 0.913 & 0.278 & 0.245 & $-$0.005 & 0.201 & $-0.006$ \\
\texttt{add\_after\_percent\_50} & 1p/g  & 0.918 & 0.281 & 0.228 & $-$0.022 & 0.202 & $-0.022$ \\
\midrule
\multicolumn{8}{l}{\textit{Vulnerable-only training (\texttt{func\_before}, \texttt{func\_before\_1000\_2000}): all 14 pairs not certified}} \\
\texttt{func\_before}            & 10p/g & 0.946 & 0.079 & 0.069 & $-$0.181 & 0.033 & $-0.181$ \\
\texttt{func\_before\_1000\_2000}& 10p/g & 0.931 & 0.078 & 0.067 & $-$0.183 & 0.022 & $-0.183$ \\
\texttt{func\_before}            & 1p/g  & 0.946 & 0.054 & 0.038 & $-$0.212 & 0.011 & $-0.212$ \\
\texttt{func\_before\_1000\_2000}& 1p/g  & 0.926 & 0.058 & 0.040 & $-$0.210 & 0.004 & $-0.210$ \\
\bottomrule
\multicolumn{8}{l}{\cellcolor{green!10}\phantom{X} = Certified.}
\end{tabular}%
}
\end{table*}

\paragraph{Finding 10: Training data balance is the primary determinant of VD certification.}
The most striking aspect of the VD results is their predictive clarity: certification status is almost entirely determined by whether the training data includes non-vulnerable (\texttt{func\_after}) examples alongside vulnerable ones.
All 33 model/scenario pairs from the \texttt{add\_after\_ratio\_*}, \texttt{add\_after\_percent\_70}, and \texttt{add\_after\_10epoch} families are certified ($\RMCVD^{\pm}\in[+0.072, +0.200]$).
All 14 \texttt{func\_before\_*} pairs fail certification ($\RMCVD^{\pm}\in[-0.222, -0.170]$).
The \texttt{add\_after\_percent\_50} family is the sole ambiguous case: two configurations narrow miss certification at $-0.022$ and $-0.006$.

\paragraph{Finding 11: Vulnerable-only training induces all-vulnerable bias, mirroring SAT's all-SAT bias.}
The \texttt{func\_before} variants achieve $\theta_P\in[0.91, 0.95]$ (very high vulnerable-instance recall) but $\theta_N\in[0.04, 0.09]$ (near-zero non-vulnerable recall): they predict ``vulnerable'' for almost every input regardless of semantics.
This is the exact VD analog of the all-SAT bias observed in 18 of 24 SAT-domain configurations: a model that ignores semantics and exploits label-frequency shortcuts.
Instance accuracy for these models is $0.60$--$0.65$, and ADR is $0.004$--$0.033$---but standard metrics leave the direction of the bias uncertified.  When the evaluation distribution is imbalanced, the accuracy rows should also be read against the majority-class baseline in Eq.~\eqref{eq:acc-baseline-supp}; otherwise a high raw accuracy can simply reflect the class prior.
RMC does: the product statistic $\theta_P\cdot\theta_N$ collapses to near zero when one factor dominates the other, regardless of which side is inflated.

\paragraph{Finding 12: Global aggregate RMC is consistent across augmentation scenarios.}
The 1 pair/group and 10 pairs/group scenarios produce nearly identical $\RMCVD^{\pm}$ values for the same model (e.g., \texttt{add\_after\_ratio\_1\_v01}: $+0.197$ at 1p/g vs.\ $+0.197$ at 10p/g).
The primary benefit of the augmented 10p/g scenario is a tighter LCB on each component---a useful property when benchmarks are designed with multiple augmentation variants per base instance---but the certified/uncertified classification is unchanged.

% ============================================================
\section{Discussion}
\label{sec:discussion}
% ============================================================

\subsection{Scope of the RMC Certificate}

A positive $\RMCSAT^{\pm}$ margin certifies:

\begin{quote}
\emph{Under the specified group construction, the specified side-channel adversary class $\shortcutHypCls$, and the chosen confidence budget $\alpha$, the model's within-group category discrimination exceeds the specified side-channel explanation under the stated confidence budget.}
\end{quote}

It does \emph{not} certify that the model reasons like a SAT solver in any philosophical sense, nor that it will generalize to other problem families.
Certification is explicitly relative to $\shortcutHypCls$ (the adversary battery) and the particular minimal-edit construction.
Enlarging $\shortcutHypCls$ (running more side-channel classifiers) only raises $\hatshortcutDisc$, which raises the shortcut ceiling and makes certification \emph{harder}---a property that should be exploited in practice to obtain progressively tighter certificates.

\subsection{The Certified Capability Boundary as a New Evaluation Artifact}

The certified capability boundary $N^*$, defined as the largest $N$ at which $\RMCSAT^{\pm}>0$, is a natural single-number summary that differs fundamentally from accuracy at a fixed $N$.  Across the 24 SAT configurations, GPT-5 has the strongest commercial boundary ($N^*=25$ under the conservative per-group criterion).  Qwen3-14B reaches $N^*=15$ and is the strongest open-weight model in our experiments.  Qwen3-32B, QwQ-32B, Claude Opus 4.7 medium-thinking, and DeepSeek-Reasoner reach $N^*=10$.  The remaining 18 configurations have $N^*=0$, meaning that none is certified at any tested size.

This ranking diverges substantially from both accuracy and ADR rankings.
Most strikingly, Qwen3-14B (14B parameters, open-weight) ranks second overall with $N^*=15$, ahead of the commercial Anthropic and DeepSeek flagship reasoning models (both $N^*=10$).
Qwen3-32B and QwQ-32B tie with the certified commercial models at $N^*=10$.
Parameter scale alone gives an incomplete explanation of the result: QwQ-32B (32B parameters) achieves the same $N^*=10$ as Qwen3-32B despite more parameters, while Qwen3-14B ($N^*=15$) outperforms both with fewer parameters.
The determinant appears to be the quality of the SAT-positive discrimination gate ($\thetaP$): Qwen3-14B's $\thetaP=0.568$ is more balanced than Qwen3-32B's $\thetaP=0.464$ or QwQ-32B's $\thetaP=0.471$, despite all three having near-perfect $\thetaN\approx 1$.

\subsection{Implications for Benchmark Design}

\paragraph{Minimum group size.}
The Clopper--Pearson LCB with Bonferroni split requires a non-trivial $|P_{\mathrm{eval}}|$ and $|N_{\mathrm{eval}}|$ to produce a useful lower bound.
For $\alpha'=0.025$, with a true $\theta_P=0.9$, $|P_{\mathrm{eval}}|=10$ yields LCB$\approx0.62$; $|P_{\mathrm{eval}}|=2$ yields LCB$\approx0.26$---barely above zero.
This implies that benchmarks with fewer than $\sim10$ instances per side per group are too small to reliably certify genuine ability.

\paragraph{The pseudo-replication trap.}
The gap between the correct factorized LCB and the naive single-binomial LCB grows with group size imbalance.
For a group with $|P|=|N|=10$ and $\shat=0.8$ (realistic for GPT-5 at $N=5$), the single-binomial LCB on $k=80$, $m=100$ is approximately $0.71$ (severely overestimated), while the factorized LCB is approximately $0.51\times0.51\approx0.26$ (corrected).
The overestimation by a factor of $\sim2.7$ would lead to spurious certifications.

% ============================================================
\section{Threats to Validity}
\label{sec:threats}
% ============================================================


\paragraph{Case A, Case B, Case C, and Case D interpretation.}
The Case A results use the ideal-decoupled-limit ceiling $\shortcutCeiling=0.25$.
This setting asks what can be certified when the positive and negative arms within each minimal-edit group are treated as sufficiently matched with respect to the side-channel features controlled by the design.
The Case B results then replace this fixed ceiling with an estimated ceiling from the reported CNF side-channel adversary battery.
Case C keeps the same residual ceiling transform as Case B but estimates the residual distinguishability after \emph{feature-level} residualization against generic structural-complexity covariates (Frisch--Waugh--Lovell partial-out): each feature in the final Case-B residual battery is first regressed on the raw linear covariates (num\_vars, num\_clauses, num\_literals, file\_size\_bytes), the per-instance residual replaces the raw feature value, the residual max-KS adversary is permutation-calibrated, and the resulting residualized estimate is then blended 50/50 with the Case A reference $\delta=0$.
Case D removes both of those adjustments: it sets the shortcut ceiling to zero, removes the witness min-gate, and reports the factorized separability lower bound itself.  Thus, the paper assigns the four treatments fixed roles: Case A gives the clean construction-level reference point that defines $N^*$, Case B---the primary shortcut-adjusted certificate in the main paper---reports how the certificates tighten when the stated residual shortcut class is explicitly estimated from the final residual battery, Case C is a supplementary robustness check that keeps that same battery but removes the complexity-predictable component of the estimated leakage before recomputing the ceiling, and Case D is a separability-only upper-envelope diagnostic that intentionally does \emph{not} charge any shortcut account.  The reported Case B ceiling uses the permutation-calibrated point estimate of $\hatshortcutDisc$ (Remark~\ref{rem:caseb-variants}); the strict UCB-inflated variant is uniformly harder to pass, so negative Case B margins are robust to it while positive margins could shrink.
All four interpretations are relative to their assumptions and should not be read as absolute proof that no unmeasured side-channel exists.

\paragraph{Certification relative to $\shortcutHypCls$.}
The certificate is valid only against the shortcut classes in $\shortcutHypCls$.
A clever side-channel not in the battery, such as a syntactic fingerprint specific to a particular CNF generator, could still reduce the strength of the certificate.
We partially mitigate this through our minimal-edit construction design and the calibrated Case B adversary battery, but broader shortcut batteries would provide stronger future certificates.

\paragraph{UNSAT witness asymmetry.}
ADR$^{+w}$ audits only the SAT side with a machine-checkable assignment.  This is useful because satisfying assignments are efficiently verifiable, but it is also asymmetric: UNSAT proofs or refutation certificates are substantially harder to elicit and verify from LLM outputs.  For this reason, ADR$^{+w}$ should be interpreted as a generous SAT-side witness-aware proxy rather than a full proof-aware SAT/UNSAT reasoning metric.  RMC handles this limitation differently.  It estimates shortcut-adjusted category discrimination over matched SAT and UNSAT arms from final categories.  The comparison between RMC and ADR$^{+w}$ is therefore diagnostic rather than definitional.  When RMC is below ADR$^{+w}$, especially at larger $N$, a plausible explanation is that some category successes, particularly on the UNSAT side, fall below the paired, shortcut-adjusted, lower-confidence test.

\paragraph{Independence and clustering.}
The pooled Clopper--Pearson and Wilson bounds in \Cref{sec:table_estimators} and the VD global aggregate bound in \Cref{sec:vd_rmc} treat instance-level outcomes as independent binomial draws.
Two clustering sources challenge this assumption: instances are generated within seed families (minimal-edit groups), and all instances share a fixed zero-shot prompt template.
The prompt template is a fixed evaluation condition rather than a sampled factor, so it defines the scope of the claim rather than violating exchangeability; seed-family clustering, however, can induce positive intra-cluster correlation, under which pooled intervals may be anti-conservative.
Note that this threat is distinct from the pseudo-replication issue addressed by the factorized statistic (\Cref{sec:lcb}), which concerns reusing observations across cross-arm comparisons rather than correlation within seed families.
The group remains the declared inference unit, and the per-group macro diagnostics provide a clustering-respecting cross-check, but a group-level (cluster) bootstrap or cluster-robust interval is the appropriate formal robustness check; we flag this assumption explicitly and recommend cluster-level intervals for future RMC reports.

\paragraph{Prompt sensitivity.}
Results may depend on the specific zero-shot prompt formulation.
Models evaluated with different prompts, such as chain-of-thought elicitation or few-shot examples, may exhibit different certification boundaries.
We evaluate only zero-shot descriptive prompts in this work.

\paragraph{Data limitations.}
Some configurations, such as Claude Opus 4.7 full thinking, have very few groups ($n_{\mathrm{groups}}=2$), making RMC trivially uncertifiable due to the conservatism of the lower confidence bound.
These entries are included in tables for completeness but should not be used for model comparisons.

% ============================================================

\paragraph{Conventional-metric baseline under class imbalance.}
A final measurement threat concerns ordinary accuracy on imbalanced evaluation sets.  Raw accuracy can exceed $0.5$ by a large margin even for an always-majority predictor.  We therefore use the majority-class baseline $\accBase$ when interpreting accuracy, while leaving RMC unchanged because its arm-normalized product already penalizes one-sided behavior.  The remaining effect of imbalance is uncertainty: smaller arms produce weaker confidence bounds.

\section{Related Work}
\label{sec:related}
% ============================================================

\paragraph{Shortcut learning and spurious correlations.}
The connection between high benchmark accuracy and shallow pattern matching is well-documented in machine learning and NLP~\cite{geirhos2020shortcut,gururangan2018annotation,mccoy2019right,niven2019probing}.
Our framework provides the first quantitative certification of whether a model's discrimination is \emph{certifiably} above the shortcut ceiling in the combinatorial reasoning domain.

\paragraph{Metamorphic testing and differential testing.}
Metamorphic testing and test-oracle research~\cite{chen2020metamorphic,barr2015oracle} as well as behavioral, contrastive, and counterfactual evaluation~\cite{ribeiro2020checklist,gardner2020contrast,kaushik2020counterfactual} share the spirit of our minimal-edit groups, while RMC adds shortcut-adjusted statistical certification with a controlled false-positive rate.

\paragraph{Causal identifiability and interventional evaluation.}
Pearl's framework for causal identifiability~\cite{pearl2009causality} and invariant risk minimization~\cite{arjovsky2019invariant} motivate our interventional framing.
RMC operationalizes identifiability for finite-sample statistical certification.

\paragraph{Underspecification and predictive multiplicity.}
D'Amour et al.'s underspecification framework~\cite{damour2022underspecification} and predictive multiplicity~\cite{marx2020predictive,watson2023predictive} formalize the non-uniqueness of models that agree on training data.
Observation~\ref{obs:nonid} is a direct consequence of underspecification in the evaluation context.

\paragraph{Conformal prediction and calibrated uncertainty.}
Conformal prediction~\cite{vovk2005algorithmic} provides coverage guarantees; our factorized LCB serves an analogous role but for a product statistic.
A conformal-style coverage statement for RMC is an open research direction.

% ============================================================

\section{Conclusion}
\label{sec:conclusion}
% ============================================================

This paper makes two connected arguments.  First, conventional class-label metrics such as accuracy and F1 provide unreliable evidence of reasoning ability.  They can be high when a model exploits label priors, one-sided sensitivity, or benchmark artifacts, and they assign the same score to a genuine solver and to a shortcut, prior-driven, or lucky model that produces the same final labels.  The empirical evidence is consistent across SAT and vulnerability detection: Acc and F1 often remain much larger than ADR$^{+w}$ and RMC precisely when the model fails to provide a witness or fails one side of the matched discrimination test.  ADR$^{+w}$ is therefore used as a SAT-side witness-aware proxy for constructive reasoning evidence and a generous reference point rather than a full proof-aware SAT/UNSAT metric.  Because it omits UNSAT refutations, ADR$^{+w}$ is generous relative to full proof-aware SAT/UNSAT reasoning; the fact that ordinary category metrics still exceed it strengthens the overestimation diagnosis.

Second, RMC provides a statistically disciplined way to ask a more defensible question: whether the model's category behavior exceeds a specified shortcut-only account under an explicit confidence budget.  This is consistent with the critique of class-label metrics.  RMC still observes output correctness, while separating isolated category accuracy from reasoning evidence and explicitly discounting the part of category correctness that can be explained by shortcut or guessing baselines.  Minimal-edit groups define the semantic contrast, the factorized LCB avoids pseudo-replication, and the shortcut ceiling states the adversarial baseline.  With $\alpha=0.05$ in Case A, positive RMC(CP) gives a conservative 95\% confidence certificate that the separability term exceeds the fixed shortcut ceiling.  In Case B, the same logic applies relative to the estimated shortcut ceiling and the chosen adversary battery.  ADR$^{+w}$ remains a separate SAT-side witness-aware reference: it helps interpret whether RMC is conservatively below constructive SAT-side evidence, while the RMC computation uses category outcomes only.

The resulting evaluation changes the object of measurement.  RMC shifts the question from whether a model looks accurate on a benchmark to whether its matched positive and negative behavior exceeds the shortcut explanation under test.  This produces a more appropriate evaluation artifact for reasoning-intensive software engineering tasks: a certified capability boundary rather than a single aggregate accuracy score.

\section{Additional RMC Details}
\label{app:rmc_details}
% ============================================================

\subsection{Case A, Case B, Case C, and Case D Treatment}

The paper reports three ceiling treatments plus one separability-only diagnostic with fixed roles.
\textbf{Case A} treats the shortcut ceiling as a construction-level constant:
for the minimal-edit SAT benchmarks, $\shortcutDisc=0$ and therefore
$\shortcutCeiling=0.25$.
The Case A columns RMC(no LCB), RMC(CP), and RMC(Wilson) in
\Cref{tab:com_pern,tab:os_pern} define the construction-level boundary $N^*$.
\textbf{Case B} estimates the ceiling from the final residual side-channel adversary
battery (\Cref{sec:ceiling_estimation}) and is the primary shortcut-adjusted
certificate reported in the main paper; the reported Case B tables use the
permutation-calibrated point estimate $\hatshortcutDisc$
(Remark~\ref{rem:caseb-variants}).
\textbf{Case C} keeps the same residual ceiling transform as Case B, but
estimates the shortcut distinguishability after a \emph{feature-level}
partialing-out of generic structural-complexity covariates (number of
variables, number of clauses, number of literals, and file size).  Concretely,
for every leakage feature $f_k$ in the final Case-B residual battery, a pooled ridge regression
$f_k\sim 1 + (\text{complexity covariates})$ is fitted pooled across all
instances; the per-instance residual $\tilde f_k$ replaces the raw $f_k$
before the max-KS adversary and permutation calibration are applied.  The implemented design uses only these raw linear covariates, not higher-order or interaction expansions.  The reported Case C estimate is then the default 50/50 blend between the calibrated residualized max-KS estimate and the Case A reference $\delta=0$.  Relative to Case B, this means that Case C keeps the same shortcut battery and the same permutation-calibrated max-KS logic, but charges only the leakage component that survives the complexity partial-out, and then shrinks that component toward the construction-level idealization.  This
is the standard Frisch--Waugh--Lovell construction, applied in feature space
rather than as a scalar post-hoc attenuation of $\hatshortcutDisc$.  Case C is
reported only as a supplementary sensitivity analysis and does not replace
the headline Case B certificate.
\textbf{Case D} removes shortcut subtraction and witness gating entirely.
In the implemented scripts, this is the \texttt{case\_d} mode: for each group
$G$, the reported value is simply
\[
\RMC_D(G)=\LCB_{\alpha/2}(\hat\theta_P(G))\cdot\LCB_{\alpha/2}(\hat\theta_N(G)).
\]
Equivalently, Case D sets $\shortcutCeiling=0$ and identifies
$\RMC_D^{\pm}(G)=\RMC_D^{+}(G)=\RMC_D(G)$.
Because no shortcut account is subtracted, Case D is intentionally easier than
Case A--C and should be read only as a separability-only upper-envelope
diagnostic, not as a shortcut-adjusted certificate.

The strict Case B variant additionally absorbs the sampling uncertainty of the
adversary side.
Let $\hatshortcutDisc$ be the calibrated distinguishability of the strongest
side-channel adversary and let $\UCB_\beta(\hatshortcutDisc)$ be a one-sided upper
confidence bound.
The strict certification inequality becomes
\begin{equation}\label{eq:formula-66}
\LCB_\alpha(\shat)
>
\left(
\frac{1+\UCB_\beta(\hatshortcutDisc)}{2}
\right)^2.
\end{equation}
This direction is essential: the model side is pushed downward by an LCB, while
the adversary side is pushed upward by a UCB.
Using an LCB on the ceiling would make false certification easier and is
therefore statistically backward.
The strict variant is uniformly harder to pass than the reported calibrated
form: every negative Case B margin in this paper remains negative under the
strict variant, while positive margins could shrink.
We recommend the strict variant whenever a positive Case B certificate supports
a high-stakes claim.

\subsection{Per-Group and Global-Pooled Quantities in the Code}

The SAT analysis code maintains two families of reported quantities.
The RMC family is computed in
\texttt{compute\_rmc\_metrics}: it forms
\texttt{RMC\_theta\_P\_hat}, \texttt{RMC\_theta\_N\_hat},
\texttt{RMC\_lcb\_s\_factorized}, \texttt{RMC\_discrimination\_margin}, and
\texttt{RMC\_signed}.  These per-group columns provide group-level certificate diagnostics; the $N^*$ boundary itself is defined from the global-pooled RMC(CP) column under the Case A ceiling (\Cref{sec:table_estimators}).  Assignment-derived quantities are stored separately and are used for ADR$^{+w}$ diagnostics, not for the RMC margin.

The global-pooled family is added by \texttt{filter\_rmc\_columns.py}.
It pools SAT-side and UNSAT-side predictions across groups at a fixed model and
$N$, then computes:
\begin{equation}\label{eq:formula-67}
\begin{array}{ll}
\texttt{global\_pooled\_RMC\_without\_LCB} &\leftrightarrow \text{RMC(no LCB)},\\
\texttt{global\_pooled\_RMC\_signed} &\leftrightarrow \text{RMC(CP)},\\
\texttt{wilson\_pooled\_RMC\_signed} &\leftrightarrow \text{RMC(Wilson)}.
\end{array}
\end{equation}
The normalized companions divide the clipped positive part by $0.75$, the
available headroom under $\shortcutCeiling=0.25$.

\subsection{Why the Main Tables Put F1 Before ADR}

The per-$N$ tables are ordered to follow the logic of the framework.
F1$_{\mathrm{SAT}}$ and F1$_{\mathrm{UNSAT}}$ first expose one-sided label bias.
ADR and ADR$^{+w}$ then show pair-level decision and witness-sensitive
performance.
The three RMC blocks finally subtract the shortcut ceiling, first without a
statistical discount, then with exact CP and approximate Wilson lower bounds.
This layout makes the empirical argument readable from left to right:
standard binary metrics, pair-level behavior, and finally certified excess over
the shortcut-only reference baseline.


\bibliographystyle{IEEEtran}
\bibliography{paper}

\end{document}
